
\chapter{Spatial Perception}
\label{chapter:3:sec:6}
\label{sec:perception}
% \label{sec:6}


% \lettrine[lines=3]{\initfamily\textcolor{darkgreen}{S}}{patial perception} 
\FramedInitial{S}{patial} perception is the foundational layer of spatial intelligence: it turns raw sensory input into the structured spatial evidence that everything downstream depends on, answering what is present, where it lies, and how it moves. It is what stands between an agent and a scene, and the reliability of every later inference is bounded by the reliability of what perception delivers.

The chapter follows a single progression, from the elemental to the structured. It begins with recognizing the discrete entities that populate a scene, grounds them in metric geometry, then binds them into relational and scene-level structure, extends that structure across time, and finally broadens it beyond vision by fusing complementary modalities. Each step presupposes the ones before it, so that perception builds from isolated percepts toward a dynamic, multimodal model of space that reasoning and generation can act on.



\section{Spatial Entity Recognition}
\label{sec:entity_recognition}

\FramedInitial{B}{efore} space can be measured or organized, its contents must be named. Spatial entity recognition is that entry point: it answers \emph{what is present}, and every later geometric and relational step builds on the entities it names. Recognition operates at several granularities, and the section is ordered accordingly: detecting discrete object instances, delineating dense regions of pixels or points by segmentation, pinning down sparse keypoints and landmarks, and finally assigning a holistic scene label. What unites them is that they recover identity and category, not yet geometry; where each entity sits and how it is oriented is taken up next.


\subsection{Object Detection}
\label{sec:6-1-2}
% logic: (Xiangtai Li)
% 1, definition
% 2, method before
% 3, method after deep learning
%  % single stage detectors
%  % two-stage detectors
%  % transformer-based detectors. (Detr)
%  % open-vocabulary / open-world detectors.
%  % video object detection methods.
% 4, dataset and benchmark info
% 5, trends and future work.

\subsubsection{Task Definition}
Given a 2D image $I$, \emph{2D object detection} identifies instances of predefined object classes and localizes each one in the image plane. The output is a set of detections
\begin{equation}
\mathcal{D}=\{(c_i, b_i, s_i)\}_{i=1}^{N},
\end{equation}
where $c_i$ is the predicted class, $b_i=(x_i,y_i,w_i,h_i)$ a 2D bounding box parameterized by center and size, and $s_i\in[0,1]$ a confidence score. In short, the task answers \emph{what objects are where}. Figure~\ref{fig:p3_s6_1_2_objdet_overview} summarizes the 2D task and its detector families.

\emph{3D object detection} keeps the same set-of-detections output but places each box in metric 3D space. The input is a LiDAR point cloud $\mathcal{P}=\{\mathbf{p}_j\}_{j=1}^{M}$ with $\mathbf{p}_j\in\mathbb{R}^{3}$ (often with an intensity channel), one or more calibrated images, or both, and each box becomes
\begin{equation}
b_i=(x_i,y_i,z_i,\;l_i,w_i,h_i,\;\theta_i),
\end{equation}
a 3D center, the physical length, width, and height, and a heading angle $\theta_i$ about the vertical axis; indoor benchmarks often use axis-aligned boxes and drop $\theta_i$. Because the box is expressed in the sensor or ego frame rather than the image plane, a 3D detection already carries metric position and heading, which is why driving and robotics stacks consume it directly.
% \end{definition}

\begin{figure}[!t]
    \centering
    \includegraphics[width=\linewidth]{figures/part3/section6/object_detection.pdf}
    \caption{2D object detection: each detected instance is a tuple $(c_i, b_i, s_i)$ of class, bounding box, and confidence score (top), and the progression of detector families with representative methods under each (bottom).}
    \label{fig:p3_s6_1_2_objdet_overview}
\end{figure}

\subsubsection{Methodology}

\paragraph{Before Deep Learning.}
Early detectors~\cite{zou2019object} ran a sliding window over multi-scale image pyramids and classified each candidate region with hand-crafted features and shallow models. Viola--Jones~\cite{viola_jones_2001} combined Haar-like features with an AdaBoost cascade for real-time face detection; HOG features with linear SVMs~\cite{dalal_triggs_2005} became the standard pipeline for pedestrians and generic objects; and deformable part-based models (DPMs)~\cite{felzenszwalb_dpm_2010} represented an object as parts connected by deformable constraints. Manual feature design, dense-window evaluation, and limited robustness to pose, scale, and occlusion capped this line and motivated the move to learned features.

\paragraph{Two-stage Detectors.}
R-CNN~\cite{girshick2014rcnn} first generates category-independent region proposals and then classifies each with a CNN, a large accuracy gain over hand-crafted pipelines bought with high computational cost. Fast R-CNN~\cite{girshick2015fast} shares convolutional computation across proposals through ROI pooling and trains classification and box regression jointly; Faster R-CNN~\cite{ren2015faster} adds a Region Proposal Network so that proposals are nearly free and the whole detector trains end-to-end. Sparse R-CNN~\cite{sun2021sparsercnn} later replaced dense anchors with a small set of learnable proposals, reaching end-to-end detection without a transformer.

\paragraph{One-stage Detectors.}
One-stage detectors drop the proposal step and predict classes and boxes directly from the feature map. YOLO~\cite{redmon2016yolo} casts detection as a single regression over the whole image, trading localization accuracy for speed; YOLOv2 and YOLOv3~\cite{redmon2017yolo9000, redmon2018yolov3} add anchors, multi-scale prediction, and stronger backbones, and later versions~\cite{bochkovskiy2020yolov4, jocher2020yolov5} refine training and deployment. SSD~\cite{liu2016ssd} predicts from multiple feature maps at different scales, and RetinaNet~\cite{lin2017focal} introduces the focal loss to counter the extreme foreground--background imbalance, bringing one-stage accuracy level with two-stage methods. Recent YOLO models move to anchor-free, decoupled heads (YOLOv8~\cite{jocher2023yolov8}) and NMS-free end-to-end training (YOLOv10~\cite{wang2024yolov10}), converging with the set-prediction detectors below.

\paragraph{Transformer-based Detectors.}
DETR~\cite{carion2020end} reformulates detection as direct set prediction with a transformer encoder--decoder, removing anchors and non-maximum suppression at the price of slow convergence and weak small-object performance. Deformable DETR~\cite{Zhu2021DeformableDETR} attends to a sparse set of sampling points, which fixes both; Conditional DETR~\cite{meng2021conditional} conditions decoder attention on spatial queries for faster training; DINO~\cite{zhang2023dino} adds denoising training and better query initialization; and RT-DETR~\cite{zhao2024detrs} brings the family to real-time speed.

\paragraph{Open-Vocabulary Detection.}
Open-vocabulary detection (OVD) recognizes objects outside a fixed label set by aligning region features with text embeddings. The line evolved from embedding distillation to large-scale vision--language pretraining to data-driven vocabulary expansion: ViLD~\cite{gu2021open} and OVR-CNN~\cite{zareian2021ovrcnn} distill or learn a shared region--text embedding space from image--text pairs; GLIP~\cite{li2022grounded} unifies detection with phrase grounding under grounded language--image pretraining; and Detic~\cite{zhou2022detecting} shows that image-level labels over a large vocabulary suffice to train the classifier of an open-vocabulary detector. OWL-ViT~\cite{minderer2022simple}, Grounding DINO~\cite{liu2024grounding}, and YOLO-World~\cite{cheng2024yoloworld} push the same idea toward simpler transfer, tighter grounding, and real-time speed; see the surveys~\cite{zhu2024survey, wu2024towards} for a fuller account.

\paragraph{3D Object Detection.}
3D detectors differ mainly in how they turn sensor data into something a network can process; the detection heads on top largely reuse the 2D families above. On LiDAR, VoxelNet~\cite{zhou2018voxelnet} learns features inside voxels and runs 3D convolutions over the grid, and PointPillars~\cite{lang2019pointpillars} collapses the vertical axis into pillars so that an ordinary 2D CNN on the bird's-eye-view (BEV) plane does the rest at real-time speed. Point-based detectors keep the raw points: PointRCNN~\cite{shi2019pointrcnn} generates proposals from PointNet++ features~\cite{qi2017pointnetplus} of segmented foreground points, PV-RCNN~\cite{shi2020pvrcnn} combines voxel and point features. CenterPoint~\cite{yin2021center} replaces anchors with BEV center heatmaps, the 3D counterpart of anchor-free 2D detection. Indoor scenes are cluttered and object centers often fall in empty space, which favors voting: in VoteNet~\cite{qi2019votenet} surface points vote for object centers before grouping, and 3DETR~\cite{misra20213detr} carries DETR-style set prediction over to point clouds.

Camera-only 3D detection has to infer the depth that LiDAR measures. Monocular detectors regress depth, size, and orientation on top of a 2D detection (Section~\ref{sec:height_distance} treats them as metric estimation). Pseudo-LiDAR~\cite{wang2019pseudolidar} instead converts a predicted depth map into a point cloud and runs a LiDAR detector on it, showing that much of the camera--LiDAR gap came from the representation rather than from depth quality alone. Multi-camera systems lift image features into a shared BEV space, either bottom-up, as Lift-Splat-Shoot~\cite{philion2020lss} does by predicting a depth distribution per pixel and splatting features onto the ground plane, or top-down, as DETR3D~\cite{wang2022detr3d} and BEVFormer~\cite{LiWLXSLQD22} do by projecting 3D queries or BEV cells into the images to sample features. Fusion combines both sensors, from Frustum PointNets~\cite{qi2018frustum}, which crops the point cloud with 2D image detections, to BEVFusion~\cite{liu2023bevfusion}, which merges camera and LiDAR features on a common BEV map. Cube R-CNN~\cite{Omni3D-Brazil0SR0G23} trains a single image-based detector across indoor and outdoor domains on Omni3D. Deployment constraints for driving are discussed in Section~\ref{sec:p4_s9_2}, and video detection, which adds cross-frame association to the same output set, under tracking in Section~\ref{sec:object_tracking}.

\begin{table*}[!t]
\centering
\caption{Summary of object detection datasets across modalities and domains, with image and bounding-box counts per official split; ImageNet VID counts refer to video clips.}
\label{tab:object_detection_dataset}
\fontsize{9}{11}\selectfont
\setlength{\tabcolsep}{1.5mm}
\begin{tabular}{l cc cc cc cc }
\toprule
\multirow{2}{*}{\textbf{Dataset}}   & \multirow{2}{*}{\textbf{Modality, Domain}}   & \multirow{2}{*}{\textbf{Class}}   & \multicolumn{2}{c}{\textbf{train}} & \multicolumn{2}{c}{\textbf{validation}} & \multicolumn{2}{c}{\textbf{test}}  \\
\cmidrule(r){4-5}\cmidrule(r){6-7}\cmidrule(r){8-9}
  & & & \textbf{images} & \textbf{boxes} & \textbf{images} & \textbf{boxes} & \textbf{images} & \textbf{boxes} \\
\midrule
VOC-2007~\cite{pascal-voc-2007} & image, general & 20 & 2,501 & 6,301 & 2,510 & 6,307 & 4,952 & 14,976\\
VOC-2012~\cite{pascal-voc-2012} & image, general & 20 & 5,717 & 13,609 & 5,823 & 13,841 & 10,991 & - \\
ILSVRC-2014~\cite{ImageNet} & image, general & 200 & 456,567 & 478,807 & 20,121 & 55,502 & 40,152 & -\\
MS-COCO-2015~\cite{lin2014microsoft} & image, general & 80 & 82,783 & 604,907 & 40,504 & 291,875 & 81,434 & - \\
MS-COCO-2017~\cite{lin2014microsoft} & image, general & 80 & 118,287 & 860,001 & 5,000 & 36,781 & 40,670 & - \\
Objects365-2019~\cite{shao2019objects365} & image, general & 365 & 600,000 & 9,623,000 & 39,000 & 479,000 & 100,000 & 170,000\\
OpenImages~\cite{kuznetsova2020open} & image, general & 600 & 1,743,042 & 14,610,229 & 41,620 & 303,980 & 125,436 & 937,327 \\
ImageNet VID~\cite{ImageNet} & video, general & 30 & 3,862 & - & 555 & - & 937 & -\\
LVIS~\cite{gupta2019lvis} & image, general & 1,203 & 100,000 & - & - & - & - & - \\
Omni3D~\cite{Omni3D-Brazil0SR0G23} & 3D (image), general  & 98 & 175,573 & - & 19,127 & - & 39,452 & -  \\
\bottomrule
\end{tabular}
\end{table*}

\begin{table*}[!t]
\centering
\caption{Reported COCO performance of representative detector configurations, with evaluation splits, input sizes, and training schedules differing across sources; AP$_S$ / AP$_M$ / AP$_L$ break down by object size (small / medium / large). SSD512* denotes the image-expansion-augmented variant.}
\label{tab:objdet_coco_results}
\fontsize{9}{11}\selectfont
\setlength{\tabcolsep}{1.5mm}
\begin{tabular}{l l c c c c c c c c c}
\toprule
\textbf{Method} & \textbf{Backbone} & \textbf{\#P (M)} & \textbf{FLOPs (G)} & \textbf{AP} & \textbf{AP$_{50}$} & \textbf{AP$_{75}$} & \textbf{AP$_S$} & \textbf{AP$_M$} & \textbf{AP$_L$} & \textbf{Year} \\
\midrule
\multicolumn{11}{l}{\emph{Two-stage detectors}} \\
Faster R-CNN~\cite{ren2015faster}            & R50-FPN     & 41.5 & 207   & 37.4 & 58.1 & 40.4 & 21.2 & 41.0 & 48.1 & 2015 \\
Sparse R-CNN~\cite{sun2021sparsercnn}        & R50-FPN     & 106  & 174   & 45.0 & 64.1 & 49.4 & 26.9 & 47.2 & 59.5 & 2021 \\
\midrule
\multicolumn{11}{l}{\emph{One-stage CNN detectors}} \\
SSD512*~\cite{liu2016ssd}                     & VGG-16      & 27   & 88    & 28.8 & 48.5 & 30.3 & 10.9 & 31.8 & 43.5 & 2016 \\
RetinaNet~\cite{lin2017focal}                & R50-FPN     & 38   & 198   & 36.5 & 55.0 & 39.1 & 19.9 & 40.3 & 48.7 & 2017 \\
YOLOv3~\cite{redmon2018yolov3}               & DarkNet-53  & 62   & 156   & 33.0 & 57.9 & 34.4 & 18.3 & 35.4 & 41.9 & 2018 \\
\midrule
\multicolumn{11}{l}{\emph{Modern YOLO}} \\
YOLOv5-s~\cite{jocher2020yolov5}             & CSPDarkNet  & 7.2  & 16.5  & 37.4 & --   & --   & --   & --   & --   & 2020 \\
YOLOv8-x~\cite{jocher2023yolov8}             & C2f-based    & 68.2 & 257.8 & 53.9 & --   & --   & --   & --   & --   & 2023 \\
YOLOv10-l~\cite{wang2024yolov10}             & C2f/CIB    & 24.4 & 120.3 & 53.2 & --   & --   & --   & --   & --   & 2024 \\
\midrule
\multicolumn{11}{l}{\emph{Transformer set-prediction detectors}} \\
DETR~\cite{carion2020end}                    & R50, 500\,ep & 41   & 86    & 42.0 & 62.4 & 44.2 & 20.5 & 45.8 & 61.1 & 2020 \\
Deformable DETR~\cite{Zhu2021DeformableDETR} & R50, 50\,ep  & 40   & 173   & 43.8 & 62.6 & 47.7 & 26.4 & 47.1 & 58.0 & 2021 \\
DINO~\cite{zhang2023dino}                    & R50, 4-sc., 12\,ep & 47 & 279 & 49.0 & 66.6 & 53.5 & 32.0 & 52.3 & 63.0 & 2023 \\
RT-DETR~\cite{zhao2024detrs}                 & R50          & 42   & 136   & 53.1 & 71.3 & 57.7 & 34.8 & 58.0 & 70.0 & 2024 \\
\bottomrule
\end{tabular}
\end{table*}




\subsubsection{Datasets and Metrics}
Table~\ref{tab:object_detection_dataset} summarizes the benchmarks that shaped the field. \textbf{PASCAL VOC}~\cite{everingham2010pascal} fixed the evaluation protocol and remained the reference through the early deep-learning era despite its modest scale. \textbf{MS COCO}~\cite{lin2014microsoft}, with complex everyday scenes and segmentation masks alongside boxes, then became the standard for both detection and segmentation. \textbf{Open Images}~\cite{kuznetsova2020open}, \textbf{Objects365}~\cite{shao2019objects365}, and \textbf{LVIS}~\cite{gupta2019lvis} scale the vocabulary to hundreds and then over a thousand categories, with LVIS deliberately emphasizing rare classes, which turned it into the canonical long-tail and open-world benchmark.

Evaluation follows the \emph{precision--recall} framework. A detection $(\hat{c}, \hat{b}, \hat{s})$ is correct when it matches a previously unmatched ground-truth instance of class $\hat{c}$, whose box $b$ overlaps $\hat{b}$ at or above the required \emph{Intersection over Union} (IoU) threshold
\begin{equation}
\mathrm{IoU}(\hat{b}, b) = \frac{|\hat{b} \cap b|}{|\hat{b} \cup b|}.
\end{equation}
Sorting the detections of a class by confidence traces a precision--recall curve whose area is the \emph{Average Precision} (AP); the \emph{mean Average Precision} (mAP) averages AP over classes. Benchmarks differ chiefly in the IoU threshold: PASCAL VOC uses a single $\mathrm{IoU} = 0.5$ (AP$_{50}$), whereas COCO averages over ten thresholds from $0.5$ to $0.95$ to obtain its primary AP and reports AP$_{50}$ and AP$_{75}$ as diagnostic slices. COCO further stratifies by object size into AP$_S$ (area $<32^2$ pixels), AP$_M$ ($32^2$--$96^2$), and AP$_L$ ($>96^2$), which exposes the persistent difficulty of small objects.

3D detection is benchmarked mainly on driving data, KITTI~\cite{KITTI-GeigerLU12}, nuScenes~\cite{caesar2020nuscenes}, and the Waymo Open Dataset~\cite{WaymoOpen-SunKDCPTGZCCVHN20}, and indoors on SUN RGB-D~\cite{song2015sunrgbd} and ScanNet~\cite{dai2017scannet}; Omni3D (Table~\ref{tab:object_detection_dataset}) pools image-based 3D annotations from both. The IoU above is computed between 3D boxes or their BEV footprints: KITTI reports 3D and BEV AP at IoU $0.7$ for cars and $0.5$ for pedestrians and cyclists, and indoor benchmarks report mAP at 3D IoU $0.25$ and $0.5$. nuScenes matches detections by center distance on the ground plane ($0.5$, $1$, $2$, and $4$\,m) instead of IoU and folds translation, scale, orientation, velocity, and attribute errors into its nuScenes Detection Score (NDS); Waymo weights AP by heading accuracy (APH).

Table~\ref{tab:objdet_coco_results} collects reported COCO results across the detector families.
On the canonical ResNet-50 backbone, AP has climbed from $37.4$ for Faster R-CNN to $53.1$ for RT-DETR over roughly a decade, with much of the recent gain attributable to transformer-based set prediction (DETR $\rightarrow$ Deformable DETR $\rightarrow$ DINO) and improved denoising-style query initialization.
The AP$_S$ column highlights the well-documented small-object weakness of vanilla DETR ($20.5$); deformable attention and denoising progressively close this gap (Deformable DETR $26.4$, DINO $32.0$, RT-DETR $34.8$).
On the efficiency frontier, the modern YOLO line (YOLOv8-x at $53.9$, YOLOv10-l at $53.2$) matches RT-DETR in accuracy at substantially higher throughput, so the traditional accuracy--speed divide between CNN one-stage detectors and transformer set prediction has largely closed.



% \noindent
\subsubsection{Challenges}
\textbf{The accuracy--efficiency trade-off remains unresolved.}
The transformer-based detectors that top COCO and open-vocabulary leaderboards are still expensive at inference: even with deployment stacks such as TensorRT, their latency and parameter counts rule out most mobile and embedded targets, while the lightweight real-time detectors that fit those budgets trail them in accuracy. The same tension recurs for every recognition task in this section.
\textbf{Open-world detection degrades on the hardest instances.}
Open-vocabulary detectors extend recognition beyond a fixed label set, yet still miss small objects and mislabel rare categories, precisely the long-tail regime that motivated LVIS.
\textbf{Detection inside multimodal large language models is still immature.} Recent systems fold detection into MLLMs as a grounding primitive that ties image regions to language, but even the strongest closed-source models post markedly lower numbers than specialist detectors on standard benchmarks~\cite{zhu2024survey, wu2024towards}.
\textbf{3D detection is bounded by depth.} For camera-only detectors, localization error is dominated by depth error, which grows with distance~\cite{ma2021delving}; LiDAR measures depth directly but its returns thin out with range, so distant and small objects are seen through a handful of points.





\subsection{Segmentation}
\label{sec:segmentation}
% logic: (xiangtai li)
% 1, definition of various segmentation (two or three paragraph)
% 2, methods before deep learning.
% 3, methods using deep learning (FCN-based/U-Net based)
% 4, transformer-based methods.
%    image/video/3D segmentation (transformers)
% 5, datasets.
% 6，trends and challenges.

Segmentation partitions an image, a video, or a 3D scan into meaningful regions. Before deep learning it was a mid-level grouping problem, clustering pixels by similarity somewhere between low-level edge detection and high-level recognition; learned models turned it into dense prediction, with class labels, instance identities, and temporal consistency attached to every region.
Figure~\ref{fig:p3_s6_1_3_segmentation_examples} shows the principal segmentation map types produced on a common scene.

\begin{figure}[!ht]
    \centering
    \includegraphics[width=0.99\linewidth]{figures/part3/section6/segmentation.pdf}
    \caption{Examples of different segmentation maps: semantic, instance, and panoptic on a shared scene.}
    \label{fig:p3_s6_1_3_segmentation_examples}
\end{figure}

\subsubsection{Task Definition}

\textit{Semantic Segmentation}:
Given an input image $I \in \mathbb{R}^{H \times W \times 3}$, \emph{semantic segmentation} assigns a class label to every pixel, producing a dense label map
\begin{equation}
    S \in \{1, \dots, C\}^{H \times W},
\end{equation}
where $C$ is the number of classes; instances of the same class are not distinguished.

\textit{Instance and Panoptic Segmentation}:
\emph{Instance segmentation} detects and segments each object instance separately. The output is a set of labeled masks
\begin{equation}
\mathcal{Y} = \{(m_i, c_i)\}_{i=1}^{N},
\end{equation}
where $m_i \in \{0,1\}^{H \times W}$ is the binary mask of the $i$-th instance and $c_i$ its class; masks may overlap, and only countable \emph{thing} classes are covered. \emph{Panoptic segmentation} keeps the same output form but requires the $N$ masks to form a non-overlapping partition of the image that covers \emph{thing} instances and \emph{stuff} regions alike, each stuff class as a single mask.

\textit{Video Segmentation}:
Given a video clip $V \in \mathbb{R}^{T \times H \times W \times 3}$, \emph{video segmentation} predicts temporally consistent masks across frames, output as a set of mask tubes
\begin{equation}
\mathcal{Y} = \{(m_i, c_i)\}_{i=1}^{N}, \quad m_i \in \{0,1\}^{T \times H \times W},
\end{equation}
each tracking an object or region over time, with semantic, instance, and panoptic variants as above.

\textit{3D Segmentation}:
Given a point cloud $\mathcal{P}=\{\mathbf{p}_j\}_{j=1}^{M}$ with $\mathbf{p}_j\in\mathbb{R}^{3}$ and optional color or intensity, \emph{3D semantic segmentation} labels every point,
\begin{equation}
S\in\{1,\dots,C\}^{M},
\end{equation}
and 3D instance and panoptic segmentation predict point masks $m_i\in\{0,1\}^{M}$ under the same set form and overlap constraints as their image counterparts; voxels or mesh vertices take the place of points for other 3D representations. The definitions carry over unchanged. What changes is the input: the points are unordered, unevenly spaced, and sparse, and that is where 3D methods depart from image methods.


\subsubsection{Methodology}
Figure~\ref{fig:p3_s6_1_3_segmentation_overview} summarizes the method families surveyed below.

\begin{figure}[!t]
    \centering
    \includegraphics[width=\linewidth]{figures/part3/section6/segmentation-method.pdf}
    \caption{Methodological progression of image and video segmentation, from pre-deep-learning methods to foundation and promptable segmenters, with representative methods under each stage.}
    \label{fig:p3_s6_1_3_segmentation_overview}
\end{figure}

% \noindent
\paragraph{Before Deep Learning.}
Classical segmentation grouped pixels either by appearance or by minimizing an energy. The first family includes \emph{thresholding} of pixel intensities~\cite{otsu1979threshold}, \emph{region growing} that merges similar neighbors~\cite{adams1994region}, \emph{$k$-means clustering} in feature space~\cite{macqueen1967kmeans}, and the \emph{watershed} transform, which treats the image as a topographic surface and segments by catchment basins~\cite{vincent1991watersheds}. The second family includes \emph{active contours (snakes)} that evolve a curve toward object boundaries~\cite{kass1988snakes}, \emph{graph cuts} that solve certain segmentation energies to global optimality~\cite{boykov2001graphcuts}, and \emph{Markov} and \emph{conditional random fields} that impose spatial coherence probabilistically~\cite{geman1984mrf,lafferty2001crf}; the last of these carried over directly into early learning-based segmenters.



\paragraph{Fully Convolutional Networks.}
FCNs~\cite{long2015fcn} replace fully connected layers with convolutions, enabling end-to-end dense prediction at arbitrary input size. U-Net~\cite{ronneberger2015unet} adds a symmetric encoder--decoder with skip connections that fuse low- and high-level features, SegNet~\cite{badrinarayanan2017segnet} reuses max-pooling indices for cheap upsampling and boundary recovery, and PSPNet~\cite{zhao2017pspnet} aggregates global and local context through pyramid pooling.

\paragraph{Dilated and Multi-Scale Context Modeling (DeepLab Series).}
The DeepLab series introduces atrous (dilated) convolution to enlarge the receptive field without losing resolution (v1~\cite{chen2015deeplab}), Atrous Spatial Pyramid Pooling for multi-scale context (v2~\cite{chen2017deeplabv2}, refined with batch normalization and stronger backbones in v3~\cite{chen2017deeplabv3}), and an encoder--decoder that recovers boundary detail from low-level features (v3+~\cite{chen2018deeplabv3plus}).

\paragraph{Region-Based Mask Prediction (Mask R-CNN Family).}
Mask R-CNN~\cite{he2017mask} extends Faster R-CNN with a parallel mask branch and ROIAlign, which removes the quantization error of ROI pooling and makes pixel-accurate masks possible. PANet~\cite{liu2018panet}, Hybrid Task Cascade~\cite{chen2019htc}, and PointRend~\cite{kirillov2020pointrend} refine it through better feature flow, cascaded box--mask prediction, and point-wise boundary rendering respectively; this two-stage, region-based family dominated instance segmentation until transformers arrived.



% \noindent
\paragraph{Transformer Backbones and Representation Learning.}
Transformer-based segmentation, following DETR, favors unified query-based formulations that bridge detection and segmentation, predict masks end-to-end without hand-crafted post-processing, and allow joint learning across tasks and modalities. The first step was the backbone: SETR~\cite{zheng2021setr} applies a pure ViT to semantic segmentation with a simple decoder, SegFormer~\cite{xie2021segformer} pairs a lightweight hierarchical transformer with an MLP decoder for a favorable accuracy--efficiency trade-off, and Swin Transformer~\cite{liu2021swin} uses shifted-window attention to produce CNN-like multi-scale features and has become a default segmentation backbone. Self-supervised pretraining such as MAE~\cite{he2022masked} supplies ViT backbones that transfer well to dense prediction.


\paragraph{Query-based Transformer Segmentation.}
Query-based methods cast segmentation as set prediction with object or mask queries. MaX-DeepLab~\cite{wang2021maxdeeplab} is the first purely mask-based panoptic framework without box supervision; MaskFormer~\cite{cheng2021maskformer} unifies semantic and panoptic segmentation by predicting mask embeddings with class labels; K-Net~\cite{zhang2021knet} treats segmentation as kernel prediction with iterative refinement; and Mask2Former~\cite{cheng2022mask2former} adds masked cross-attention and multi-scale decoding, setting a strong baseline across all three tasks. QueryInst~\cite{fang2021queryinst} and Panoptic SegFormer~\cite{li2022panopticsegformer} are variants that respectively generate instance-specific kernels from queries and decouple semantic from instance queries.


\paragraph{Video and Multimodal Transformer-based Segmentation.}
VisTR~\cite{wang2021vistr} predicts spatio-temporal mask tubes end-to-end from object queries; Video K-Net~\cite{li2022videoknet} unifies video semantic, instance, and panoptic segmentation by associating queries across time; and TubeFormer~\cite{wang2022tubeformer} adds tube-level attention for long-term consistency. X-Decoder~\cite{zou2023xdecoder} unifies image segmentation with language-based prediction, while OMGSeg~\cite{OMGSeg} covers image and video segmentation across open-vocabulary and class-agnostic settings within one model.


\paragraph{SAM series.} SAM~\cite{kirillov2023sam} made segmentation promptable: trained on the SA-1B dataset of over one billion masks, it decouples a heavy image encoder from a light prompt-conditioned mask decoder and produces masks from points, boxes, or text with strong zero-shot generalization. SAM~2~\cite{ravi2024sam2} extends this to video, propagating masks across frames with a memory mechanism while keeping interactive control.

\paragraph{3D Point Cloud Segmentation.}
Point-cloud segmenters split by how they impose structure on the points (Section~\ref{sec:rep_3d_point}). Point-based networks consume raw points: PointNet~\cite{qi2017pointnet} applies a shared MLP to every point and pools symmetrically, PointNet++~\cite{qi2017pointnetplus} adds hierarchical neighborhoods, DGCNN~\cite{wang2019dgcnn} builds $k$-nearest-neighbor graphs that are recomputed in feature space at every layer, KPConv~\cite{thomas2019kpconv} defines convolution kernels by points placed in space, and RandLA-Net~\cite{hu2020randlanet} relies on random sampling to scale to scenes of millions of points. Sparse voxel networks such as MinkowskiNet~\cite{choy2019minkowski} and SPVNAS~\cite{tang2020searching} quantize the points and convolve only occupied voxels, which keeps 3D convolution affordable on room- and street-scale scans. Projection methods take the opposite route: RangeNet++~\cite{milioto2019rangenet} unrolls a LiDAR sweep into a range image, runs a 2D segmenter in real time, and projects the labels back onto the points. Transformers entered with Point Transformer~\cite{zhao2021point}, and Point Transformer V3~\cite{wu2024ptv3} replaces exact neighbor search with serialized point orderings so that attention scales to large scenes.

Instance segmentation in 3D retraces the 2D history. PointGroup~\cite{jiang2020pointgroup} predicts per-point offsets toward instance centers and clusters the shifted points, a bottom-up grouping; Mask3D~\cite{schult2023mask3d} applies Mask2Former-style queries to sparse voxel features and predicts instance masks directly. Because 3D labels remain scarce, open-vocabulary 3D segmentation mostly borrows from 2D foundation models: OpenScene~\cite{peng2023openscene} distills CLIP-aligned pixel features into points, and OpenMask3D~\cite{takmaz2023openmask3d} classifies class-agnostic 3D instance masks with CLIP features aggregated over multi-view images.



\begin{table*}[!t]
\centering
\caption{Summary of representative segmentation datasets across 2D and 3D modalities; split sizes count images, scans, or video sequences as indicated. Parenthesized class counts denote common reduced evaluation label sets.}
\label{tab:segmentation_dataset}
\fontsize{9}{11}\selectfont
\setlength{\tabcolsep}{3.7mm}
\begin{tabular}{l c cc ccc }
\toprule
\multirow{2}{*}{\textbf{Dataset}}   & \multirow{2}{*}{\textbf{Purpose}}   & \multirow{2}{*}{\textbf{Classes}} & \multirow{2}{*}{\textbf{Resolution}} & \multicolumn{3}{c}{\textbf{Samples}}  \\
\cmidrule(r){5-7}
% \cmidrule(r){5-6}\cmidrule(r){7-8}
  & & & & \textbf{Train} & \textbf{Valid} & \textbf{Test} \\
\midrule
\rowcolor{lightgrey}\multicolumn{7}{l}{$\bullet$ \textit{2D}}\\
Pascal VOC \cite{everingham2010pascal} & Generic & 21 & Variable & 1464 & 1449 & Private\\
MS COCO \cite{lin2014microsoft} & Generic & 80+ & Variable & 82,783 & 40,504 & 81,434\\
ADE20K \cite{zhou2017ade20k} & Generic & 150 & Variable & 20,210 & 2,000 & 3,000 \\
COCO-stuff \cite{COCO-Stuff-CaesarUF18} & Generic & 172 & Variable & 118,000 & 5,000 & 20,000 \\
IDD \cite{IDD-VarmaSNCJ19} & Urban/Driving & 34 & Variable & 6,993 & 981 & 2,029\\
Dark Zurich \cite{SakaridisDG19} & Urban/Driving & 19 & 1920$\times$1080& 8,377  & 50 & 151\\
ACDC \cite{ACDC-SakaridisDG21} & Driving & 19 & 2048$\times$ 1024& 1,600  & 406 & 2,000\\
PartImageNet \cite{PartImageNet-HeYYKYCLYYY22} & Generic-Part & 158 & Variable & 16,540 & 2,957 & 4,598\\
SiftFlow \cite{SiftFlow-LiuYT09} & Outdoor & 33 & 256$\times$256 & 2,488 & - & 200\\
BSDS300 \cite{BSD-MartinFTM01}& Boundary/Region & -- & Variable & 200 & - & 100 \\
Youtube-Objects \cite{PrestLCSF12} & Generic/Video & 10 & 480$\times$360 & 126 & - & - \\
Cityscapes~\cite{cordts2016cityscapes}  & Urban & 30 (19) & 2048 $\times$1024 & 2975 & 500 & 1525\\
CamVid \cite{CamVid-BrostowFC09} & Urban/Driving & 32 (11) & 960$\times$720 & 367 & 101 & 233 \\
Mapillary Vistas~\cite{neuhold2017mapillary} & Street & 65 & - & 18,000 & - & - \\
\hline
\rowcolor{lightgrey}\multicolumn{7}{l}{$\bullet$ \textit{3D}}\\
NYU-D V2~\cite{silberman2012nyudv2} & Indoor & 40 (13) & 640$\times$480 & 795 & - & 654 \\
SUN-3D~\cite{xiao2013sun3d} & SfM/Video & - & 640$\times$480 & 415 seqs & - & - \\
SUN RGB-D~\cite{song2015sunrgbd} & Indoor & 37 & Variable & 5,285 & - & 5,050 \\
ScanNet~\cite{dai2017scannet} & Indoor & 20 & Variable & 1,201 & 312 & 100 \\
SemanticKITTI~\cite{SemanticKITTI-BehleyGMQBSG19} & LiDAR/Driving & 28 & Point cloud & 19,130 & 4,071 & 20,351 \\
\bottomrule
\end{tabular}
\end{table*}

% \noindent
\subsubsection{Datasets and Metrics}
Table~\ref{tab:segmentation_dataset} summarizes the main benchmarks. \textbf{PASCAL VOC}~\cite{everingham2010pascal} anchored the classical semantic-segmentation era with 20 foreground classes plus background. \textbf{MS COCO}~\cite{lin2014microsoft}, with semantic, instance, box, and panoptic annotations over complex object-centric scenes, is the shared testbed for all three modern task types. \textbf{Cityscapes}~\cite{cordts2016cityscapes} and \textbf{Mapillary Vistas}~\cite{neuhold2017mapillary} specialize to street scenes, the latter stressing long-tail classes and diverse geographies, and \textbf{ADE20K}~\cite{zhou2017ade20k} targets holistic scene parsing over 150 object and stuff classes. The 3D rows of the table cover indoor RGB-D scans (NYU-D V2, SUN RGB-D, ScanNet) and outdoor LiDAR (SemanticKITTI), and S3DIS~\cite{armeni2016s3dis} is a further indoor standard.

Segmentation evaluation reuses IoU (Section~\ref{sec:6-1-2}) with masks in place of boxes. For semantic segmentation the headline metric is the \emph{mean Intersection over Union} (mIoU),
\begin{equation}
\mathrm{mIoU} = \frac{1}{C} \sum_{k=1}^{C} \frac{\mathrm{TP}_k}{\mathrm{TP}_k + \mathrm{FP}_k + \mathrm{FN}_k},
\end{equation}
where $C$ is the number of classes and $\mathrm{TP}_k$, $\mathrm{FP}_k$, $\mathrm{FN}_k$ count the true-positive, false-positive, and false-negative pixels for class $k$.
Because the average is unweighted, mIoU is sensitive to rare classes, which is one motivation behind ADE20K and Mapillary Vistas; pixel accuracy is dominated by frequent classes (\eg, sky, road) and is less informative.

For \emph{instance segmentation}, the COCO mask-AP protocol matches predictions to ground truth by mask IoU and averages AP over IoU thresholds $0.5$--$0.95$.
For \emph{panoptic segmentation}, \emph{Panoptic Quality}~\cite{kirillov2019panoptic} is $\mathrm{PQ} = \mathrm{SQ} \times \mathrm{RQ}$, the average IoU of matched segments times the F1 score over matched, false-positive, and false-negative segments, treating things and stuff uniformly.
For \emph{video object segmentation}, the DAVIS $\mathcal{J}\&\mathcal{F}$ metric averages region similarity ($\mathcal{J}$, mask IoU) and contour accuracy ($\mathcal{F}$, computed on mask boundaries).
3D semantic segmentation reports the same mIoU computed over points or voxels, and 3D instance segmentation on ScanNet reports AP at mask IoU thresholds averaged over $0.5$--$0.95$, with AP$_{50}$ and AP$_{25}$ as looser slices.

\begin{table*}[!t]
\centering
\caption{Reported performance of representative segmentation methods on standard benchmarks: ADE20K mIoU and Cityscapes mIoU (semantic), COCO Mask AP (instance), and COCO PQ (panoptic). Numbers follow each paper's own protocol (split, single-/multi-scale) and are indicative only; ``--'' marks a metric not reported.}
\label{tab:segmentation_results}
\fontsize{9}{11}\selectfont
\setlength{\tabcolsep}{1.4mm}
\begin{tabular}{l l c c c c c}
\toprule
\textbf{Method} & \textbf{Backbone} & \textbf{ADE20K mIoU} & \textbf{Cityscapes mIoU} & \textbf{COCO Mask AP} & \textbf{COCO PQ} & \textbf{Year} \\
\midrule
\rowcolor{lightgrey}\multicolumn{7}{l}{$\bullet$\emph{FCN family and dilated convolution}} \\
FCN-8s~\cite{long2015fcn}                  & VGG-16        & 29.4 & 62.7 & --   & --   & 2015 \\
PSPNet~\cite{zhao2017pspnet}               & R101          & 43.3 & 80.2 & --   & --   & 2017 \\
DeepLab v3+~\cite{chen2018deeplabv3plus}   & Xception-65   & 45.6 & 82.1 & --   & --   & 2018 \\
\midrule
\rowcolor{lightgrey}\multicolumn{7}{l}{$\bullet$ \emph{Region-based instance segmentation}} \\
Mask R-CNN~\cite{he2017mask}               & R50-FPN       & --   & --   & 35.4 & --   & 2017 \\
HTC~\cite{chen2019htc}                     & R50-FPN       & --   & --   & 38.4 & --   & 2019 \\
\midrule
\rowcolor{lightgrey}\multicolumn{7}{l}{$\bullet$ \emph{Transformer backbones}} \\
SETR-MLA~\cite{zheng2021setr}              & ViT-L         & 50.3 & 82.2 & --   & --   & 2021 \\
SegFormer-B5~\cite{xie2021segformer}       & MiT-B5        & 51.0 & 84.0 & --   & --   & 2021 \\
Swin-L + UperNet~\cite{liu2021swin}        & Swin-L        & 53.5 & --   & --   & --   & 2021 \\
\midrule
\rowcolor{lightgrey}\multicolumn{7}{l}{$\bullet$ \emph{Universal query-based segmentation}} \\
MaskFormer~\cite{cheng2021maskformer}      & Swin-L        & 55.6 & --   & --   & 52.7 & 2021 \\
Mask2Former~\cite{cheng2022mask2former}    & Swin-L        & 57.7 & --   & 50.1 & 57.8 & 2022 \\
\bottomrule
\end{tabular}
\end{table*}

Table~\ref{tab:segmentation_results} reports results across these four protocols. The clearest trend is on semantic segmentation, where ADE20K mIoU has roughly doubled in under a decade, from FCN-8s's $29.4$ to Mask2Former's $57.7$; the largest single jumps track two shifts, first to transformer backbones (SETR, SegFormer, Swin-L) and then to universal query-based formulations. That second shift is what unifies the table: MaskFormer is the first entry to report semantic mIoU and panoptic PQ within one framework, and Mask2Former extends the same architecture to all three tasks with competitive results on each. Instance segmentation followed a partly separate path, climbing from Mask R-CNN ($35.4$) and HTC ($38.4$) on R50-FPN to $50.1$ once Mask2Former is paired with a Swin-L backbone, which confirms that transformer backbones lift mask quality at the instance level as well, not only semantic accuracy.


% \noindent
\subsubsection{Challenges}
The cost of self-attention at high resolution, over long video, and on dense 3D data sharpens the accuracy--speed tension noted for detection, and two problems are specific to segmentation.
\textbf{Long-video, multi-object consistency is hard.} Holding spatial precision and temporal coherence together remains difficult in cluttered scenes, particularly for video panoptic segmentation and segmentation-based multi-object tracking, where objects enter, leave, and occlude one another over long horizons.
\textbf{Segmentation is moving toward the open world.} The field is shifting from fixed label sets to class-agnostic, open-vocabulary, and multimodal settings, in which language and other auxiliary signals are folded in to improve spatial reasoning and adaptability, a direction that SAM, X-Decoder, and OMGSeg already anticipate.









\subsection{Keypoint and Landmark Detection}
\label{sec:keypoint_landmark}
Keypoint and landmark detection localizes a set of semantically consistent, precisely positioned points that characterize object geometry in images or videos: eye corners, skeletal joints, or structurally distinctive parts of man-made objects.
Compared with region-based representations, keypoints give a sparser, explicitly geometric encoding that often survives deformation, rotation, and partial occlusion.
Figure~\ref{fig:p3_s6_1_1_keypoint_overview} summarizes the task families and the method families surveyed below.



\subsubsection{Task Definition}
Given an input image $I$, a parametric model $f_{\theta}$ estimates a set of point locations
\begin{equation}
f_{\theta}(I) = \{\mathbf{p}_i\}_{i=1}^{N}, \quad \mathbf{p}_i \in \mathbb{R}^2,
\end{equation}
where each $\mathbf{p}_i$ is the image coordinate of a salient visual or structural element. Landmark detection fixes a predefined semantic identity for each point (a particular eye corner, a particular joint); the two terms are otherwise used interchangeably.

\emph{3D keypoint detection} predicts $\mathbf{p}_i\in\mathbb{R}^{3}$ instead. From images, the depth of each point must be inferred; from depth maps or point clouds $\mathcal{P}$, surface geometry is observed, but occluded or internal landmarks still require inference. For human pose, 3D joints are usually expressed in camera coordinates relative to a root joint (the pelvis), which leaves global depth unconstrained, as a single image cannot fix it.

The task varies along three axes. The input modality sets what geometry is directly available: most detectors work on 2D images (RGB, grayscale, or infrared), where keypoints must be inferred from appearance alone, whereas point clouds expose spatial structure directly; RGB-D or RGB-infrared fusion adds robustness in clutter and poor lighting, and video supplies temporal continuity through motion blur and transient occlusion.
The prediction target ranges from planar image coordinates, which underpin most benchmarks, to 3D keypoints required in robotics and autonomous driving, to 4D trajectories tracked over time.
The application scenario fixes what the points mean. Research has concentrated on humans: body joints for pose estimation, facial landmarks for expression recognition and identity verification, and finger joints for gesture-based interaction and sign language recognition~\cite{ko2019neural}. Keypoints are also defined on everyday items, industrial parts, and vehicles for manipulation, inspection, and driving (\eg, monocular 3D detection from vehicle keypoints~\cite{li2020rtm3d}), on animals for behavior analysis (multi-animal DeepLabCut~\cite{lauer2022multi}), and on anatomical structures in medical imaging, where universal landmark detectors cover several modalities and anatomies~\cite{zhu2021you}.


\begin{figure}[!t]
    \centering
    \includegraphics[width=\linewidth]{figures/part3/section6/keypoint_detection_v2.pdf}
    \caption{Keypoint and landmark detection: representative task families (top) and the methodological progression, mainly of 2D detectors, from hand-crafted feature detectors to generative models, with representative methods under each stage (bottom).}
    \label{fig:p3_s6_1_1_keypoint_overview}
\end{figure}


\subsubsection{Methodology}
Classical pipelines combined local feature detectors with statistical shape models that handled occlusion and deformation through explicit geometric constraints.
Moravec~\cite{moravec1980obstacle}, Harris~\cite{harris1988combined}, and Shi--Tomasi~\cite{shi1994good} corners were followed by feature detectors and descriptors such as SIFT~\cite{lowe2004distinctive}, SURF~\cite{bay2006surf}, FAST~\cite{rosten2006machine}, BRIEF~\cite{calonder2010brief}, and ORB~\cite{rublee2011orb}, which improved scale invariance, speed, and robustness.
For faces and hands, active shape models~\cite{cootes1995asm}, active appearance models~\cite{cootes1998active}, and constrained local models with regularized mean-shift fitting~\cite{saragih2011deformable} provided the parametric frameworks.


Learning-based detection began with supervised descent for face alignment~\cite{xiong2013supervised} and, in 2014, with DeepPose~\cite{toshev2014deeppose}, which regressed human joints directly with a deep network.
Heatmap regression then dominated through 2020: Tompson \etal~\cite{tompson2014joint,tompson2015efficient} coupled convolutional networks with graphical models, convolutional pose machines~\cite{wei2016convolutional} and stacked hourglass networks~\cite{newell2016stacked} refined heatmaps over repeated stages, and HRNet~\cite{wang2020deep} kept a high-resolution branch throughout the network for precise localization.


Multi-person estimation splits into two strategies. Top-down methods detect each person first and then estimate keypoints inside the box, which is accurate but scales with the number of people; bottom-up methods predict all keypoints at once and then group them, which is faster but makes grouping the hard part.
Part Affinity Fields~\cite{cao2017realtime}, consolidated as OpenPose~\cite{CaoHSWS21}, are the canonical bottom-up design; YOLO-Pose~\cite{maji2022yolo}, RTMPose~\cite{jiang2023rtmpose}, and RTMO~\cite{lu2024rtmo} advance real-time multi-person estimation through one-stage and top-down designs.


Vision transformers and foundation models made keypoint detection general-purpose: ViTPose~\cite{xu2022vitpose} set simple transformer baselines for human pose in 2022, and Sapiens~\cite{khirodkar2024sapiens} in 2024 provides a single foundation model for human-centric vision tasks.

3D human pose estimation grew out of these 2D detectors. Lifting methods take detected 2D keypoints as input and regress the 3D joints: a fully connected network on 2D joints alone was a strong baseline~\cite{martinez2017simple}, temporal convolutions over 2D pose sequences (VideoPose3D~\cite{pavllo2019videopose3d}) and spatio-temporal transformers (MotionBERT~\cite{zhu2023motionbert}) use motion to reduce the depth ambiguity, and DiffPose~\cite{gong2023diffpose} casts lifting as diffusion. Direct methods predict 3D from pixels, as in the volumetric heatmaps of Pavlakos \etal~\cite{pavlakos2017coarse}. With a calibrated multi-view rig most of the ambiguity disappears, and learnable triangulation~\cite{iskakov2019learnable} fuses per-view heatmaps into a 3D estimate. Recovering a full body mesh rather than a joint set is taken up with human motion in Section~\ref{sec:p3_s8_4}.


Classical corner detectors were in turn replaced by learned ones: SuperPoint~\cite{detone2018superpoint} learns detection and description jointly under self-supervised homographic adaptation, DISK~\cite{tyszkiewicz2020disk} selects keypoints with a reinforcement-style objective, and SuperGlue~\cite{sarlin2020superglue} matches sparse keypoints; the detector-free LoFTR~\cite{Sun2021LoFTR} instead produces semi-dense correspondences.
These serve as correspondence rather than semantic points and are the workhorses of the structure-from-motion, visual-odometry, and SLAM pipelines in the spatial-localization section that follows.

On point clouds, keypoints mostly serve registration and pose estimation. Intrinsic Shape Signatures~\cite{zhong2009iss} select points whose local neighborhood has a distinctive eigenvalue structure; USIP~\cite{li2019usip} learns, without labels, detectors that stay stable under rigid transformation, and D3Feat~\cite{bai2020d3feat} predicts a detection score and a descriptor for every point, the 3D analogue of SuperPoint. Category-level 3D keypoints are a different kind: Suwajanakorn \etal~\cite{suwajanakorn2018discovery} discover keypoints that are consistent across instances of a category through a relative-pose objective, and kPAM~\cite{manuelli2019kpam} uses such keypoints as the interface for robotic manipulation. Keypoint-based 6D object pose, which predicts the 2D projections of 3D model keypoints and solves PnP, belongs to pose estimation (Section~\ref{sec:object_pose}).

Table~\ref{tab:keypoint_coco_results} reports COCO Keypoints (val2017) results for the CNN, transformer, and real-time families.

\begin{table*}[!t]
\centering
\caption{Reported performance of representative keypoint detection methods on COCO Keypoints (val2017). AP is computed under the OKS-based protocol.}
\label{tab:keypoint_coco_results}
\fontsize{9}{11}\selectfont
\setlength{\tabcolsep}{2.0mm}
\begin{tabular}{l l c c c c c c c c}
\toprule
\textbf{Method} & \textbf{Backbone} & \textbf{Input} & \textbf{\#Params (M)} & \textbf{GFLOPs} & \textbf{AP} & \textbf{AP$_{50}$} & \textbf{AP$_{75}$} & \textbf{AR} & \textbf{Year} \\
\midrule
HRNet-W32~\cite{wang2020deep}  & HRNet-W32   & 256$\times$192 & 28.5  & 7.1   & 74.4 & 90.5 & 81.9 & 79.8 & 2020 \\
HRNet-W48~\cite{wang2020deep}  & HRNet-W48   & 384$\times$288 & 63.6  & 32.9  & 76.3 & 90.8 & 82.9 & 81.2 & 2020 \\
YOLOv5s6-Pose~\cite{maji2022yolo} & YOLOv5-s   & 960$\times$960 & 15.1  & --  & 63.8 & 87.6 & 69.6 & 70.4 & 2022 \\
YOLOv5m6-Pose~\cite{maji2022yolo} & YOLOv5-m & 960$\times$960 & 41.4 & -- & 67.4 & 89.1 & 73.7 & 73.9 & 2022 \\
YOLOv5l6-Pose~\cite{maji2022yolo} & YOLOv5-l & 960$\times$960 & 87.0 & -- & 69.4 & 90.2 & 76.1 & 75.9 & 2022 \\
ViTPose-B~\cite{xu2022vitpose} & ViT-B       & 256$\times$192 & 86    & --  & 75.8 & -- & -- & 81.1 & 2022 \\
ViTPose-L~\cite{xu2022vitpose} & ViT-L       & 256$\times$192 & 307   & --  & 78.3 & -- & -- & 83.5 & 2022 \\
ViTPose-H~\cite{xu2022vitpose} & ViT-H       & 256$\times$192 & 632   & -- & 79.1 & -- & -- & 84.1 & 2022 \\
RTMPose-s~\cite{jiang2023rtmpose} & CSPNeXt-s & 256$\times$192 & 5.5  & 0.68  & 69.6 & -- & -- & -- & 2023 \\
RTMPose-m~\cite{jiang2023rtmpose} & CSPNeXt-m & 256$\times$192 & 13.6  & 1.93  & 73.6 & -- & -- & -- & 2023 \\
RTMPose-l~\cite{jiang2023rtmpose} & CSPNeXt-l & 256$\times$192 & 27.7  & 4.16  & 74.8 & -- & -- & -- & 2023 \\
RTMO-s~\cite{lu2024rtmo} & CSPDarkNet & 640$\times$640 & 9.9 & -- & 67.7 & 87.8 & 73.7 & 71.5 & 2024 \\
RTMO-m~\cite{lu2024rtmo} & CSPDarkNet & 640$\times$640 & 22.6 & -- & 70.9 & 89.0 & 77.8 & 74.7 & 2024 \\
RTMO-l~\cite{lu2024rtmo} & CSPDarkNet & 640$\times$640 & 44.8 & -- & 72.4 & 89.9 & 78.8 & 76.2 & 2024 \\
\bottomrule
\end{tabular}
\end{table*}

\subsubsection{Datasets and Metrics}
Keypoint datasets are mostly embedded in broader pose and action benchmarks rather than built in isolation; Table~\ref{tab:keypoint_dataset_summary} lists the most used.
Human body datasets are the most mature and trace the field's expansion: MPII~\cite{andriluka20142d} and COCO Keypoints~\cite{lin2014microsoft} fixed in-the-wild 2D evaluation, PoseTrack~\cite{andriluka2018posetrack} added temporal association in video, Human3.6M~\cite{IonescuPOS14} supplies controlled multi-view 3D ground truth for 2D-to-3D lifting, DensePose~\cite{guler2018densepose} replaces sparse joints with dense surface correspondences, and CrowdPose~\cite{li2019crowdpose} stresses crowded, occluded scenes.
Face datasets moved from large pose variation (AFLW~\cite{koestinger2011annotated}) and high-density annotation (HELEN~\cite{le2012interactive}) to occlusion (COFW~\cite{burgos2013robust}), the standardized 68-point protocol of 300-W~\cite{sagonas2013300} and its large-pose synthetic extension 300-W-LP~\cite{zhu2016face}, and the 98-point, hard-condition WFLW~\cite{wu2018look}.
Hand datasets contend with fine articulation and self-occlusion through multi-view capture (CMU Panoptic HandDB~\cite{simon2017hand}), synthetic 3D labels (RHD~\cite{zimmermann2017learning}), real 3D pose and shape (FreiHAND~\cite{zimmermann2019freihand}), and interacting hands (InterHand2.6M~\cite{moon2020interhand2}).
Animal-Pose~\cite{cao2019cross} and AP-10K~\cite{yu2021ap} carry the human recipe to animals across species, while ApolloCar3D~\cite{song2019apollocar3d} and KeypointNet~\cite{you2020keypointnet} define keypoints on rigid vehicles and 3D object categories.

\begin{table*}[!t]
\centering
\caption{Representative datasets for keypoint and landmark detection across human body, face, hand, animal, and object domains. ``KP'' denotes the number of annotated keypoints per instance; ``Mod.'' denotes modality (2D image, 3D pose, or multi-view). Sample counts follow the official splits where available; ``--'' marks values not provided by the source paper.}
\label{tab:keypoint_dataset_summary}
\fontsize{9}{11}\selectfont
\setlength{\tabcolsep}{1.0mm}
\begin{tabular}{L{0.22\textwidth} C{0.06\textwidth} L{0.23\textwidth} C{0.11\textwidth} C{0.08\textwidth} L{0.23\textwidth}}
\toprule
\textbf{Dataset} & \textbf{Year} & \textbf{Domain} & \textbf{Mod.} & \textbf{KP} & \textbf{\#Samples (train\,/\,val\,/\,test)} \\
\midrule
\rowcolor{lightgrey} \multicolumn{6}{l}{$\bullet$ \emph{Human body}} \\
MPII~\cite{andriluka20142d}            & 2014 & In-the-wild body & 2D & 16 & 25K images (train+test) \\
COCO Keypoints~\cite{lin2014microsoft} & 2014 & In-the-wild multi-person & 2D & 17 & 57K / 5K / 20K images \\
PoseTrack2018~\cite{andriluka2018posetrack} & 2018 & Video body tracking & 2D & 15 & 593 / 170 / 375 videos \\
Human3.6M~\cite{IonescuPOS14}   & 2014 & Indoor multi-view body & 3D & 32 & 3.6M frames, 11 subjects \\
DensePose~\cite{guler2018densepose}    & 2018 & Dense body surface & 2D & SMPL surface & $\sim$50K annotated persons (COCO subset) \\
CrowdPose~\cite{li2019crowdpose}       & 2019 & Crowded multi-person & 2D & 14 & 10K / 2K / 8K images \\
\midrule
\rowcolor{lightgrey} \multicolumn{6}{l}{$\bullet$ \emph{Face}} \\
AFLW~\cite{koestinger2011annotated}    & 2011 & In-the-wild face & 2D & 21 & 25K faces (22K images) \\
HELEN~\cite{le2012interactive}         & 2012 & High-resolution face & 2D & 194 & 2K / -- / 330 images \\
COFW~\cite{burgos2013robust}           & 2013 & Occluded face & 2D & 29 & 1.3K / -- / 507 images \\
300-W~\cite{sagonas2013300}            & 2013 & Aggregated face & 2D & 68 & 3.1K / -- / 689 images \\
300-W-LP~\cite{zhu2016face}            & 2016 & Large-pose synthetic face & 2D & 68 & 61K images \\
WFLW~\cite{wu2018look}                 & 2018 & In-the-wild face (challenging) & 2D & 98 & 7.5K / -- / 2.5K images \\
\midrule
\rowcolor{lightgrey} \multicolumn{6}{l}{$\bullet$ \emph{Hand}} \\
CMU Panoptic HandDB~\cite{simon2017hand} & 2017 & Multi-view hand & Multi-view & 21 & 14K manual + bootstrapped \\
RHD~\cite{zimmermann2017learning}       & 2017 & Synthetic hand & 2D + 3D & 21 & 41K / -- / 2.7K frames \\
FreiHAND~\cite{zimmermann2019freihand}  & 2019 & Real hand pose + shape & 3D & 21 & 33K / -- / 4K samples \\
InterHand2.6M~\cite{moon2020interhand2} & 2020 & Two-hand interaction & 3D & 42 & 2.6M frames \\
\midrule
\rowcolor{lightgrey} \multicolumn{6}{l}{$\bullet$ \emph{Animal \& Object}} \\
Animal-Pose~\cite{cao2019cross}         & 2019 & 5-class mammal pose & 2D & 20 & $\sim$4K images \\
AP-10K~\cite{yu2021ap}                  & 2021 & 54-species animal pose & 2D & 17 & 10K images \\
ApolloCar3D~\cite{song2019apollocar3d} & 2019 & Vehicle keypoints (driving) & 2D & 66 & 5.3K images \\
KeypointNet~\cite{you2020keypointnet}   & 2020 & Multi-category 3D objects & 3D & 5--23 & 8.3K models \\
\bottomrule
\end{tabular}
\end{table*}

Metrics measure raw or scale-normalized distances between predicted and ground-truth points. The End-Point Error (EPE) is the raw Euclidean distance. The Normalized Mean Error (NME) divides it by a reference length:
\begin{equation}
\mathrm{NME} = \frac{1}{N} \sum_{i=1}^{N} \frac{\lVert \hat{\mathbf{x}}_i - \mathbf{x}_i \rVert_2}{d},
\end{equation}
where $\hat{\mathbf{x}}_i$ and $\mathbf{x}_i$ denote the predicted and ground-truth locations of the $i$-th keypoint, and $d$ is a task-dependent normalization factor (\eg, inter-ocular distance for faces).
The Percentage of Correct Keypoints (PCK) instead thresholds the error,
\begin{equation}
\mathrm{PCK@}\alpha = \frac{1}{N} \sum_{i=1}^{N}
\mathbb{I}\!\left( \lVert \hat{\mathbf{x}}_i - \mathbf{x}_i \rVert_2 \le \alpha \cdot s \right),
\end{equation}
where $\alpha$ controls the tolerance and $s$ is the reference scale, giving an accuracy rate that is easy to compare across methods; the Cumulative Error Distribution (CED) curve and its Area Under the Curve (AUC) summarize behavior across thresholds.

For large-scale benchmarks, Object Keypoint Similarity (OKS) incorporates object scale and keypoint-specific uncertainty,
\begin{equation}
\mathrm{OKS} =
\frac{\sum_i \exp\!\left( -\frac{\lVert \hat{\mathbf{x}}_i - \mathbf{x}_i \rVert_2^2}{2 s^2 \kappa_i^2} \right)\delta(v_i > 0)}
{\sum_i \delta(v_i > 0)}.
\end{equation}
Here $s$ denotes the object scale (the square root of the segmented area), $\kappa_i$ is a per-keypoint normalization constant calibrating the fall-off, and $v_i \in \{0, 1, 2\}$ is the ground-truth visibility flag for keypoint $i$, where $v_i = 0$ indicates unlabeled, $v_i = 1$ indicates labeled but occluded, and $v_i = 2$ indicates labeled and visible.
OKS plays the role of IoU in computing average precision and recall on MS COCO~\cite{lin2014microsoft}. For 3D human pose, the Mean Per-Joint Position Error (MPJPE) averages the Euclidean joint error in millimeters after aligning the root joints; PA-MPJPE applies a Procrustes alignment of scale, rotation, and translation first, and 3DPCK thresholds the error at $150$\,mm, as on MPI-INF-3DHP~\cite{mehta2017monocular}. Point-cloud keypoint detectors are scored by repeatability under rigid transformation and, downstream, by registration recall on benchmarks such as 3DMatch~\cite{zeng20173dmatch}.

\subsubsection{Challenges}
Robustness in unconstrained scenes is still the first problem~\cite{fisal2025reported, guo2025survey}: occlusion, extreme pose, and non-rigid deformation leave partial or ambiguous configurations, and symmetric parts (left versus right limbs, the fingers of a self-occluded hand) are routinely swapped.
Generalization is the second~\cite{kappan2026survey, ma20252d, jiang2022animal}: keypoint definitions differ across datasets even for the same object class, so models trained on one benchmark transfer poorly to another, and the gap widens for animals, objects, and anatomical structures, where annotation is scarce and demographic or species bias enters the training data.
The accuracy--speed tension discussed for detection applies here as well~\cite{khabarlak2021fast, alvarez2021real}: high-precision models are too heavy for embedded targets, and lightweight ones lose accuracy precisely under occlusion and fast motion.










\subsection{Scene Classification}
\label{sec:scene_classification}

% 最开始是简单的scene classification，即简单的分类任务，但是后续变得更加的fine-grained的任务，变成multi-label classification，再更加详细的话，就会变成scene （image/video/3D）captioning task.

Going beyond objects, regions, and keypoints, scene classification (SC) assigns a single semantic label to a whole image describing the environment it portrays~\cite{zeng2021deep}.
Where object recognition names the things in a picture (``car'', ``tree'', ``person''), scene classification names the \emph{where}: ``a busy street'', ``a serene beach'', ``a modern kitchen''.
It asks ``\emph{What kind of place is this?}'', a judgment that rests on the global layout and gist of the image rather than on any single object~\cite{henderson1999high,epstein2005cortical}, and it supplies the holistic prior on which an embodied or virtual agent conditions its later behavior~\cite{liao2016understand,ni2022improved,muhammad2022vision,Zhang2024UrbanVI}.


\subsubsection{Task Definition}

Formally, scene classification learns a mapping $f: I \mapsto y \in \mathcal{Y}$ from an image $I$ to a label $y$ in a predefined category set $\mathcal{Y}$, or, in the multi-label case, to a subset $Y \subseteq \mathcal{Y}$ covering several facets of the scene~\cite{boutell2004learning} (``urban'' for area, ``parking lot'' for place, ``night'' for condition).
Beyond the single-/multi-label split, the task is instantiated by domain, with indoor~\cite{quattoni2009recognizing} and outdoor~\cite{6724252} scenes, remote-sensing imagery~\cite{nogueira2017towards,cheng2017remote}, and driving contexts such as highways and tunnels~\cite{ni2022improved}, and by modality, extending from still images to video~\cite{huang1999integration} and acoustic scenes~\cite{mesaros2018multi}.


\subsubsection{Methodology}

\begin{figure}[!t]
    \centering
    \includegraphics[width=0.99\linewidth]{figures/part3/section6/scene_classification.pdf}
    \caption{Scene classification: four task formulations, from single-label classification to compositional captioning output (top), and the methodological progression from hand-crafted descriptors to multimodal LLMs, with representative methods under each stage (bottom).}
    \label{fig:p3_s6_1_4_scene_classification_overview}
\end{figure}

Figure~\ref{fig:p3_s6_1_4_scene_classification_overview} summarizes the progression from hand-crafted global descriptors through deep features to vision--language models.

Early methods summarized the whole image into one discriminative vector of low-level texture, color, and gradient statistics.
The Spatial Envelope model~\cite{oliva2001modeling} describes a scene by a few perceptual dimensions (naturalness, openness, roughness, expansion, ruggedness), a descriptor commonly called GIST.
Bag-of-Visual-Words (BoVW)~\cite{csurka2004visual} treats the image as a histogram of quantized local features such as SIFT~\cite{lowe2004distinctive}, discarding most spatial information~\cite{7025553}; Spatial Pyramid Matching (SPM)~\cite{lazebnik2006beyond} restores coarse layout by computing BoVW histograms over a pyramid of sub-regions and concatenating them.

Deep networks learn scene features end to end.
CNNs~\cite{krizhevsky2012imagenet,o2015introduction} such as AlexNet~\cite{krizhevsky2012imagenet}, VGG~\cite{simonyan2014very}, ResNet~\cite{he2016deep}, and DenseNet~\cite{huang2017densely}, pre-trained on ImageNet~\cite{ImageNet, krizhevsky2012imagenet}, proved to be strong global feature extractors for scenes, and Places-CNN~\cite{zhou2014learning} re-trained them on the large-scale Places benchmark to establish the standard transfer recipe.
Vision Transformers~\cite{vit} match or exceed CNNs with a class-token embedding that attends over all patches, in the manner of BERT~\cite{devlin2019bert}.


A separate line models layout explicitly, since a stadium is recognized by its global oval while a classroom differs from an office by the mid-scale arrangement of desks and blackboard.
Multi-scale pooling over the CNN feature hierarchy (FPN-style~\cite{lin2017feature}) captures structure at several scales, and scene graphs~\cite{johnson2015image} make objects and their relations explicit, so that propagating over them with graph convolutional networks~\cite{kipf2017semisupervised} yields layout-aware scene features~\cite{Liang2020ADN}.

Vision--language models (VLMs) turn scene classification into open-vocabulary, zero-shot matching.
Contrastive models such as CLIP~\cite{radford2021learning} embed the image and the textual names of all candidate classes in one space and pick the most similar class~\cite{li2023rs}, needing no fine-tuning beyond prompt design~\cite{sahoo2024systematic}; the recipe extends to 3D point clouds~\cite{huang2023clip2point}.
Autoregressive multimodal large language models (MLLMs; LLaVA~\cite{liu2023llava}, Qwen-VL~\cite{bai2023qwen}, Gemma~3~\cite{team2025gemma}) instead generate the label as text in response to an instruction~\cite{Saha_2024_CVPR}.
With rich semantics available for free, the research focus has shifted from one global label toward compositional understanding of what is inside the scene~\cite{chen2023clip2scene}.

\subsubsection{Datasets and Metrics}
Table~\ref{tab:sceneclf_datasets} lists the standard benchmarks. The small early sets, Scene 15~\cite{lazebnik2006beyond} and the indoor MIT 67~\cite{quattoni2009recognizing}, were superseded by SUN 397~\cite{XiaoHEOT10,Xiao2014SUNDE}, which fixed a near-exhaustive category list, and by Places~\cite{zhou2014learning, zhou2017places}, which scaled scene recognition to millions of images across its Places205, Places88, Places365-Standard, and Places365-Challenge subsets. NYU Depth v2~\cite{silberman2012nyudv2} and SUN RGB-D~\cite{song2015sunrgbd} add aligned depth and a scene category per image for indoor RGB-D classification.

\begin{table*}[!t]
\centering
\caption{Representative scene classification datasets. ``Mod.'' denotes modality; ``--'' marks a value not fixed by the source (scene-category subsets vary across protocols).}
\label{tab:sceneclf_datasets}
\fontsize{9}{11}\selectfont
\setlength{\tabcolsep}{2.2mm}
\begin{tabular}{l c l c c l}
\toprule
\textbf{Dataset} & \textbf{Year} & \textbf{Domain} & \textbf{Mod.} & \textbf{\#Classes} & \textbf{\#Images (protocol)} \\
\midrule
Scene 15~\cite{lazebnik2006beyond}                    & 2006 & Indoor + outdoor, grayscale & 2D & 15  & 4,448 \\
MIT Indoor 67~\cite{quattoni2009recognizing}               & 2009 & Indoor                      & 2D & 67  & 15,620 (80 / 20 per class) \\
SUN 397~\cite{XiaoHEOT10,Xiao2014SUNDE}       & 2010 & Indoor + outdoor            & 2D & 397 & 108,754 (50 test per class) \\
Places205~\cite{zhou2014learning}          & 2014 & Indoor + outdoor            & 2D & 205 & 2.5M \\
Places365-Standard~\cite{zhou2017places}   & 2017 & Indoor + outdoor            & 2D & 365 & $\sim$1.8M \\
Places365-Challenge~\cite{zhou2017places}  & 2017 & Indoor + outdoor            & 2D & 365 & $\sim$8M \\
NYU Depth v2~\cite{silberman2012nyudv2}    & 2012 & Indoor                      & RGB-D & -- & 1,449 labeled pairs \\
SUN RGB-D~\cite{song2015sunrgbd}           & 2015 & Indoor                      & RGB-D & -- & 10,335 \\
\bottomrule
\end{tabular}
\end{table*}

Scene classification is evaluated primarily with top-1 accuracy (the fraction of images for which the highest-scoring predicted class matches the ground truth) and top-5 accuracy (the fraction of images where the ground-truth class is among the top five predictions), with the latter especially useful on long-tail benchmarks such as Places-365 and SUN 397.
For multi-label settings, the mean Average Precision (mAP) and class-wise F1 are reported instead, mirroring the protocol used in multi-label image recognition.

\begin{table*}[!t]
\centering
\caption{Reported performance of representative scene classification methods on Places-365 val and SUN 397. ``--'' marks a metric not reported by the corresponding source.}
\label{tab:sceneclf_results}
\fontsize{9}{11}\selectfont
\setlength{\tabcolsep}{5.4mm}
\begin{tabular}{l l c c c}
\toprule
\textbf{Method} & \textbf{Backbone} & \textbf{Places-365 Top-1} & \textbf{SUN 397 Top-1} & \textbf{Year} \\
\midrule
\rowcolor{lightgrey} \multicolumn{5}{l}{$\bullet$ \emph{Supervised}} \\
Places365-AlexNet~\cite{zhou2017places}    & AlexNet   & 53.2 & --   & 2017 \\
Places365-VGG~\cite{zhou2017places}        & VGG-16    & 55.2 & --   & 2017 \\
Places365-ResNet~\cite{zhou2017places}     & R152      & 54.7 & --   & 2017 \\
\midrule
\rowcolor{lightgrey} \multicolumn{5}{l}{$\bullet$ \emph{Vision-language zero-shot}} \\
CLIP~\cite{radford2021learning}            & ViT-B/32  & --   & 63.2 & 2021 \\
\bottomrule
\end{tabular}
\end{table*}

Table~\ref{tab:sceneclf_results} reports results on the two most-cited scene benchmarks. Supervised CNNs on Places-365 cluster in a narrow band of roughly $53$--$55\%$ top-1, with diminishing returns from AlexNet to ResNet, while CLIP reaches $63.2\%$ top-1 on SUN 397 zero-shot~\cite{radford2021learning}, on par with supervised scene classifiers without any target-set training.




\subsubsection{Challenges}
Closed-set accuracy on Places-365, SUN 397, and MIT 67 is now recovered zero-shot by CLIP~\cite{radford2021learning} and by MLLMs (LLaVA~\cite{liu2023llava}, Qwen-VL~\cite{bai2023qwen}, Gemma~3~\cite{team2025gemma}), so ``what kind of place is this'' is no longer the frontier.
The task has climbed a granularity ladder, from one label to attribute sets to grounded, compositional description across image, video, and 3D, and the open problems sit at that richer end, where MLLMs are weaker than their coarse-label accuracy suggests.

\textbf{Fine-grained and long-tail discrimination.} MLLMs still trail dedicated classifiers and CLIP-style models on fine-grained and long-tail categories, where the distinction between, \eg, a ``bistro'' and a ``brasserie'' turns on subtle cues~\cite{NEURIPS2024_5c702404}.

\textbf{Compositional and spatial grounding.} The natural output is no longer a label but a description whose claims can be grounded in specific objects and layouts~\cite{chen2023clip2scene}.
MLLMs describe fluently but ground weakly: objects are referenced vaguely rather than by region or box, and the persistent failure is structural rather than semantic, since models name the constituents of a scene correctly yet err on relative layout, depth ordering, orientation, and cross-view or 3D consistency.
The cause is representational: 2D patch tokens are optimized for semantic alignment with language, not geometry, so fine spatial cues are lost once flattened into the language stream, and spatial-cognition benchmarks show even strong MLLMs degrading sharply as tasks move from naming toward layout and 3D reasoning~\cite{xu2025spatialbench,daxberger2025mmspatial}.

\textbf{Hallucination and faithfulness.} Free-form description reintroduces a failure mode that closed-set classification never had: MLLMs assert objects, attributes, or relations absent from the image, and their reasoning traces rarely point to verifiable spatial evidence~\cite{daxberger2025mmspatial}.






% \newpage

\section{Spatial Localization and Position Estimation}
\label{sec:p3-sec6-2-spatial-localization-and-position-estimation}

\FramedInitial{R}{ecognition} says what is present; it does not say where. Spatial localization and position estimation supply the missing geometry, situating recognized entities in a metric frame and recovering how they are oriented within it. The section climbs a ladder of geometric abstraction: from the per-pixel quantities of depth and surface normals, through metric attributes such as height and distance, to full SE(3) pose, first of the moving camera and then of rigid objects in the scene. Throughout, the concern is the static geometry of a single instant; how these quantities evolve is deferred to spatio-temporal perception.


\subsection{Depth Estimation}
\label{sec:depth_estimation}
Depth estimation predicts, for every pixel, the distance from the camera to the surface it observes, turning a 2D image into the 3D structure that reconstruction, navigation, and manipulation build on. Figure~\ref{fig:p3_s6_2_1_depth_overview} summarizes the input configurations and the method families surveyed below.

\subsubsection{Task Definition}
Given an image $I$ and optional auxiliary observations $\mathcal{V}$ (a second view, a multi-view set, a sparse sensor measurement, or neighboring video frames), a depth estimator predicts a depth map
\begin{equation}
f_\theta(I,\mathcal{V}) = D \in \mathbb{R}_{+}^{H\times W},
\end{equation}
where $D(u)$ is the distance along the optical axis from the camera center to the surface seen at pixel $u$. The task splits along three axes.

\textbf{Sensing.} Active methods emit light (structured light, time-of-flight, LiDAR) and measure its return, providing reliable metric depth even on textureless surfaces but tying accuracy to the sensor and degrading under strong ambient light or at long range. Passive methods infer depth from ordinary images alone; they need no special hardware and have become the dominant research line, at the cost of sensitivity to texture and illumination.

\textbf{Scale.} Metric depth gives distances in physical units, which 3D reconstruction and navigation require. Relative depth recovers only pixel ordering, or depth up to an affine transform; it cannot support metric measurement but is easier to supervise and generalizes better because heterogeneous datasets can be mixed without a consistent scale.

\textbf{Density and Input Configuration.} Classical pipelines produced sparse depth at feature points or semi-dense depth along high-gradient edges; dense per-pixel prediction is now the norm. The input configuration fixes how much geometry is available: active sensing measures depth directly; stereo recovers it from the disparity between two calibrated views; multi-view estimation enforces consistency across many views, which improves robustness to occlusion and matching ambiguity at the price of computation and known camera poses; monocular estimation lacks direct multi-view triangulation constraints and relies on learned statistical, structural, and semantic cues, which is why it is ill-posed in scale yet the most widely applicable setting.

% \paragraph{Extend Tasks}

%  \item Depth Completion (RGB + sparse depth)
%  \item Depth Upsampling
%  \item Depth Refinement
%  \item Depth Fusion

\subsubsection{Methodology}
% 深度估计方法整体脉络可整理为早期依据几何方法直接求解三维的深度数值，而后进入深度学习领域，从早期的有监督方法发展至无监督和自监督方法，其中又由传统的CNN网络转变为更强大的Transformer架构。目前最新的深度估计也在探索具有更好Zero-Shot表现、泛化性能更强的方法。

\begin{figure}[!t]
    \centering
    \includegraphics[width=\linewidth]{figures/part3/section6/depth_estimation_v3.pdf}
    \caption{Depth estimation: four input configurations (top) and five overlapping method families with representative methods on an approximate emergence timeline (bottom). Cost volumes and attention are distinct mechanisms; denoising applies to the generative branch. The depth maps are schematic illustrations, with brighter values indicating nearer surfaces.}
    \label{fig:p3_s6_2_1_depth_overview}
\end{figure}

\paragraph{Geometric Methods.}
Classical depth estimation solves geometry explicitly. In binocular stereo, two rectified cameras with focal length $f$ and baseline $B$ observe a point whose projections differ by the disparity $d = x_1 - x_2$; similar triangles give
\begin{equation}
Z = \frac{fB}{d},
\end{equation}
so depth follows once correspondences are found. Stereo matching is therefore a correspondence problem, usually posed as an energy that balances data fidelity against smoothness; Semi-Global Matching approximates the global optimum by aggregating costs along several 1D paths and remains a standard dense baseline. Structure from Motion (SfM) extends triangulation to many uncalibrated views, recovering camera poses and a sparse point cloud up to global scale, and supplies the poses that Multi-View Stereo (MVS) densifies. MVS typically proceeds by plane sweep: for a reference pixel $u_r$ and a hypothesized depth $d_k$, back-project with the intrinsics $K_r$, transform into a source view with the relative pose $(R, t)$, reproject, and keep the depth whose reprojections are most photometrically consistent across source views,
\begin{equation}
u_s \sim K_s\left(R\, d_k K_r^{-1} u_r + t\right), \qquad
D(u_r) = \arg\min_{d_k} \sum_s \left\| I_s(u_s) - I_r(u_r) \right\|.
\end{equation}
Plane sweep is thus triangulation turned into a search over depth planes, and it is the template that MVSNet~\cite{yao2018mvsnet} and CasMVSNet~\cite{gu2020cascade} later made differentiable as learned cost volumes.

Two single-image cues complete the classical picture. Depth from defocus exploits the thin-lens relation $1/f = 1/z + 1/z'$: a point off the focal plane blurs into a circle of confusion whose diameter grows with its distance from that plane, so measured blur can be inverted for depth. Shape from shading inverts a reflectance model, typically Lambertian, to recover surface orientation from intensity and integrates it into depth; it is treated under surface normals in \S\ref{sec:surface_normal}. Both are ill-posed under unknown illumination, reflectance, or optics and survive today mainly as auxiliary cues.

\paragraph{Active Depth Sensing.}
Active sensors measure depth rather than infer it. Time-of-flight systems time the round trip of emitted light: LiDAR scans the scene with laser beams and returns a sparse but long-range point cloud, whereas ToF cameras flood the scene and recover a dense depth map at every pixel, best suited to short-range indoor use. Structured light projects a coded pattern (the Kinect v1 used a pseudo-random infrared speckle) and triangulates from its deformation; active stereo projects texture onto the scene only to make ordinary stereo matching stable on featureless surfaces.

\paragraph{Early Learning-Based Methods.}
% 早期的机器学习方法
Data-driven depth began in the mid-2000s with monocular estimation. Saxena \etal~\cite{saxena2005learning} regressed depth from hand-crafted multi-scale features with a Markov random field enforcing spatial coherence, and Make3D~\cite{saxena2008make3d} extended the model to the 3D position and orientation of scene patches while releasing a laser-scanned outdoor dataset that served as an early benchmark.

\paragraph{Supervised CNN Regression.}
Eigen \etal~\cite{eigen2014depth} first applied CNNs to monocular depth with a coarse-to-fine two-scale network, and their evaluation protocol became the standard; a follow-up~\cite{eigen2015predicting} showed that the same architecture predicts depth, normals, and semantic labels with minor changes. Deep Convolutional Neural Fields~\cite{liu2015deep} paired CNN unaries with a continuous CRF, FCRN~\cite{laina2016deeper} introduced a fully convolutional encoder--decoder with efficient upsampling of ResNet features, and DORN~\cite{fu2018deep} recast depth as ordinal regression over a spacing-increasing discretization.

\paragraph{Learned Stereo and Multi-View Stereo.}
Learning entered the geometric pipelines first as a replacement for individual components: MC-CNN~\cite{zbontar2015computing} computed stereo matching costs with a Siamese network in place of hand-crafted similarity. Mayer \etal~\cite{mayer2016large} then showed that networks can regress disparity and scene flow end to end, and GC-Net~\cite{kendall2017end} and PSMNet~\cite{chang2018pyramid} embedded the cost volume and its regularization inside the network. The same move reached multi-view geometry with SurfaceNet~\cite{ji2017surfacenet} and MVSNet~\cite{yao2018mvsnet}, which infer depth from several posed images, and R-MVSNet~\cite{yao2019recurrent} and CasMVSNet~\cite{gu2020cascade} improved their memory footprint and resolution.

\paragraph{Self-Supervised and Multi-Task Learning.}
Ground-truth depth is costly, so a second line replaced it with geometric consistency. SfM-Net~\cite{vijayanarasimhan2017sfm} and SfM-Learner~\cite{zhou2017unsupervised} learn depth and camera motion jointly from unlabeled video by minimizing photometric reprojection error between adjacent frames; Monodepth~\cite{godard2017unsupervised} obtains the same signal from stereo pairs, and Monodepth2~\cite{godard2019digging} unifies both sources in one framework. GeoNet~\cite{yin2018geonet} adds optical flow to the joint estimate. Because all of these train a monocular network with stereo or multi-view supervision, the old separation between input configurations began to dissolve. Alongside, multi-task sharing (depth with semantic segmentation~\cite{wang2015towards}), ImageNet-pretrained encoders with structural losses (DenseDepth~\cite{alhashim2018high,deng2009imagenet,huang2017densely}), local planar guidance (BTS~\cite{lee2019big}), and mixed-dataset training for zero-shot transfer (MiDaS~\cite{MiDaS}) each pushed accuracy and generalization.

\paragraph{Transformers and Foundation Models.}
Self-attention~\cite{attention} gives every pixel access to global context, which helps precisely where local convolutions fail: textureless regions, occlusions, and large scenes. DPT~\cite{dpt} established the pattern of a Vision Transformer~\cite{vit} encoder with a light dense decoder, and DepthFormer~\cite{li2023depthformer} combined global attention with local convolutional refinement. Scaling data and model size then produced depth foundation models that are trained once and transferred zero-shot: ZoeDepth~\cite{zoedepth} and Depth Anything~\cite{DepthAnything} learn from large mixed corpora with flexible supervision and generalize across scenes and capture conditions without per-dataset adaptation.

\paragraph{Video and Generative Depth.}
Video adds temporal constraints. ManyDepth~\cite{watson2021temporal} aggregates neighboring frames into a cost volume at test time, whereas MonoViT~\cite{zhao2022monovit} keeps the self-supervised video objective but strengthens single-frame reasoning with a Transformer. A separate line treats depth as conditional generation: VPD~\cite{zhao2023unleashing} showed that text-to-image diffusion backbones carry usable geometric features, Marigold~\cite{ke2024repurposing} reframed monocular depth as conditional denoising and obtained globally coherent maps under hard conditions, and GeoWizard~\cite{fu2024geowizard} extended the recipe to joint depth and normal prediction. The same diffusion-as-perception idea recurs for surface normals and object pose later in this section.

\paragraph{Event-Based Depth.}
% 早期几何驱动 → 学习驱动方法和自监督硬件优化方法 → 综述 (Event-based Stereo Depth Estimation: A Survey 2024)
Event cameras record asynchronous per-pixel brightness changes~\cite{mahowald1992vlsi, lichtsteiner2008128}, offering microsecond latency and high dynamic range where frame cameras fail. Early event-based depth relied on geometric constraints and temporal accumulation~\cite{rebecq2018emvs, ieng2018neuromorphic,zhou2018semi, xie2017event}; learned event stereo~\cite{tulyakov2019learning, zhang2022discrete}, event--image fusion~\cite{mostafavi2021event, ghosh2022multi}, and self-supervised or hardware-aware frameworks~\cite{cho2023learning, de2025derd, cho2024temporal} followed. The line remains a niche complement for extreme dynamics; see the survey~\cite{ghosh2025event}.

\subsubsection{Datasets and Metrics}
Table~\ref{tab:depth_estimation_datasets} groups the main datasets into Type I, which ship explicit depth maps for direct supervision, and Type II, which require depth to be derived from LiDAR, stereo, or SfM. Three waves are visible. Indoor RGB-D collections (NYU Depth v2~\cite{silberman2012nyudv2}, TUM RGB-D~\cite{Sturm2012TUMRGBD}, SUN RGB-D~\cite{song2015sunrgbd}, later ScanNet~\cite{dai2017scannet} and Matterport3D~\cite{chang2017matterport3d}) and the LiDAR-based KITTI~\cite{Geiger2013KITTI,uhrig2017sparsity} fixed the indoor and outdoor protocols that still anchor the field, with Make3D~\cite{saxena2008make3d} as the laser-scanned precursor. Relative-depth supervision from internet photos (DIW~\cite{chen2016single}, ReDWeb~\cite{xian2018monocular}, MegaDepth~\cite{li2018megadepth}) and synthetic data (Scene Flow~\cite{mayer2016large}) then relaxed the need for metric sensors, while ETH3D~\cite{Schoeps2017ETH3D}, iBims-1~\cite{koch2018evaluation}, ApolloScape~\cite{wang2019apolloscape}, and Cityscapes stereo pairs~\cite{cordts2016cityscapes} served precision evaluation and self-supervised training. The current wave is driven by autonomous driving and rendering: multi-sensor fleets (nuScenes~\cite{caesar2020nuscenes}, Waymo Open~\cite{WaymoOpen-SunKDCPTGZCCVHN20}, Argoverse~2~\cite{wilson2023argoverse}, KITTI-360~\cite{liao2022kitti360}, DDAD~\cite{guizilini20203d}) supply millions of LiDAR-calibrated frames, photorealistic synthetic sets (Hypersim~\cite{roberts2021hypersim}, Virtual KITTI~2~\cite{Cabon2020VirtualKITTI2}, UrbanSyn~\cite{UrbanSyn}, TartanAir~\cite{TartanAir}) supply exact ground truth even for transparent surfaces and hard lighting, and DIODE~\cite{vasiljevic2019diode}, BlendedMVS~\cite{yao2020blendedmvs}, ARKitScenes~\cite{baruch2021arkitscenes}, Mannequin~\cite{li2019learning}, and WildRGB-D~\cite{xia2024rgbd} cover high-precision scans, handheld LiDAR, and dynamic or cluttered capture.

\begin{table}[!t]
\centering
\caption{Depth estimation datasets grouped by supervision type: Type~I datasets provide explicit depth maps for direct supervision; Type~II datasets require depth to be derived from raw LiDAR, stereo pairs, or SfM, often through projection or self-supervision.}
\label{tab:depth_estimation_datasets}
\fontsize{9}{11}\selectfont
\setlength{\tabcolsep}{2.5mm}
\begin{tabular}{lccccccr}
\toprule
\textbf{Name} & \textbf{Year} & \textbf{Scene} & \textbf{Type} & \textbf{Metric} & \textbf{Setup} & \textbf{Volume} & \textbf{Data Source} \\
\midrule
\rowcolor{lightgrey} \multicolumn{8}{l}{$\bullet$ \textit{Type I: Datasets with Explicit Depth Maps (Direct Supervision})} \\
Make3D~\cite{saxena2008make3d} & 2009 & Outdoor & Real & Metric & Mono & 534 & Laser Scanner \\
NYU Depth v2~\cite{silberman2012nyudv2} & 2012 & Indoor & Real & Metric & Mono & 1.4K & Kinect (SL) \\
TUM RGB-D~\cite{Sturm2012TUMRGBD} & 2012 & Indoor & Real & Metric & Mono & 39 Seqs & Kinect (SL) \\
SUN RGB-D~\cite{song2015sunrgbd} & 2015 & Indoor & Real & Metric & Mono & 10K & RGB-D Cameras \\
DIW~\cite{chen2016single} & 2016 & Varied & Real & Relative & Mono & 495K & Internet Photos \\
KITTI (Depth)~\cite{uhrig2017sparsity} & 2017 & Outdoor & Real & Metric & Stereo & 93K & LiDAR (Velodyne) \\
ScanNet~\cite{dai2017scannet} & 2017 & Indoor & Real & Metric & Mono & 2.5M & Structure Sensor \\
Matterport3D~\cite{chang2017matterport3d} & 2017 & Indoor & Real & Metric & Mono & 194K & Matterport (RGB-D) \\
ETH3D~\cite{Schoeps2017ETH3D} & 2017 & Both & Real & Metric & Multi-V & 13 Scans & Laser Scanner \\
ApolloScape~\cite{wang2019apolloscape} & 2018 & Outdoor & Real & Metric & Stereo & 140K & LiDAR + Stereo \\
iBims-1~\cite{koch2018evaluation} & 2018 & Indoor & Real & Metric & Mono & 100 & Camera + Laser \\
MegaDepth~\cite{li2018megadepth} & 2018 & Outdoor & Real & Relative & Mono & 130K & Internet + SfM \\
ReDWeb~\cite{xian2018monocular} & 2018 & Varied & Real & Relative & Mono & 3.6K & Internet + Stereo \\
DIODE~\cite{vasiljevic2019diode} & 2019 & Both & Real & Metric & Mono & 26K & Camera + LiDAR \\
BlendedMVS~\cite{yao2020blendedmvs} & 2020 & Both & Syn+Real & Metric & Multi-V & 17K & MVS Reconstruction \\
Hypersim~\cite{roberts2021hypersim} & 2020 & Indoor & Syn & Metric & Mono & 77K & Simulation \\
DDAD~\cite{guizilini20203d} & 2020 & Outdoor & Real & Metric & Multi-V & 76K & HD LiDAR \\
Virtual KITTI 2~\cite{Cabon2020VirtualKITTI2} & 2020 & Outdoor & Syn & Metric & Stereo & 21K & Simulation (Unity) \\
TartanAir~\cite{TartanAir} & 2020 & Varied & Syn & Metric & Stereo & $\sim$1M & Simulation (UE4) \\
ARKitScenes~\cite{baruch2021arkitscenes} & 2021 & Indoor & Real & Metric & Mono & 5K Clips & iPad LiDAR \\
UrbanSyn~\cite{UrbanSyn} & 2023 & Outdoor & Syn & Metric & Mono & 7.5K & Simulation \\
WildRGB-D~\cite{xia2024rgbd} & 2024 & Varied & Real & Metric & Mono & 20K & RGB-D Cameras \\
\midrule
\rowcolor{lightgrey} \multicolumn{8}{l}{$\bullet$ \textit{Type II: Datasets Requiring Depth to be Derived from Raw Geometry (LiDAR, Stereo or SfM)}} \\
KITTI (Raw)~\cite{Geiger2013KITTI} & 2012 & Outdoor & Real & Metric & Stereo & 42K & LiDAR (Velodyne) \\
Cityscapes~\cite{cordts2016cityscapes} & 2016 & Outdoor & Real & Metric & Stereo & 25K & Stereo Camera \\
Scene Flow~\cite{mayer2016large} & 2016 & Varied & Syn & Metric & Stereo & 39K & Simulation (Blender) \\
NuScenes~\cite{caesar2020nuscenes} & 2019 & Outdoor & Real & Metric & Multi-V & 1.4M & LiDAR \\
Waymo Open~\cite{WaymoOpen-SunKDCPTGZCCVHN20} & 2019 & Outdoor & Real & Metric & Multi-V & 200K & LiDAR \\
Mannequin~\cite{li2019learning} & 2019 & Varied & Real & Relative & Mono & 170K & Internet + SfM \\
KITTI-360~\cite{liao2022kitti360} & 2020 & Outdoor & Real & Metric & Stereo & 100K+ & LiDAR + Stereo \\
Argoverse 2~\cite{wilson2023argoverse} & 2022 & Outdoor & Real & Metric & Multi-V & 6M & LiDAR \\
\bottomrule
\end{tabular}

\end{table}

% 经典指标：RMSE、RMSE Log；强调相对偏差时：Abs Rel、Sq Rel；针对尺度模糊、引入不变性：Scale-Invariant Log Error；特定于深度估计，评估相对准确率：δ Threshold Accuracy；相对误差/排序指标主要用于相对深度估计的深度顺序关系
Depth is evaluated with pixel-wise regression metrics. The root mean squared error and its logarithmic variant,
\begin{equation}
\mathrm{RMSE} = \sqrt{ \frac{1}{N} \sum_{i=1}^{N} \left( \hat{d}_i - d_i \right)^2 }, \qquad
\mathrm{RMSE}_{\log} = \sqrt{ \frac{1}{N} \sum_{i=1}^{N} \left( \log \hat{d}_i - \log d_i \right)^2 },
\end{equation}
measure absolute deviation, the log form damping the dominance of far-range errors. Because an error of fixed size matters more at close range, relative errors normalize by the true depth,
\begin{equation}
\mathrm{Abs\ Rel} = \frac{1}{N} \sum_{i=1}^{N} \frac{ \left| \hat{d}_i - d_i \right| }{ d_i }, \qquad
\mathrm{Sq\ Rel} = \frac{1}{N} \sum_{i=1}^{N} \frac{ \left( \hat{d}_i - d_i \right)^2 }{ d_i },
\end{equation}
and Abs Rel is the primary metric on most benchmarks. For monocular prediction, whose scale is ambiguous, the scale-invariant log error discounts a global scale shift,
\begin{equation}
\mathrm{SI\text{-}Log} =
\frac{1}{N} \sum_{i=1}^{N} \left( \log \hat{d}_i - \log d_i \right)^2
-
\frac{1}{N^2} \left( \sum_{i=1}^{N} \left( \log \hat{d}_i - \log d_i \right) \right)^2 .
\end{equation}
Threshold accuracy reports the fraction of pixels within a multiplicative tolerance,
\begin{equation}
\delta =
\frac{1}{N}
\left|
\left\{
i \;\middle|\;
\max\!\left(
\frac{\hat{d}_i}{d_i},
\frac{d_i}{\hat{d}_i}
\right)
< \tau
\right\}
\right|, \qquad \tau \in \{1.25,\, 1.25^2,\, 1.25^3\}.
\end{equation}
Ordinal metrics, which score only the pairwise depth ordering, are used where supervision is relative; they track perceived layout but say nothing about metric precision.

On KITTI, supervised regressors such as BTS~\cite{lee2019big} and DPT~\cite{dpt} reach Abs Rel around $0.06$ with $\delta_{1.25}$ near $0.96$, while the self-supervised Monodepth2~\cite{godard2019digging} reaches Abs Rel near $0.10$ without any ground-truth depth. On NYU Depth v2, ZoeDepth~\cite{zoedepth} and Depth Anything~\cite{DepthAnything} reach $\delta_{1.25}$ of $0.96$ and $0.98$ after fine-tuning, and UniDepth~\cite{UniDepth} matches them zero-shot, which makes transfer a credible alternative to per-benchmark training. Generative models such as Marigold~\cite{ke2024repurposing} match or exceed regression on scale-invariant metrics while producing sharper boundaries, so recent gains come more from representation and data scale than from architectural depth. Boundary sharpness, temporal stability, and latency are still not evaluated in a standard way.

\subsubsection{Challenges}
Monocular depth is ill-posed: distinct 3D scenes project to the same image, so a single view alone does not uniquely determine 3D structure or metric scale and predictions drift across scenes. Scene content adds its own failures. Moving objects break the static-world assumption behind self-supervised and multi-view methods, and in video the result is flicker between frames and scale drift over time. Transparent and reflective materials and textureless regions give weak or misleading appearance cues; thin structures and sharp boundaries lose detail; low light, adverse weather, and underwater imagery degrade models trained on clean data. These scene effects compound into a domain gap: models are tied to the geometry and depth distribution of their training data, the indoor--outdoor divide is especially wide, real depth labels are scarce and noisy, and synthetic data still transfers imperfectly.

Current work attacks these problems from several sides. Depth foundation models~\cite{DepthAnything,UniDepth,zoedepth} supply strong zero-shot priors, and Video Depth Anything~\cite{chen2025video} adds temporal modules to one of them for consistent video depth, while video depth diffusion (DepthCrafter~\cite{hu2025depthcrafter}, ChronoDepth~\cite{shao2025learning}, RollingDepth~\cite{ke2025video}) pursues the same goal generatively. Coupling depth with 3D scene representations (DepthSplat~\cite{xu2024depthsplat}) improves spatial consistency, language cues help anchor metric scale~\cite{zeng2024rsa}, and lightweight architectures such as MonoLENS~\cite{higashiuchi2025monolens} and LMDepth~\cite{long2025lmdepth} target the real-time budgets of driving, robotics, and AR/VR.





\subsection{Surface Normal Estimation}
\label{sec:surface_normal}

Surface normal estimation predicts the local orientation of visible surfaces from images or videos.
Where depth answers ``how far is this point from the camera?'', a normal answers ``which way does the surface face?''.
That makes normals a compact, scale-invariant description of local 3D geometry: a normal field exposes planes, folds, object boundaries, and contact surfaces even when metric depth is ambiguous, which is why it serves reconstruction~\cite{wang2021neus}, relighting and intrinsic decomposition~\cite{munkberg2022extracting,barron2015shape}, physical interaction modeling~\cite{Bear2021Physion}, and egocentric perception~\cite{grauman2022ego4d}.

The learning problem looks simple: from an RGB image, RGB-D frame, multi-view set, or video, output one unit vector per valid pixel.
The inverse problem underneath is badly constrained, since the same intensity can be explained by different combinations of shape, reflectance, and illumination, and camera intrinsics change how image gradients map to 3D orientation.
Progress has therefore come from three directions at once: more and better supervision, stronger geometric inductive biases, and visual priors inherited from large dense-prediction and diffusion models; since 2024 the last of these has turned normal estimation into one output head of broader geometric foundation models.



\subsubsection{Task Definition}
Let $I\in\mathbb{R}^{H\times W\times 3}$ be an image, $\mathcal{V}$ optional auxiliary observations (neighboring frames, stereo or multi-view images, depth, intrinsics, or poses), and $\Omega$ the set of valid pixels.
The goal is to estimate
\begin{equation}
    f_\theta(I,\mathcal{V}) = \hat{\mathbf{N}}, \qquad
    \hat{\mathbf{N}}(u) = \hat{\mathbf{n}}_u \in \mathbb{S}^2,\quad u\in\Omega,
\end{equation}
where each $\hat{\mathbf{n}}_u$ is a unit vector in the camera frame or in a gravity- or world-aligned frame.
Supervision comes from depth sensors, fitted meshes, rendered assets, multi-view reconstruction, or indirectly from photometric and intrinsic-image losses, and camera metadata matters because a pixel's viewing ray depends on focal length and field of view.
Monocular prediction from one RGB image is the canonical benchmark setting; RGB-D variants use measured or predicted depth as input or supervision through the differential relation between depth gradients and normals, multi-view and video variants enforce viewpoint or temporal consistency, object-centric variants target CAD-like objects, hands, faces, and bodies, and panoramic variants predict under equirectangular or fisheye geometry, where pinhole assumptions break.


\begin{figure}[!t]
    \centering
    \includegraphics[width=\linewidth]{figures/part3/section6/surface_normal_estimation.pdf}
    \caption{Surface normal estimation. \textit{Top}: an RGB, video, or multi-view input is mapped by a dense prediction network ((a) ViT-like or (b) diffusion variant) to a per-pixel unit normal field $\hat{\mathbf{n}}\in\mathbb{S}^2$, scored by angular error. \textit{Bottom}: method timelines along three threads: joint depth--normal prediction, intrinsic/shading decomposition, and backbone architecture.}
    \label{fig:p3_s6_2_3_normal_overview}
\end{figure}

\subsubsection{Methodology}
Surface normal estimation has moved from physics-based inverse rendering to supervised dense prediction and, most recently, to foundation-model and diffusion-based geometry estimation.
Figure~\ref{fig:p3_s6_2_3_normal_overview} summarizes the method families; the last two below are grouped as one family in the figure.

\paragraph{Photometric and intrinsic-decomposition methods.}
The classical origin is shape from shading and photometric stereo.
Photometric stereo recovers orientation from several images taken under different lighting; Woodham's formulation~\cite{woodham1980photometric} shows that normals follow once illumination and reflectance are sufficiently constrained.
Single-image shape from shading is harder because a change in intensity may signal a change in orientation, albedo, illumination, or material; SIRFS~\cite{barron2015shape} is the canonical optimization-era method, jointly recovering shape, illumination, and reflectance under natural-image and geometric priors.
These methods identify the physical ambiguities that every later learned model inherits, and physically based rendering datasets such as Hypersim~\cite{roberts2021hypersim}, where shading, materials, depth, and normals are jointly available, keep their logic in play.

\paragraph{Supervised CNN and multi-task dense prediction.}
Eigen and Fergus~\cite{eigen2015predicting} trained one multi-scale convolutional network to predict depth, normals, and semantic labels, establishing that normals can be regressed directly from RGB and that geometric tasks benefit from shared representations.
Later CNN systems improved encoder capacity, decoder resolution, and boundary preservation, usually with normals as one branch of a multi-task perception model, and Taskonomy~\cite{zamir2018taskonomy} and Omnidata~\cite{eftekhar2021omnidata} systematized the paradigm by training dense predictors across many 3D scans and tasks.
A CNN normal estimator runs in one forward pass and drops into real-time reconstruction, AR, or robotics pipelines; its local features, however, blur thin structures and lack the global context needed for consistent orientation across large planes.

\paragraph{Transformer and camera-aware discriminative predictors.}
Vision Transformers let every output pixel attend to the whole image.
DPT~\cite{dpt} showed that Transformer features support strong dense prediction across depth and semantics, and the same global context helps normals: co-planar walls, floors, and object faces may be far apart in the image yet geometrically coupled.
Pure attention discards one geometric fact, that pixels correspond to camera rays; DSINE~\cite{bae2024dsine} injects ray-direction information and shows that camera geometry remains valuable even for large learned backbones.

\paragraph{Diffusion and foundation-model priors.}
A more recent line casts normal estimation as conditional generation or as one head of a general geometric foundation model.
VPD~\cite{zhao2023unleashing} showed that text-to-image diffusion models carry useful features for dense perception, and GeoWizard~\cite{fu2024geowizard} extended the Marigold recipe of \S\ref{sec:depth_estimation} to joint depth and normal prediction under a shared generative prior.
Metric3Dv2~\cite{hu2024metric3dv2} took the discriminative route, training a camera-aware monocular model for metric depth and normals across heterogeneous datasets.
StableNormal~\cite{ye2024stablenormal} addressed a practical weakness of diffusion-based normals: stochastic denoising yields high-quality but high-variance maps, so reducing variance matters as much as improving the mean.
The generative formulation samples uncertain or textureless regions from a learned distribution instead of forcing a single local regression, at the cost of iterative-denoising latency, which motivates distillation and single-step variants.

\paragraph{Temporal, multi-view, and reconstruction-coupled methods.}
Normals matter most when predictions must persist across time or support 3D reconstruction.
Video normal estimation enforces temporal consistency by warping predictions with optical flow or camera poses and penalizing flicker; NormalCrafter~\cite{wang2025normalcrafter} targets temporally consistent sequences for downstream video reconstruction and editing rather than per-frame accuracy alone.
Multi-view settings derive normals from reconstructed meshes or regularize them through cross-view photometric consistency, and neural surface reconstruction such as NeuS~\cite{wang2021neus} consumes normal-like differential geometry implicitly through signed-distance gradients, so normal estimation and reconstruction now supervise each other.

\paragraph{Material- and domain-specific normal estimation.}
A smaller branch revisits what generic models still get wrong.
Transparent, reflective, and laboratory objects violate Lambertian assumptions because their gradients come from refraction, reflection, or the background rather than the surface; TransNormal~\cite{li2026transnormal} targets transparent objects with a dedicated benchmark and model, making explicit that open-world normal prediction cannot be solved by indoor RGB-D supervision alone.



\begin{table}[!t]
\centering
\caption{Representative datasets and benchmarks for surface normal estimation. The ``Normal supervision/use'' column indicates whether normals are directly rendered, derived from depth or meshes, or indirectly related through intrinsic-image supervision. Counts are indicative of commonly used versions and may vary by split or preprocessing protocol.}
\label{tab:normal_dataset}
\fontsize{9}{11}\selectfont
\setlength{\tabcolsep}{4pt}
\begin{tabular}{L{0.15\textwidth}C{0.1\textwidth} C{0.15\textwidth}C{0.15\textwidth}L{0.35\textwidth}}
\toprule
\textbf{Dataset} & \textbf{Year} & \textbf{Domain} & \textbf{Source} & \textbf{Normal supervision / use} \\
\midrule
NYU Depth v2~\cite{silberman2012nyudv2} & 2012 & Indoor & RGB-D & Normals derived from Kinect depth; canonical indoor monocular normal benchmark. \\
SUN RGB-D~\cite{song2015sunrgbd} & 2015 & Indoor & RGB-D & Depth-derived normals across diverse indoor scenes and sensors. \\
ScanNet~\cite{dai2017scannet} & 2017 & Indoor & RGB-D video + mesh & Mesh- or depth-derived normals with camera trajectories; useful for video and reconstruction-linked supervision. \\
Matterport3D~\cite{chang2017matterport3d} & 2017 & Indoor panorama & RGB-D scans + mesh & High-resolution reconstructed geometry; supports room-scale and panoramic normal evaluation. \\
Taskonomy~\cite{zamir2018taskonomy} & 2018 & Indoor & 3D scans & Large multi-task supervision including surface normals, depth, reshading, and semantics. \\
Replica~\cite{straub2019replica} & 2019 & Indoor synthetic / scanned & Rendered geometry & Clean normals from high-quality reconstructed scenes; useful for controlled evaluation. \\
DIODE~\cite{vasiljevic2019diode} & 2019 & Indoor / outdoor & LiDAR depth & Depth-derived normals; stresses cross-domain generalization beyond NYU-style indoor scenes. \\
Hypersim~\cite{roberts2021hypersim} & 2021 & Indoor synthetic & PBR rendering & Ground-truth depth, normals, materials, and lighting for geometry + intrinsic learning. \\
Omnidata~\cite{eftekhar2021omnidata} & 2021 & Indoor multi-source & 3D scans & Large-scale mid-level vision training data with normals and many related tasks. \\
\bottomrule
\end{tabular}
\end{table}


\subsubsection{Datasets and Metrics}
\label{sec:normal_performance}

Table~\ref{tab:normal_dataset} shows that normal supervision is usually inherited from another representation rather than annotated directly.
RGB-D datasets (NYU Depth v2~\cite{silberman2012nyudv2}, SUN RGB-D~\cite{song2015sunrgbd}, ScanNet~\cite{dai2017scannet}) supply measured depth from which local planes or mesh normals are computed, which fixes the indoor evaluation regime but also means the labels are only as good as the depth sensor and are dominated by man-made indoor structure.
Scanned and synthetic collections address coverage and label quality from opposite sides: Matterport3D~\cite{chang2017matterport3d}, Taskonomy~\cite{zamir2018taskonomy}, and Omnidata~\cite{eftekhar2021omnidata} scale supervision across many environments and tasks, whereas Replica~\cite{straub2019replica} and Hypersim~\cite{roberts2021hypersim} give clean rendered normals with physically meaningful lighting at the price of a synthetic-to-real gap.
DIODE~\cite{vasiljevic2019diode} pushes toward indoor--outdoor generalization, yet outdoor normals remain far less standardized than outdoor depth because sparse LiDAR is a poor match for dense orientation supervision; intrinsic-image sets such as IIW~\cite{bell2014intrinsic} and MPI Sintel Intrinsics~\cite{chen2013simple} constrain shading and reflectance for inverse-rendering formulations without being normal benchmarks themselves.
The most visible gap is temporal: although ScanNet and other RGB-D videos contain sequences, there is no widely adopted video normal benchmark, so temporal stability is reported qualitatively or through method-specific warping losses, even though flicker often hurts reconstruction and AR consumers more than a small rise in mean angular error.




Evaluation is angular because normals are unit vectors.
With $\hat{\mathbf{n}}_i,\mathbf{n}_i\in\mathbb{S}^2$ the predicted and ground-truth normals at valid pixel $i$ and $N$ the number of valid pixels, the per-pixel error in degrees is
\begin{equation}
   e_i = \frac{180}{\pi}\arccos\!\left(\mathrm{clip}(\hat{\mathbf{n}}_i^\top \mathbf{n}_i,-1,1)\right),
\end{equation}
summarized by its mean (sensitive to boundary noise and invalid-depth artifacts), its median (robust to a few severe failures), occasionally its RMSE, and by threshold accuracies
\begin{equation}
    \mathrm{Acc}_{\tau} = \frac{1}{N}\sum_{i=1}^{N}\mathbf{1}(e_i < \tau),
    \qquad \tau\in\{11.25^\circ,22.5^\circ,30^\circ\},
\end{equation}
which expose changes in the fraction of practically usable pixels that a small shift in mean error can hide.
Training objectives usually minimize the cosine loss $1-\hat{\mathbf{n}}_i^\top\mathbf{n}_i$, monotonic in angle for unit vectors, but benchmarks report angles for interpretability.
Beyond per-pixel angle, joint depth--normal models are checked for agreement between predicted normals and normals differentiated from depth, reconstruction-oriented work asks whether normal discontinuities align with object boundaries and whether large planes stay coherent, and video methods report temporal consistency by warping predictions between adjacent frames with flow or pose,
\begin{equation}
    \mathrm{TempCons} = \frac{1}{M}\sum_{(i,t)}
    \left\|\hat{\mathbf{n}}_{i}^{\,t} -
    \mathcal{W}_{t+1\rightarrow t}(\hat{\mathbf{n}}^{\,t+1})_i\right\|_2,
\end{equation}
after rotating camera-frame normals into a common frame.
Protocol details move rankings: depth-derived normals are unstable at discontinuities, papers mask invalid pixels differently, and some evaluate in the camera frame while others rotate into a gravity-aligned frame, so the table below records which paper produced each number.

\begin{table}[!t]
    \centering
    \caption{Representative surface-normal estimators on NYU Depth v2 (lower mean / median angular error and higher $\mathrm{Acc}_{11.25^\circ}$ are better). Rows marked $\dagger$ are taken from the re-evaluation in StableNormal~\cite{ye2024stablenormal}, under whose protocol DSINE scores $18.6^\circ$ / $9.9^\circ$ / $56.1\%$; the other rows are as reported by the cited paper (Omnidata as re-evaluated by DSINE~\cite{bae2024dsine}). Numbers across protocols are indicative, not strictly comparable.}
    \label{tab:normal_performance}
    \fontsize{9}{11}\selectfont
    \setlength{\tabcolsep}{2.2mm}
    \begin{tabular}{L{0.26\textwidth} L{0.28\textwidth} c c c c}
    \toprule
    \textbf{Method} & \textbf{Family} & \textbf{Year} & \textbf{Mean}$\downarrow$ & \textbf{Median}$\downarrow$ & $\mathbf{Acc}_{11.25^\circ}\uparrow$ \\
    \midrule
    Eigen \& Fergus~\cite{eigen2015predicting} & Multi-scale CNN, multi-task & 2015 & 20.9 & 13.2 & 44.4 \\
    Omnidata~\cite{eftekhar2021omnidata} & Large-scale mid-level supervision & 2021 & 23.1 & 12.9 & 45.8 \\
    DSINE~\cite{bae2024dsine} & Ray-aware discriminative & 2024 & 16.4 & 8.4 & 59.6 \\
    Metric3Dv2 (ViT-g, zero-shot)~\cite{hu2024metric3dv2} & Geometric foundation model & 2024 & 13.3 & 7.0 & 66.4 \\
    Marigold (normal variant)$^\dagger$~\cite{ke2024repurposing} & Diffusion prior & 2024 & 20.9 & 11.1 & 50.5 \\
    GeoWizard$^\dagger$~\cite{fu2024geowizard} & Joint depth + normal diffusion & 2024 & 20.4 & 11.9 & 47.0 \\
    StableNormal$^\dagger$~\cite{ye2024stablenormal} & Low-variance diffusion & 2024 & 19.7 & 10.5 & 53.0 \\
    \bottomrule
    \end{tabular}
\end{table}

Table~\ref{tab:normal_performance} makes two things visible. Across the reported settings, camera-aware models achieve low angular errors: DSINE roughly halves the median error of the 2015 baseline, and Metric3Dv2 reaches $13.3^\circ$ mean error zero-shot, so projective inductive bias and large camera-diverse pretraining still pay off even with strong backbones. Diffusion-based estimators trail these numbers on NYU Depth v2 under the StableNormal protocol, but they trade accuracy on clean indoor scenes for coherence in textureless or ambiguous regions and for a usable uncertainty, which the single-frame benchmark does not reward; the post-GeoWizard line has accordingly shifted the target from single-frame accuracy toward stable, usable geometry, StableNormal controlling diffusion variance and NormalCrafter~\cite{wang2025normalcrafter} enforcing temporal coherence.

\subsubsection{Challenges}

An RGB image alone does not determine shape, lighting, and reflectance: a dark pixel may be a tilted surface, a dark material, a cast shadow, or a reflection of a dark region, and learned models resolve the ambiguity with priors rather than removing it, which is why estimators look confident on common indoor layouts and fail systematically on rare materials, unusual lighting, and out-of-distribution scenes.
The persistent failure modes follow from this.
\textit{Reflective and transparent surfaces} produce gradients caused by the environment rather than the surface.
\textit{Textureless and weakly lit regions} give little evidence, and networks hallucinate planar sheets.
\textit{Thin structures and depth discontinuities} mix receptive fields and inherit depth-derived label noise, blurring normals around chair legs, plants, railings, and silhouettes.
\textit{Wide-field and panoramic cameras} break models trained under narrow pinhole assumptions, which is what makes ray-aware designs such as DSINE~\cite{bae2024dsine} necessary.
\textit{Dynamic and deformable objects} flicker because most supervision comes from static scans or rigid synthetic scenes.

Dataset bias compounds the problem: canonical benchmarks are indoor and sensor-derived, so outdoor surfaces, hands, faces, deformable and transparent objects, and egocentric manipulation stay underrepresented, and no accepted benchmark jointly measures cross-domain generalization, temporal stability, and downstream reconstruction quality. TransNormal~\cite{li2026transnormal} is a corrective for one slice of the non-Lambertian regime (glass, liquid, mirror, metal, wet roads), which remains far from solved.
The directions now taking shape answer these gaps directly: geometric foundation models that merge depth, normals, segmentation, material, and calibration into one dense-prediction backbone, following Omnidata~\cite{eftekhar2021omnidata} through Metric3Dv2~\cite{hu2024metric3dv2} and GeoWizard~\cite{fu2024geowizard}; diffusion distillation and variance control that keep the coherence and uncertainty of generative models while recovering real-time, deterministic inference~\cite{ye2024stablenormal}; temporally consistent, reconstruction-coupled normals~\cite{wang2025normalcrafter}; and physically grounded models for non-Lambertian materials, where data scale is no substitute for modeling refraction and reflection.



\subsection{Height and Distance Estimation}
\label{sec:height_distance}
Height and distance estimation recovers absolute, real-world lengths: the dimensions of an object (the height of a person), the separation between scene entities (the gap between two vehicles), or the distance traveled by a moving object over a known time interval, from which average speed follows.
Whereas depth estimation (\S\ref{sec:depth_estimation}) recovers a dense per-pixel distance from the sensor to every visible surface, height and distance estimation targets a few metric quantities defined between scene entities rather than relative to the sensor, and is therefore tied to specific object instances and to a downstream measurement requirement.
Radar, ultrasonic, and radio ranging measure such distances directly and are mature; this subsection covers the vision-based problem.

\subsubsection{Task Definition}
The input is one or more images or a video together with a specification of what to measure: the height of an object, the distance between two objects, or the distance an object travels. The output is a length in physical units, or a traveled distance over a known time interval, whose ratio gives average speed (\textbf{speed estimation}).

Every vision-based route to a metric length passes through the camera model, shown in Figure~\ref{fig:p3_s6_2_2_camera_model}. Under the pinhole model a world point $P$ with homogeneous coordinates $\mathbf{P}_w$ projects to the pixel
\begin{equation}
    \mathbf{p} \sim K\,[\,R \mid \mathbf{t}\,]\,\mathbf{P}_w,
\end{equation}
where the extrinsics $(R,\mathbf{t})$ move the point into the camera frame and the intrinsic matrix $K$ (focal length, principal point, pixel aspect) maps normalized image-plane coordinates to pixels. Real lenses deviate from this model. A common forward form of Brown's radial distortion model~\cite{brown1971close} is
\begin{equation}
    \begin{pmatrix} x_{\text{dist}} \\ y_{\text{dist}} \end{pmatrix} =
    \begin{pmatrix} x(1 + k_1 r^2 + k_2 r^4 + k_3 r^6) \\ y(1 + k_1 r^2 + k_2 r^4 + k_3 r^6) \end{pmatrix},
\end{equation}
where $(x,y)$ are undistorted normalized coordinates and $r^2=x^2+y^2$. Tangential terms model sensor--lens misalignment, and camera calibration~\cite{zhang2000flexible} recovers $K$ together with the distortion coefficients. Without calibration, every pixel-to-length conversion inherits an unknown scale.

\begin{figure}[!ht]
    \centering
    \includegraphics[width=0.5\linewidth]{figures/part3/section6/camera-1.png}
    \caption{Pinhole camera model: world coordinates map to camera, image-plane, and pixel coordinates through the extrinsic, projection, and affine matrices; the last two compose into the intrinsic matrix $M_{\mathrm{intrinsic}}$ and all three into the full projection $M$.}
    \label{fig:p3_s6_2_2_camera_model}
\end{figure}

Beyond calibration, geometric pipelines rely on a small set of recurring intermediate quantities, illustrated in Figure~\ref{fig:p3_s6_2_2_subtasks}: vanishing points, where parallel world lines converge and from which camera orientation follows; the horizon line, which fixes the camera's tilt relative to the ground; and a ground-plane homography, which lets lengths on a dominant plane such as a road or floor be measured once a single reference length is known. Stereo pipelines add rectification, which aligns the two image planes so that correspondences lie on the same scanline, and disparity matching. Learning-based pipelines keep object detection (Section~\ref{sec:6-1-2}) as the front end and replace the explicit geometry with regression of depth, 3D center, and dimensions from region features, optionally conditioned on a dense depth map and refined over time when video is available.

\begin{figure}[!t]
    \centering
    \includegraphics[width=\linewidth]{figures/part3/section6/height_estimation_1.pdf}
    \caption{Geometric subtasks supporting height and distance estimation: (a) parallel street edges determine a vanishing point; (b) multiple views of a known planar pattern estimate camera intrinsics and lens distortion; (c) the ground-plane horizon passes through horizontal-direction vanishing points. Street geometry and the separate camera-pitch inset are illustrative.}
    \label{fig:p3_s6_2_2_subtasks}
\end{figure}

\subsubsection{Methodology}
Vision-based height and distance estimation is organized along two axes: the sensor configuration (stereo versus monocular), which decides whether metric depth is available by triangulation, and the inference paradigm (geometry-based versus learning-based), which decides whether scene geometry is solved from projective constraints or inferred from data. Figure~\ref{fig:p3_s6_2_2_method_overview} lays out the resulting families; it draws the uncalibrated case as a third branch for readability, but it is best read as the weakly constrained regime of monocular or stereo inference in which some geometric information is missing and must be supplied by priors.

\paragraph{Stereo Vision.}
A calibrated stereo rig recovers metric depth directly from disparity, $Z = fB/d$ with focal length $f$, baseline $B$, and disparity $d$.
Classical pipelines run calibration, rectification, matching, disparity estimation, and triangulation~\cite{zhang2000flexible,Scharstein2002Stereo,hartley2003multiple}; their accuracy rests on calibration quality and on matching reliability in textureless, reflective, or occluded regions, and object distance is read off by aggregating depth over the detected region. Multi-view stereo adds views for consistency without changing the triangulation principle~\cite{Seitz2006Middlebury,yao2018mvsnet}.
Learned stereo replaces hand-crafted matching costs with learned features and cost volumes~\cite{kendall2017end,chang2018pyramid,lipson2021raftstereo} and, in recent systems, folds distance estimation into stereo 3D object detection, predicting category, 2D extent, depth, and 3D box jointly~\cite{li2019stereorcnn,chen2020dsgn,guo2021ligastereo}. Learning improves robustness, but the metric scale still comes from the calibrated baseline and intrinsics.

\paragraph{Monocular Vision.}
A single image does not fix absolute scale, so monocular methods recover it from explicit geometric constraints, learned priors, or external reference information~\cite{hartley2003multiple,criminisi2000single}.
Classical monocular metrology uses projective geometry: vanishing points of parallel world lines give camera orientation~\cite{caprile1990using,hartley2003multiple}, and a homography between a dominant ground plane and the image enables metric measurement on that plane once one reference length is known~\cite{criminisi2000single,liebowitz1999metric}. The approach is interpretable and accurate in controlled scenes but needs clean geometric cues, reliable calibration, and visible parallel lines or planar supports.
Learning-based monocular methods detect objects in 2D and regress their depth, 3D center, and physical dimensions directly from image features~\cite{chen2016monocular,mousavian2017deep,manei2019monodis,brzozowski2019m3drpn,liu2020smoke,wang2021fcos3d,liu2022monocon,ma2021delving}, implicitly learning the inverse relation between image extent and distance together with perspective, context, and occlusion cues, without an explicit vanishing-point or homography step.

\paragraph{Foundation-Model-Assisted Metric Estimation.}
A newer route uses pretrained metric or metric-aligned depth models as an intermediate representation. DPT, ZoeDepth, Metric3D, and Depth Anything predict dense depth from a single image with strong cross-domain generalization~\cite{dpt,zoedepth,yin2024metric3d,DepthAnything}; pooling the predicted depth inside a detected region or instance mask then gives object distance, and combining masks, depth, and simple geometric constraints gives object height. This removes the need for explicit vanishing-point detection and manual scene calibration and offers a practical path to metric estimation on in-the-wild images.
In the fully uncalibrated case, exact metric recovery from one image remains ill-posed~\cite{hartley2003multiple}; practical systems inject scale through a known reference object, a known camera height, a ground-plane assumption, category-level size priors, or one of the pretrained metric-depth models above~\cite{criminisi2000single,mousavian2017deep,zoedepth,yin2024metric3d}.

\begin{figure}[!t]
    \centering
    \includegraphics[width=\linewidth]{figures/part3/section6/height_estimation_2.pdf}
    \caption{Sources of metric scale in vision-based measurement: (a) calibrated stereo uses focal length, disparity, and a known baseline; (b) monocular metrology transfers a reference height using ground-plane and vertical-direction geometry; (c) a reference-assisted example of weakly calibrated monocular inference uses a known planar rectangle to rectify a shared plane and set scale. Scenes are illustrative; the panels do not imply a universal accuracy ranking.}
    \label{fig:p3_s6_2_2_method_overview}
\end{figure}

\subsubsection{Datasets and Metrics}
Table~\ref{tab:height_distance_estimate_benchmark} summarizes the benchmarks, which fall into three groups by sub-task. Object-size estimation is evaluated on the indoor RGB-D sets CAMERA25 and REAL275~\cite{Wang0HVSG19}, where category-level dimensions are scored jointly with pose on controlled tabletop scenes. Sensor-to-object distance is rarely benchmarked on its own; it is evaluated through the 3D-detection annotations (box centers, sizes, velocities) of the driving datasets KITTI~\cite{Geiger2013KITTI}, nuScenes~\cite{caesar2020nuscenes}, Waymo Open~\cite{WaymoOpen-SunKDCPTGZCCVHN20}, and Argoverse~2~\cite{wilson2023argoverse}, which test metric localization at scale across multi-camera rigs and scenes, with BDD100K~\cite{yu2020bdd100k} as a monocular video source. Speed and displacement estimation is dominated by fixed-camera traffic video: AI City Challenge 2018~\cite{Naphade_2018_CVPR_Workshops}, BrnoCompSpeed~\cite{SochorJSMSHZ17}, UTFPR~\cite{LuvizonNM17}, and QMUL~\cite{moazzam2019image}, plus the UAV-mounted Beihang videos~\cite{KeKLW15}; these are the only benchmarks that score a physically meaningful motion quantity rather than static geometry.

\begin{table*}[!ht]
    \centering
    \caption{Height and distance estimation benchmarks, grouped by the sub-task they score (object size, sensor-to-object distance via 3D boxes, or vehicle speed). BDD100K is included as an auxiliary video source.}
    \label{tab:height_distance_estimate_benchmark}
    \fontsize{9}{11}\selectfont
    \setlength{\tabcolsep}{1.0mm}
    \begin{tabular}{L{0.19\textwidth} C{0.07\textwidth} C{0.15\textwidth} C{0.09\textwidth} C{0.09\textwidth} C{0.12\textwidth} C{0.21\textwidth}}
    \toprule
    \textbf{Dataset} & \textbf{Task} & \textbf{Modality} & \textbf{View} & \textbf{Domain} & \textbf{Scale} & \textbf{Resolution} \\
    \midrule
    CAMERA25 \cite{Wang0HVSG19} & Size  & RGB-D image & Monocular & Mixed-reality indoor & 25K eval. images & - \\
    REAL275 \cite{Wang0HVSG19} & Size  & RGB-D image & Monocular & Indoor & 2,750 test images & - \\
    KITTI 3D~\cite{Geiger2013KITTI} & Distance & RGB + LiDAR     & Monocular     & Vehicle     & 15K frames     & - \\
    nuScenes~\cite{caesar2020nuscenes}     & Size / distance     & RGB + LiDAR + radar     & Multi-view     & Vehicle     & 1,000 scenes & - \\
    Waymo Open Dataset~\cite{WaymoOpen-SunKDCPTGZCCVHN20}     & Size / distance   & RGB + LiDAR     & Multi-view     & Vehicle     & 1,150 segments & - \\
    Argoverse 2 Sensor~\cite{wilson2023argoverse}     & Size / distance  & RGB + LiDAR     & Multi-view     & Vehicle     & 1,000 scenes & - \\
    AI City 2018 Track 1 \cite{Naphade_2018_CVPR_Workshops} & Speed  & Video & Monocular & Vehicle & 27 one-min videos & 1920$\times$1080 @30\,Hz \\
    BrnoCompSpeed \cite{SochorJSMSHZ17} & Speed  & Video & Monocular & Vehicle & 18 videos $\times$ 1 hour & 1920$\times$1080 @50\,Hz\\
    UTFPR \cite{LuvizonNM17} & Speed  & Video & Monocular & Vehicle & 5 hours & 1920$\times$1080 @30\,Hz\\
    QMUL \cite{moazzam2019image} & Speed & Video & Monocular & Vehicle & 1 hour & 360$\times$288 @25\,Hz \\
    Beihang UAV Videos \cite{KeKLW15} & Speed & Image & Monocular & Vehicle & 70 images & 960$\times$540\\
    BDD100K~\cite{yu2020bdd100k}     & Auxiliary     & Video     & Monocular     & Vehicle     & 100K videos & 720p @30\,Hz\\
    \bottomrule
    \end{tabular}
\end{table*}

Each sub-task is scored with the same triplet: an absolute error in physical units, a relative error for cross-range comparison, and a threshold accuracy for downstream binary decisions. For object size, the metric is the mean absolute error in physical units (centimeters for human height, millimeters for category-level pose),
\begin{equation}
    \mathrm{MAE} = \frac{1}{N}\sum_{i=1}^{N} |\hat{s}_i - s_i|,
\end{equation}
with the mean absolute percentage error $\mathrm{MAPE} = \tfrac{1}{N}\sum_i |\hat{s}_i - s_i|/s_i \times 100\%$ alongside for cross-class comparison; when size is scored jointly with pose, the 3D IoU of predicted and ground-truth boxes is thresholded, the NOCS protocol~\cite{Wang0HVSG19} reporting mAP at 3D IoU thresholds of $25\%$ and $50\%$, and pose mAP under joint rotation/translation thresholds such as $(5^\circ,5\,\mathrm{cm})$, $(10^\circ,5\,\mathrm{cm})$, and $(10^\circ,10\,\mathrm{cm})$. Inter-object distance reuses MAE, RMSE, and MAPE on distances $\hat{d}_i$ and the depth-style ratio accuracy $\delta_\tau$ of \S\ref{sec:depth_estimation} with $\tau \in \{1.10,\,1.25\}$; because stereo error grows as $\Delta Z \propto Z^2/(fB)$, KITTI-style protocols~\cite{Geiger2013KITTI} stratify by range (near $<10$\,m, mid $10$--$50$\,m, far $>50$\,m) so that far-field degradation is not averaged away. Speed estimation reports MAE and RMSE in physical speed units, and MAPE as a percentage. AI City Challenge 2018~\cite{Naphade_2018_CVPR_Workshops} reports speed in mph and defines $S_1=\mathrm{DR}(1-\mathrm{NRMSE})$ from vehicle detection rate and speed RMSE min--max normalized across submissions, while BrnoCompSpeed~\cite{SochorJSMSHZ17} also reports median absolute error to reduce sensitivity to outliers. Two protocol details change rankings: per-vehicle averaging (BrnoCompSpeed) weighs trajectories differently from per-frame averaging over localized ground-truth vehicles (AI City), and monocular and uncalibrated methods are sensitive to the assumed focal length, principal point, and reference length, so their reported accuracy is only meaningful together with the calibration assumptions.

Performance is sub-task specific. Category-level size recovery on REAL275~\cite{Wang0HVSG19} remains far from millimeter accuracy in cluttered scenes. On BrnoCompSpeed~\cite{SochorJSMSHZ17}, calibration-aware pipelines reach low single-digit km/h errors on the original sessions but degrade when transferred to unseen camera placements, exposing the brittleness of fixed-camera homography assumptions, and on the AI City 2018 speed track~\cite{Naphade_2018_CVPR_Workshops} the leading entries pair strong detection-and-tracking front ends with multi-frame speed regression, so the S1 score reflects detection recall as much as raw speed accuracy. Across all three sub-tasks the bottleneck is calibration, of the camera and of the scene-plane reference, rather than the regression backbone.

\subsubsection{Challenges}
% 说明这个问题其实在当前的空间理解中，是非常重要的一个任务评估形式，它主要在与空间理解相关的下游任务中，作为一个重要的评估指标，来衡量模型对于空间理解的能力。比如在自动驾驶领域，车辆的速度估计是一个非常重要的任务，它直接关系到车辆的安全性和行驶效率；在室内场景中，物体尺寸的估计也是一个重要的任务，它可以帮助机器人更好地与环境进行交互。因此，准确的高度和距离估计对于提升空间理解模型的实用性和可靠性具有重要意义。
Calibration sensitivity is the dominant open problem: small errors in intrinsics or in the scene-plane reference propagate non-linearly into every recovered length, and methods that treat focal length and reference size as latent variables, or refine calibration online from multi-view consistency, generalize across camera setups where fixed-calibration pipelines fail. The second gap is coverage: benchmarks are confined to indoor tabletop size estimation and vehicle-traffic speed, leaving pedestrian speed in crowds, package dimensions in logistics, and indoor inter-object distance essentially unmeasured. Fusing visual estimates with LiDAR, radar, or inertial sensing would lift robustness in low light and under occlusion, but the vision-only setting treated here still lacks a benchmark that scores metric accuracy jointly with calibration uncertainty.



\subsection{Visual Odometry and Camera Pose Estimation}
\label{sec:visual_odometry}

Visual odometry (VO) and camera pose estimation recover the position and orientation of a camera from visual observations.
They provide the coordinate scaffold that lets a system stitch individual images into a persistent spatial representation: depth, surface normals, scene flow, maps, and 3D reconstructions all become more useful once they can be expressed in a shared camera or world frame.
VO usually means estimating \emph{incremental} motion between consecutive frames, often as a trajectory $\{\mathbf{T}_t\}_{t=1}^{T}$ in $SE(3)$; camera pose estimation is broader and may recover a query image's absolute 6-DoF pose in a known scene, the relative pose between two views, or camera motion and scene geometry jointly from an unordered image collection.
The field was built on geometric optimization and SLAM, and its center of gravity is now shifting toward learned correspondence, differentiable optimization, dense geometric priors, and feed-forward geometric foundation models that predict cameras and geometry together.

\subsubsection{Task Definition}

Let $\mathbf{T}_t\in SE(3)$ denote the $4\times4$ homogeneous camera-to-world transformation for frame $t$, with rotation $\mathbf{R}_t$ and translation $\mathbf{t}_t$.
Given images $\{I_t\}$ (monocular, stereo, RGB-D, multi-view, or video), optional intrinsics $\{\mathbf{K}_t\}$, optional inertial readings, and sometimes a prior map or database, the goal is to estimate relative camera motions $\Delta\mathbf{T}_{i\rightarrow j}=\mathbf{T}_i^{-1}\mathbf{T}_j$ expressed in frame $i$, an absolute pose $\mathbf{T}_q$ for a query image, or a full trajectory together with sparse or dense scene geometry, and occasionally the intrinsics themselves when calibration is unknown.
The variants differ in what is assumed known.
\textbf{Monocular VO} is scale ambiguous without external priors; \textbf{stereo and RGB-D VO} recover metric scale from a known baseline or measured depth; \textbf{visual-inertial odometry} uses IMU measurements to stabilize scale and fast rotation; \textbf{SLAM} adds mapping and loop closure to online odometry; \textbf{SfM} recovers poses and structure from unordered collections; \textbf{map-based relocalization} assumes an existing scene representation and places a query in it; and \textbf{map-free relocalization and feed-forward reconstruction} skip explicit map construction by predicting relative pose or globally aligned point maps directly from images.

\begin{figure}[!t]
    \centering
    \includegraphics[width=\linewidth]{figures/part3/section6/visual_odometry_and_camera_pose_estimation.pdf}
    \caption{Visual odometry and camera pose estimation. \textit{Top}: a video or unordered image set is mapped to relative motion $\Delta\mathbf{T}_{i\rightarrow j}$, an absolute pose $\mathbf{T}_q$ in a known scene, or camera poses $\{\mathbf{T}_t\}$ jointly with scene structure. \textit{Bottom}: method lanes (a) geometry-first, (b) learned and self-supervised VO, (c) hybrid learned--geometric, (d) relocalization, and (e) feed-forward foundation models. Mechanism arrows indicate operations; dated timelines summarize representative methods in publication order, with non-proportional spacing and no implied computational dependency. The dashed warp-feedback path denotes self-supervised training.}
    \label{fig:p3_s6_2_4_vo_overview}
\end{figure}

\subsubsection{Methodology}

Figure~\ref{fig:p3_s6_2_4_vo_overview} summarizes five method families, ordered from explicit geometry to learned geometric priors to feed-forward scene-level geometry.

\paragraph{Geometry-first SfM, VO, and SLAM.}
Classical pipelines detect repeatable features, match them across views, fit epipolar geometry with robust estimators, and refine poses and 3D points by bundle adjustment~\cite{triggs2000bundle,agarwal2010bundle}.
COLMAP~\cite{schonberger2016structure} made this recipe robust enough for internet photo collections, and GLOMAP~\cite{pan2024glomap} revisits the pose-graph and averaging components of global SfM for large-scale reconstruction.
For sequential input, geometric VO and SLAM track camera motion, with SLAM systems additionally maintaining maps and often closing loops: MonoSLAM~\cite{davison2007monoslam}, PTAM~\cite{KleinMurray2007PTAM}, and ORB-SLAM~\cite{mur2015orb} define the feature-based branch, LSD-SLAM~\cite{engel2014lsd} and DSO~\cite{engel2017direct} the direct branch.
Explicit geometric consistency is the family's strength; weak texture, dynamic objects, rolling shutter, illumination shifts, and small-baseline or planar degeneracies are its failure modes.

\paragraph{Learned VO and self-supervised view synthesis.}
End-to-end VO regresses camera motion from image pairs or short windows, DeepVO~\cite{Wang2017DeepVO} being the standard early example with convolutional features and recurrent temporal modeling.
Self-supervised VO instead learns depth and ego-motion together by minimizing photometric reconstruction error between adjacent frames, with SfM-Learner~\cite{zhou2017unsupervised} as the canonical starting point.
This made monocular video supervision scalable, but the geometry remains bound to photometric assumptions (static Lambertian scenes, stable intrinsics, sufficient parallax) and to monocular scale ambiguity unless stereo, IMU, or learned metric priors are added.

\paragraph{Hybrid learned geometry and differentiable optimization.}
The most successful learned systems keep geometric optimization and learn the parts that are hard to hand-design.
Learned features and matchers (SuperPoint~\cite{detone2018superpoint}, SuperGlue~\cite{sarlin2020superglue}, LoFTR~\cite{Sun2021LoFTR}) improve the correspondence front end, and VGG-SfM~\cite{Vggsfm} and ACE0~\cite{acezero} push learned features and scene-coordinate learning into SfM and relocalization.
DROID-SLAM~\cite{teed2021droid} is the representative differentiable-optimization system, coupling dense correspondence, depth, and camera motion through recurrent updates and a dense bundle-adjustment layer, and LEAP-VO~\cite{chen2024leap} stabilizes motion over longer windows through long-term point tracking.

\paragraph{Relocalization and map-based pose estimation.}
Given a query image and a known scene or database, relocalization estimates the query's 6-DoF pose.
PoseNet~\cite{kendall2015posenet} popularized absolute pose regression, which is fast but less accurate than geometry.
Scene-coordinate regression, from Shotton \etal~\cite{Scene_Coordinate} to DSAC~\cite{brachmann2017dsac} and ACE-style relocalizers~\cite{acezero}, predicts a 3D scene coordinate per pixel and solves pose with PnP and RANSAC; hierarchical localization~\cite{sarlin2019coarse} retrieves candidate database images, matches, and verifies geometrically, and remains the high-accuracy reference.
Relative-pose regressors trained on much larger collections, such as Reloc3r~\cite{reloc3r}, revisit the speed and generalization trade-off and show that regression recovers a role when paired with scale and data.

\paragraph{Feed-forward geometric foundation models.}
The main post-2024 shift is toward models that predict scene geometry and camera relations directly, often without SfM initialization.
DUSt3R~\cite{wang2024dust3r} regresses dense 3D point maps from image pairs, providing learned geometry from which correspondences and camera parameters are recovered using geometric post-processing; MASt3R~\cite{mast3r} grounds image matching in this 3D representation, and MASt3R-SfM~\cite{MASt3R_SfM} builds an unconstrained SfM pipeline on it.
MegaSaM~\cite{li2025megasam} and MonST3R~\cite{monst3r} extend the idea to casual dynamic video, where static-scene assumptions fail; CUT3R~\cite{wang2025Cut3R} adds persistent state for continuous 3D perception, Fast3R~\cite{yang2025fast3r} scales feed-forward reconstruction to very large image sets, and VGGT~\cite{wang2025vggt} predicts camera poses, depth, and point maps jointly with a visual-geometry-grounded Transformer.
Geometry does not disappear in this family; it moves from a hand-coded optimization loop into a learned prior that supplies initialization, dense correspondence, or complete camera-and-scene predictions.

\begin{table}[!t]
\centering
\caption{Representative datasets and benchmarks for visual odometry, camera pose estimation, relocalization, and feed-forward scene geometry. Counts and splits vary across protocols.}
\label{tab:vo_pose_datasets}
\fontsize{9}{11}\selectfont
\setlength{\tabcolsep}{4pt}
\begin{tabular}{L{0.15\textwidth}C{0.1\textwidth} C{0.15\textwidth}C{0.15\textwidth}L{0.35\textwidth}}
\toprule
\textbf{Dataset / benchmark} & \textbf{Year} & \textbf{Domain} & \textbf{Input} & \textbf{Primary use} \\
\midrule
KITTI Odometry~\cite{KITTI-GeigerLU12} & 2012 & Outdoor driving & Stereo + GPS/IMU & Classical and learned VO; segment drift evaluation. \\
TUM RGB-D~\cite{Sturm2012TUMRGBD} & 2012 & Indoor & RGB-D + mocap & Indoor RGB-D SLAM and trajectory accuracy. \\
7-Scenes~\cite{Scene_Coordinate} & 2013 & Indoor & RGB-D video & Camera relocalization and scene-coordinate regression. \\
Cambridge Landmarks~\cite{kendall2015posenet} & 2015 & Outdoor landmarks & RGB & Outdoor absolute pose regression and relocalization. \\
EuRoC MAV~\cite{Burri2016EuRoC} & 2016 & Aerial / industrial indoor & Stereo + IMU & Visual-inertial odometry under fast motion. \\
ScanNet~\cite{dai2017scannet} & 2017 & Indoor & RGB-D video + mesh & Learned pose, reconstruction, and dense geometry. \\
MegaDepth~\cite{li2018megadepth} & 2018 & Internet landmarks & SfM supervision & Relative pose, matching, depth, and geometric pretraining. \\
RealEstate10K~\cite{zhou2018stereo} & 2018 & Indoor / real-estate video & Video + recovered poses & Self-supervised video geometry and view synthesis. \\
Virtual KITTI 2~\cite{Cabon2020VirtualKITTI2} & 2020 & Synthetic driving & Rendered stereo / depth / pose & Controlled sim-to-real VO and depth evaluation. \\
DL3DV-10K~\cite{ling2024dl3dv} & 2024 & Large-scale scenes & Multi-view video / images & Generalizable 3D reconstruction and pose-aware training. \\
\bottomrule
\end{tabular}
\end{table}

\subsubsection{Datasets and Metrics}
\label{sec:vo_datasets}

Table~\ref{tab:vo_pose_datasets} shows three kinds of benchmark serving three different questions.
Trajectory benchmarks (KITTI Odometry~\cite{KITTI-GeigerLU12}, TUM RGB-D~\cite{Sturm2012TUMRGBD}, EuRoC MAV~\cite{Burri2016EuRoC}) measure drift under a fixed sensor setup and cover a narrow slice of deployment: urban driving, small rooms, micro-aerial-vehicle flights.
Relocalization benchmarks (7-Scenes~\cite{Scene_Coordinate}, Cambridge Landmarks~\cite{kendall2015posenet}) ask whether a query can be placed in a known scene and favor methods with scene-specific maps, retrieval, or scene-coordinate regression, so they are not comparable to online VO.
Large internet and video corpora (MegaDepth~\cite{li2018megadepth}, RealEstate10K~\cite{zhou2018stereo}, ScanNet~\cite{dai2017scannet}, DL3DV-10K~\cite{ling2024dl3dv}) serve mostly as training data for generalizable geometry models, which need diverse baselines, unknown intrinsics, unordered images, dynamic content, and dense reconstruction targets.

Evaluation is correspondingly layered.
With $\mathbf{T}_i$ and $\hat{\mathbf{T}}_i$ the ground-truth and estimated poses, the absolute trajectory error measures global discrepancy after a rigid or similarity alignment $\mathbf{S}$,
\begin{equation}
    \mathrm{ATE} =
    \sqrt{\frac{1}{N}\sum_{i=1}^{N}
    \left\|\mathrm{trans}\!\left(\mathbf{T}_i^{-1}\mathbf{S}\hat{\mathbf{T}}_i\right)\right\|^2},
\end{equation}
and the relative pose error measures local motion over an interval $\Delta$,
\begin{equation}
    \mathrm{RPE}(\Delta)=
    \frac{1}{N-\Delta}\sum_{i=1}^{N-\Delta}
    \left\|\mathrm{Log}\!\left[
    (\mathbf{T}_i^{-1}\mathbf{T}_{i+\Delta})^{-1}
    (\hat{\mathbf{T}}_i^{-1}\hat{\mathbf{T}}_{i+\Delta})
    \right]\right\|;
\end{equation}
KITTI reports the same idea as translational and rotational drift averaged over fixed-length path segments.
Whether scale is aligned before computing ATE, or scale drift is scored separately, decides whether monocular, stereo, RGB-D, and visual-inertial methods can be compared at all.
Pairwise matching and feed-forward pose methods report rotation error and translation-direction error in degrees, often aggregated as the area under the accuracy curve (AUC) below angular thresholds; relocalization reports the fraction of queries within thresholds such as $(5\,\mathrm{cm},5^\circ)$ indoors or meter-level outdoors, or median translation and rotation error; and feed-forward geometry models add depth error, point-map error, Chamfer distance, registration rate, or view-synthesis quality, since their cameras are judged by whether they support a usable 3D representation.

\begin{table*}[!t]
    \centering
    \caption{Representative methods under two shared protocols. \textit{Top}: monocular SLAM / VO, absolute trajectory error [m] averaged over the 11 EuRoC MAV sequences and the 9 TUM-RGBD freiburg1 sequences, as reported in DROID-SLAM~\cite{teed2021droid} (``fails $k$'' = lost tracking on $k$ sequences, so no average is defined). \textit{Bottom}: multi-view camera pose, AUC@30 on RealEstate10K (unseen) and CO3Dv2, as reported in VGGT~\cite{wang2025vggt}; the learned methods were trained on CO3Dv2 but not RealEstate10K. Numbers are from the cited evaluating paper and are comparable only within a block.}
    \label{tab:vo_pose_performance}
    \fontsize{9}{11}\selectfont
    \setlength{\tabcolsep}{1.5mm}
    \begin{tabular}{L{0.27\textwidth} L{0.23\textwidth} c C{0.14\textwidth} C{0.14\textwidth}}
    \toprule
    \textbf{Method} & \textbf{Family} & \textbf{Year} & \textbf{EuRoC ATE}$\downarrow$ & \textbf{TUM fr1 ATE}$\downarrow$ \\
    \midrule
    ORB-SLAM~\cite{mur2015orb} & Feature-based SLAM & 2015 & fails 2 / 11 & fails 6 / 9 \\
    DSO~\cite{engel2017direct} & Direct sparse odometry & 2017 & 0.601 & -- \\
    ORB-SLAM3~\cite{Campos2021ORBSLAM3} & Feature-based multi-map SLAM & 2021 & fails 1 / 11 & fails 5 / 9 \\
    DROID-SLAM~\cite{teed2021droid} & Differentiable dense BA & 2021 & \textbf{0.022} & \textbf{0.038} \\
    \midrule
    \hiderowcolors \textbf{Method} & \textbf{Family} & \textbf{Year} & \textbf{Re10K AUC@30}$\uparrow$ & \textbf{CO3Dv2 AUC@30}$\uparrow$ \\
    \midrule
    COLMAP + SuperPoint/SuperGlue~\cite{schonberger2016structure,detone2018superpoint,sarlin2020superglue} & Incremental SfM, learned matching & 2016 / 2020 & 45.2 & 25.3 \\
    DUSt3R~\cite{wang2024dust3r} & Pairwise point-map regression & 2024 & 67.7 & 76.7 \\
    MASt3R~\cite{mast3r} & 3D-grounded matching & 2024 & 76.4 & 81.8 \\
    VGGSfM v2~\cite{Vggsfm} & Learned differentiable SfM & 2024 & 78.9 & 83.4 \\
    CUT3R~\cite{wang2025Cut3R} & Continuous 3D perception with state & 2025 & 75.3 & 82.8 \\
    Fast3R~\cite{yang2025fast3r} & Large image-set feed-forward & 2025 & 72.7 & 82.5 \\
    VGGT~\cite{wang2025vggt} & Visual geometry grounded Transformer & 2025 & \textbf{85.3} (93.5 w/ BA) & \textbf{88.2} (91.8 w/ BA) \\
    \bottomrule
    \end{tabular}
\end{table*}

Table~\ref{tab:vo_pose_performance} puts the two regimes side by side. On trajectory benchmarks the lesson is robustness as much as accuracy: the classical systems that are highly accurate when they track lose tracking on a fraction of EuRoC flights and on most TUM freiburg1 sequences, whose rolling shutter, motion blur, and fast rotation defeat feature matching, whereas DROID-SLAM's learned dense correspondence inside a bundle-adjustment loop tracks every sequence at centimeter-level error. On multi-view pose the feed-forward models have overtaken the COLMAP pipeline by a wide margin and keep improving on unseen data (RealEstate10K), with VGGT ahead of DUSt3R-style pairwise regression at a fraction of the runtime, and a short bundle-adjustment refinement still adding several points, which is the hybrid pattern the challenges below return to. The two blocks are not comparable to each other, and that is the point: a method optimized for low drift under a calibrated rig and one optimized for zero-shot sparse-view geometry are judged by different questions.

\subsubsection{Challenges}

Geometric ambiguity sets the floor.
Monocular cameras cannot recover metric scale without priors; pure rotation, planar scenes, and very small baselines make epipolar geometry ill-conditioned; rolling shutter and unknown intrinsics violate the assumptions built into many solvers.
Learned models can hide these failures behind priors but cannot restore missing information, which is why visual-inertial fusion, calibration, and depth priors remain necessary even with large geometry models.
Deployment adds its own: dynamic agents corrupt the static-world assumption behind classical SLAM and bundle adjustment; textureless rooms, repetitive corridors, glass, water, and specular surfaces yield unreliable correspondences; appearance shift across time of day, weather, season, and camera hardware breaks place recognition; and long-horizon operation still needs loop closure, map maintenance, and uncertainty handling, none of which a feed-forward prediction supplies on its own.

The strongest trend is the convergence of SLAM, reconstruction, and foundation geometry. DROID-SLAM~\cite{teed2021droid} showed that optimization can be learned without abandoning geometric structure, and the feed-forward models above predict dense geometry, cameras, and correspondences together; the near-term frontier is hybrid, with feed-forward geometry supplying initialization, dense priors, and generalization and optimization enforcing metric consistency, loop closure, and uncertainty-aware correction.
Two further shifts matter for spatial intelligence. Pose estimation must become \textbf{dynamic-scene aware}, estimating camera and object motion jointly as MegaSaM~\cite{li2025megasam} and MonST3R~\cite{monst3r} begin to do, rather than discarding moving content as outliers; and camera pose should be judged by \textbf{downstream usefulness}, whether it supports stable depth and normal prediction, reconstruction, embodied navigation, AR anchoring, or world-model memory. As pose becomes one output of larger geometric perception systems, the boundary between VO, SLAM, scene reconstruction, and spatial memory will keep blurring.




\subsection{Pose Estimation}
\label{sec:object_pose}
% Author intent (Hao Fei): image to human-pose or 3D pose.
% Scope decision: §6.2.5 covers 6-DoF rigid-object pose;
% sparse landmark detection on humans, hands, and animals is owned by
% §6.1.1, and camera-frame pose is owned by §6.2.4.

6-DoF object pose estimation recovers a rigid transformation
\begin{equation}
\mathbf{T}_o =
\begin{bmatrix}
\mathbf{R} & \mathbf{t}\\
\mathbf{0}^{\top} & 1
\end{bmatrix}
\in SE(3),
\end{equation}
which maps an object's canonical or model frame into the current camera frame, giving a rotation $\mathbf{R}\in SO(3)$ and a translation $\mathbf{t}\in\mathbb{R}^3$ per detected instance.
It is the camera pose problem of \S\ref{sec:visual_odometry} with the reference frame reversed: the camera frame serves as the reference and each rigid object is located within it, using much of the same machinery (2D--3D correspondence, PnP solvers, RANSAC, and render-and-compare refinement) but with object-centric CAD models and pose annotations as supervision.
It is also distinct from the keypoint and landmark detection of \S\ref{sec:keypoint_landmark}: the target is the full 6-DoF pose of a rigid body, not semantic points on an articulated or deformable one.
Object pose is the geometric primitive consumed by robotic manipulation and grasping (revisited in Chapter~\ref{sec:embodied}), AR/VR object overlay, assembly guidance, bin picking, and the metric localization of other agents in driving.

Three axes organize the literature. The \emph{assumption axis} runs from instance-level methods that assume a high-fidelity CAD mesh per object, through category-level methods that assume only a canonical frame per category, to model-free methods that see a previously unseen object through a handful of reference images. The \emph{representation axis} separates direct $SE(3)$ regression from keypoint- or dense-correspondence pipelines that hand pose recovery to PnP, and from render-and-compare refinement of an initial pose. The \emph{foundation-model axis} tracks the shift from per-object training to cross-object generalization built on segmentation, vision, and depth foundation models.

\subsubsection{Task Definition}
The input is an RGB image, an RGB-D image, or a short multi-view burst of a scene containing one or more rigid objects, together with object-specific priors that depend on the setting: a CAD mesh or a small reference image set per object (instance-level and model-free), or a category label with a canonical frame (category-level), optionally with camera intrinsics and an upstream 2D detection or mask.
The output is a pose $\mathbf{T}_o\in SE(3)$ per detected object in the camera frame, augmented in the category-level setting with a metric size vector $\mathbf{s}\in\mathbb{R}_{+}^3$ (9-DoF pose and size).
For symmetric objects the prediction is properly an equivalence class of poses under the object's symmetry group, which the ADD-S and BOP metrics below account for.

\subsubsection{Methodology}

Figure~\ref{fig:p3_s6_2_5_pose_overview} schematizes the task and the three method-family trajectories, which recent methods often straddle.

\begin{figure}[!t]
    \centering
    \includegraphics[width=0.99\linewidth]{figures/part3/section6/pose_estimation.pdf}
    \caption{6-DoF object pose estimation. \textit{Top}: recover $\mathbf{T}_o\in SE(3)$ for each rigid object in an RGB(-D) image, given a CAD mesh, a canonical frame, or a few reference views. \textit{Bottom}: method timelines for (i) instance-level CAD-aware, (ii) category-level (NOCS) canonical-frame, and (iii) novel-object, foundation-era pose estimation.}
    \label{fig:p3_s6_2_5_pose_overview}
\end{figure}

\paragraph{Instance-Level 6-DoF: CAD-Aware Pipelines.}
With a mesh of each target object available, pose estimation becomes alignment of the mesh with image evidence, and synthetic renderings of the CAD models are the main training signal.
PoseCNN~\cite{xiang2017posecnn} set the modern baseline by regressing an object's image-space center and a quaternion rotation with a CNN, and introduced the YCB-Video benchmark.
A second family decouples the network from the geometric solver by predicting projections of sparse 3D keypoints and solving PnP with RANSAC; PVNet~\cite{peng2019pvnet} votes per pixel for each keypoint, which makes it robust to occlusion, while SSD-6D~\cite{kehl2017ssd6d} folds a discretized viewpoint classification into a detector and skips PnP entirely.
DenseFusion~\cite{Wang2019DenseFusion} fuses RGB-D features per pixel and directly regresses poses with iterative refinement. A third family uses dense surface correspondences instead of sparse keypoints: GDR-Net~\cite{wang2021gdrnet} couples dense coordinate regression with a geometry-guided PnP layer, ZebraPose~\cite{su2022zebrapose} encodes the surface with a coarse-to-fine binary code, and SurfEmb~\cite{haugaard2022surfemb} learns continuous correspondence distributions that stay robust under symmetry and occlusion; EPro-PnP~\cite{chen2022epro} replaces the non-differentiable PnP solver with a probabilistic surrogate so the whole pipeline trains end to end.
A fourth family refines by render-and-compare: DeepIM~\cite{li2018deepim} iteratively minimizes the residual between a rendered CAD view and the image, and CosyPose~\cite{Labbe2020CosyPose} extends this to consistent multi-view, multi-object refinement and led the BOP 2020 leaderboard.
The same machinery reaches far from the tabletop, for example spacecraft pose under extreme scale variation~\cite{hu2021widedepthrange}.

\paragraph{Category-Level 6-DoF + Size: Canonical Frame Recovery.}
Without a per-instance mesh but with a category that admits a canonical frame, pose is recovered up to that frame together with three-dimensional object size, a 9-DoF prediction.
NOCS~\cite{Wang0HVSG19} introduced the formulation, mapping every instance of a category into a Normalized Object Coordinate Space and regressing dense correspondences there; its REAL275 benchmark is the standard category-level evaluation.
Later work re-parameterizes the problem: FS-Net~\cite{chen2021fsnet} decouples rotation, translation, and size and speeds up inference through a shape-based pathway, GPV-Pose~\cite{di2022gpvpose} grounds rotation in geometry-guided point-wise voting, SAR-Net~\cite{lin2022sarnet} recovers a deformed shape prior alongside pose to handle intra-class variation, and IST-Net~\cite{liu2023istnet} drops the shape prior and learns the canonical transformation implicitly.
SecondPose~\cite{chen2024secondpose} and AG-Pose~\cite{lin2024agpose} import DINOv2 semantics and instance-adaptive keypoints for further gains on REAL275, and GenPose~\cite{zhang2023genpose} casts category-level pose as score-based diffusion sampling, the precursor of the 9-DoF diffusion model below.
The family's standing assumption, that a meaningful canonical frame exists per category, holds for mugs and cameras but fails for heterogeneous semantic groupings, which motivates the model-free line.

\paragraph{Novel-Object Pose: Model-Based and Model-Free Pipelines.}
The largest recent shift is away from per-object and per-category training toward objects first seen at test time, given only a few reference images or a CAD model.
From a short reference video, OnePose~\cite{sun2022onepose} triangulates a feature point cloud and matches test queries against it, OnePose++~\cite{he2022oneposepp} removes the dependence on repeatable keypoints with a dense-matching front end, and Gen6D~\cite{liu2022gen6d} unifies detection, viewpoint selection, and refinement without per-object fine-tuning.
From a CAD model alone, MegaPose~\cite{labbe2022megapose} trains render-and-compare on millions of synthetic objects and transfers zero-shot, FoundationPose~\cite{wen2024foundationpose} consolidates model-based and model-free operation under large-scale synthetic pretraining and leads the BOP novel-object track, GigaPose~\cite{nguyen2024gigapose} cuts the cost of render-and-compare by sampling templates in a two-DoF rotation space and recovering the remaining parameters from a single patch correspondence, and SAM-6D~\cite{lin2024sam6d} uses the Segment Anything Model~\cite{kirillov2023sam} as a zero-shot detection front end before learned pose matching.
A further sub-line drops pose-specific training altogether and matches frozen foundation features: FoundPose~\cite{ornek2024foundpose} establishes 2D--3D correspondences between DINOv2 patch descriptors of the query and pre-rendered templates, and FreeZe~\cite{caraffa2024freeze} fuses frozen 3D geometric descriptors with 2D foundation features and recovers pose by RANSAC registration; a FreeZe variant took first place for unseen-object localization and detection in the BOP Challenge 2024~\cite{nguyen2024bopchallenge}, overtaking the render-and-compare line.
Diff9D~\cite{liu2024diff9d} reformulates 9-DoF category-level pose as diffusion sampling, the object-pose instance of the diffusion-as-perception pattern seen for depth and normals, which yields calibrated multi-hypothesis predictions (natural for symmetric objects) at the cost of latency.
The foundation-model line has thus largely dissolved the instance-versus-category-versus-model-free distinction; the active question is how cheap the per-test-object supervision can become.

\paragraph{Sensing Modalities.}
Across these axes, the input modality fixes which supervisory signals exist and which symmetries can be resolved. RGB-only pipelines~\cite{xiang2017posecnn,peng2019pvnet,wang2021gdrnet,chen2022epro,sun2022onepose} suit mobile and AR deployment but face scale--depth ambiguity on texture-poor objects; RGB-D fusion (DenseFusion~\cite{Wang2019DenseFusion}, CosyPose~\cite{Labbe2020CosyPose} with depth refinement, FoundationPose~\cite{wen2024foundationpose}) is substantially more accurate on BOP at the cost of a depth sensor; depth-only methods are rare but remain competitive on highly symmetric objects where shading carries little rotational evidence. Monocular depth foundation models (\S\ref{sec:depth_estimation}) are increasingly used as plug-in depth providers at inference, which narrows the RGB / RGB-D gap without an active sensor.

\subsubsection{Datasets and Metrics}
\label{sec:pose_datasets}

Table~\ref{tab:pose_datasets} lists the recurring benchmarks. Instance-level evaluation grew from LINEMOD~\cite{hinterstoisser2012linemod}, the largely texture-poor early benchmark that introduced the ADD metric, through its occlusion-heavy re-annotation LINEMOD-Occluded~\cite{brachmann2014lmo}, which modern methods report as the more informative of the two, to the video setting of YCB-Video~\cite{xiang2017posecnn} and the texture-less, symmetric industrial parts of T-LESS~\cite{Hodan2017TLESS}, where symmetry handling decides the ranking; HomebrewedDB~\cite{kaskman2019hbdb} serves as an out-of-distribution test for methods trained on the BOP core data, and SPEED~\cite{Sharma2019SpacecraftPE} stretches the task to spacecraft under scale and illumination ranges no tabletop set covers. NOCS-REAL275 with its synthetic companion CAMERA~\cite{Wang0HVSG19} is the category-level standard. The BOP Challenge~\cite{hodan2018bop,hodan2020bop,sundermeyer2023bop} consolidates the instance-level datasets under one protocol and has been the field's headline leaderboard since 2017; its 2023 edition distinguishes seen-object and unseen-object settings, with 3D mesh models provided in both. Modern papers report at minimum on LINEMOD-Occluded, YCB-Video, and a BOP-core subset, novel-object methods add the BOP unseen-object track, and category-level methods center on REAL275.

\begin{table*}[!t]
    \centering
    \caption{Representative datasets and benchmarks for 6-DoF object pose estimation. ``--'' marks a value not fixed by the source or varying across protocol versions.}
    \label{tab:pose_datasets}
    \fontsize{9}{11}\selectfont
    \setlength{\tabcolsep}{1.2mm}
    \begin{tabular}{L{0.20\textwidth} c L{0.23\textwidth} C{0.09\textwidth} C{0.13\textwidth} L{0.19\textwidth}}
    \toprule
    \textbf{Dataset} & \textbf{Year} & \textbf{Setting} & \textbf{Objects} & \textbf{Frames} & \textbf{Role} \\
    \midrule
    LINEMOD~\cite{hinterstoisser2012linemod} & 2012 & Instance, texture-poor objects, RGB-D & 15 & $\sim$1.2K / object & Entry benchmark; introduced ADD \\
    LINEMOD-Occluded~\cite{brachmann2014lmo} & 2014 & Instance, heavy mutual occlusion & 8 & re-annotated LM scene & Occlusion stress test \\
    YCB-Video~\cite{xiang2017posecnn} & 2017 & Instance, RGB-D video, moving camera & 21 & $\sim$133K & Sequential pose and tracking \\
    T-LESS~\cite{Hodan2017TLESS} & 2017 & Instance, texture-less, symmetric industrial parts & 30 & -- & Symmetry / texture-poor stress test \\
    HomebrewedDB~\cite{kaskman2019hbdb} & 2019 & Instance, household scenes & -- & -- & Out-of-distribution generalization \\
    NOCS REAL + CAMERA~\cite{Wang0HVSG19} & 2019 & Category-level 9-DoF, RGB-D & 6 categories & $\sim$8K real (+ synthetic) & Category-level standard \\
    SPEED~\cite{Sharma2019SpacecraftPE} & 2019 & Spacecraft, extreme scale and lighting & 1 & -- & Beyond-tabletop deployment \\
    BOP Challenge~\cite{hodan2018bop,hodan2020bop,sundermeyer2023bop} & 2018-- & Unified protocol over LM-O, T-LESS, YCB-V, ITODD, IC-BIN, TUD-L, HB, HOPE & -- & -- & Headline leaderboard; novel-object track since 2023 \\
    \bottomrule
    \end{tabular}
\end{table*}

Each metric answers to a particular failure mode of pose error, and the symmetry-aware variants exist because many objects admit several valid poses.
The Average Distance of Model Points (ADD) is the standard for asymmetric objects: with predicted pose $(\hat{\mathbf{R}}, \hat{\mathbf{t}})$, ground truth $(\mathbf{R}, \mathbf{t})$, and mesh vertices $\mathcal{M}$,
\begin{equation}
    \mathrm{ADD} = \frac{1}{|\mathcal{M}|} \sum_{\mathbf{x} \in \mathcal{M}}
    \left\| (\mathbf{R}\mathbf{x} + \mathbf{t}) - (\hat{\mathbf{R}}\mathbf{x} + \hat{\mathbf{t}}) \right\|,
\end{equation}
and a pose counts as correct when $\mathrm{ADD} < 0.1\,d$ for object diameter $d$, the ADD-0.1d threshold introduced with LINEMOD~\cite{hinterstoisser2012linemod}.
For symmetric objects ADD-S replaces direct vertex correspondence with nearest-neighbor matching,
\begin{equation}
    \mathrm{ADD\text{-}S} = \frac{1}{|\mathcal{M}|} \sum_{\mathbf{x}_1 \in \mathcal{M}}
    \min_{\mathbf{x}_2 \in \mathcal{M}}
    \left\| (\mathbf{R}\mathbf{x}_1 + \mathbf{t}) - (\hat{\mathbf{R}}\mathbf{x}_2 + \hat{\mathbf{t}}) \right\|,
\end{equation}
under the same threshold; the two are reported together so that symmetric and asymmetric objects are diagnosed separately.
For AR overlay the relevant error is in image space, so the 2D projection error averages the pixel distance between mesh vertices projected under the two poses, with 5\,px as the usual threshold.
Category-level evaluation on REAL275 counts a prediction correct when
\begin{equation}
    \angle(\mathbf{R}^{-1} \hat{\mathbf{R}}) < 5^\circ
    \quad \text{and} \quad
    \| \mathbf{t} - \hat{\mathbf{t}} \| < 5\,\mathrm{cm},
\end{equation}
and adds 3D box IoU at $50\%$ and $75\%$ to score pose and size jointly.
The BOP Challenge~\cite{hodan2018bop} averages recall over three symmetry-aware errors, each thresholded at several operating points: the Visible Surface Discrepancy (VSD), which scores only the visible object surface and is therefore robust to occlusion, the Maximum Symmetry-aware Surface Distance (MSSD), the maximum surface-point deviation after selecting the best valid symmetry alignment, and its image-space analogue MSPD,
\begin{equation}
    \mathrm{AR} = \tfrac{1}{3} \bigl( \mathrm{AR}_{\mathrm{VSD}} + \mathrm{AR}_{\mathrm{MSSD}} + \mathrm{AR}_{\mathrm{MSPD}} \bigr),
\end{equation}
averaged over the BOP core datasets.
ADD and ADD-S remain the default for LINEMOD-style instance-level work, BOP-AR for cross-dataset comparison, and projection error for AR; the metrics persist side by side because no single norm captures both strict pose error and the equivalence-class structure that VSD and MSSD make explicit.

\begin{table}[!t]
    \centering
    \caption{Representative 6-DoF object pose results on LINEMOD-Occluded (higher is better). ADD(-S)-0.1d anchors the RGB methods; BOP-AR is reported for the refinement and foundation-era systems. The two columns are different metrics and are not directly comparable.}
    \label{tab:pose_benchmark}
    \fontsize{9}{11}\selectfont
    \setlength{\tabcolsep}{4pt}
    \begin{tabular}{lllllcc}
        \toprule
        \textbf{Method} & \textbf{Family} & \textbf{Modality} & \textbf{Train assumption} & \textbf{Year} & \textbf{ADD(-S)-0.1d} $\uparrow$ & \textbf{BOP-AR} $\uparrow$ \\
        \midrule
        PoseCNN~\cite{xiang2017posecnn}             & Instance / direct         & RGB   & per-object   & 2017 & 24.9 & --   \\
        PVNet~\cite{peng2019pvnet}                  & Instance / keypoint+PnP   & RGB   & per-object   & 2019 & 40.8 & --   \\
        GDR-Net~\cite{wang2021gdrnet}               & Instance / coord. reg.    & RGB   & per-object   & 2021 & 62.2 & --   \\
        % 需核查: 63.3 is CosyPose's RGB single-view LM-O result (BOP 2020); the RGB-D refined variant scores higher.
        CosyPose~\cite{Labbe2020CosyPose}           & Instance / refinement     & RGB   & per-object   & 2020 & --   & 63.3 \\
        SurfEmb~\cite{haugaard2022surfemb}          & Instance / dense corr.    & RGB-D & per-object   & 2022 & --   & 75.8 \\
        % 需核查: MegaPose LM-O BOP-AR value and modality to be confirmed against the BOP leaderboard.
        MegaPose~\cite{labbe2022megapose}           & Novel-object / render\&comp. & RGB-D & cross-object & 2022 & --   & 58.3 \\
        FoundationPose~\cite{wen2024foundationpose} & Novel-object / found.       & RGB-D & cross-object & 2024 & --   & 78.8 \\
        \bottomrule
    \end{tabular}
\end{table}

Table~\ref{tab:pose_benchmark} reports LINEMOD-Occluded, the most widely used instance-level benchmark under the BOP protocol.
Among RGB methods scored by ADD(-S)-0.1d, explicit geometry steadily closes the gap over direct regression: PoseCNN's 24.9 rises to 40.8 with per-pixel keypoint voting (PVNet) and to 62.2 with geometry-guided dense coordinate regression (GDR-Net).
On the symmetry-aware BOP-AR, the render-and-compare CosyPose (63.3) and the dense-correspondence SurfEmb (75.8) are now matched or exceeded by novel-object systems without per-object training, FoundationPose reaching 78.8 against MegaPose's earlier 58.3; this is the 2023--2024 shift that made novel-object pose a deployment target rather than a research curiosity.
The gap between the two families is widest on the symmetric objects of T-LESS, where losses that model symmetry (BOP-AR-aware objectives, ADD-S optimization) generalize better than those assuming a single asymmetric pose.

\subsubsection{Challenges}

\textbf{The onboarding bottleneck has shifted from pose to detection.}
Moving to novel objects relocated the limiting stage rather than removing it.
On the BOP Challenge 2024~\cite{nguyen2024bopchallenge}, 6-DoF \emph{localization} of unseen objects has nearly caught up with the seen-object regime (the strongest training-free entry, a FreeZe variant~\cite{caraffa2024freeze}, trails the best seen-object method by a few points of average recall), whereas \emph{detection} of unseen objects still lags its seen-object counterpart by roughly a third and is now the largest error source in the novel-object stack.
The open problem is zero-shot detection and segmentation of a never-seen instance in clutter, together with the cost of \emph{onboarding} it: BOP 2024 added a track in which an object is registered from a short reference video, Any6D~\cite{lee2025any6d} registers from a single RGB-D anchor image, and how few views, or whether a single image or a text prompt, suffice is the defining frontier question. The same challenge pivoted toward high-resolution AR/VR headset recordings, a signal that the benchmark-realism gap, including articulated and deformable objects, open-vocabulary targets named rather than modeled (POPE~\cite{fan2023pope}), and egocentric capture, is now a first-order concern.

\textbf{Absolute metric scale remains ambiguous for model-free targets.}
The true size of an unseen object is under-determined from RGB alone, since scale and depth trade off along the viewing ray.
Category-level 9-DoF methods absorb this by regressing size in a canonical frame~\cite{Wang0HVSG19}, and model-free systems can use a depth sensor or another metric reference, while CAD-based methods can use the model's known extent; without such information, scale is the residual unknown separating a plausible pose from a usable one, which Any6D~\cite{lee2025any6d} makes explicit by coupling pose recovery with scale alignment, and monocular depth priors inherit the affine ambiguity of \S\ref{sec:depth_estimation}.

\textbf{Accuracy trades sharply against latency.}
The methods that top the novel-object leaderboard are the least suited to closed-loop use: on BOP 2024~\cite{nguyen2024bopchallenge} the most accurate unseen-object method runs in tens of seconds per image, close to an order of magnitude slower than both the fastest seen-object pipelines and the practical sub-second entries, which buy speed with accuracy.
Heavy segmentation, render-and-compare, and diffusion priors (Diff9D~\cite{liu2024diff9d}) are the cause, which is why one-correspondence template matching (GigaPose~\cite{nguyen2024gigapose}), distillation, and selective invocation of expensive components have become load-bearing for robotic and AR loops at frame rate.

\textbf{Symmetry and geometric ambiguity are irreducible.}
Symmetric and texture-less parts of the kind that fill T-LESS~\cite{Hodan2017TLESS} admit a continuous or discrete family of photometrically indistinguishable poses, and any loss assuming a unique target diverges on them; transparent and specular surfaces violate the brightness-constancy and depth-sensing assumptions behind correspondence and RGB-D fusion; rotation parameterization adds degeneracies that continuous 6D or axis-angle encodings reduce but do not remove; and heavy occlusion, the regime of LINEMOD-Occluded~\cite{brachmann2014lmo}, starves the estimator of the evidence needed to resolve any of them.
Generative multi-hypothesis pose (Diff9D~\cite{liu2024diff9d}) predicts a distribution instead of a point, but turning that distribution into a \emph{calibrated} confidence a planner can trust in safety-critical manipulation and driving is still largely unaddressed.

Across these problems, pose pipelines increasingly consume foundation-model priors as plug-in modules (promptable segmentation for zero-shot detection, monocular depth for inference-time geometry) rather than training pose-specific backbones, and the boundary between pose estimation, 6-DoF tracking, and neural reconstruction is dissolving, as FoundationPose~\cite{wen2024foundationpose} folds estimation and tracking into one loop and the reconstruction discussion of Chapter~\ref{sec:generation} pursues further. Promptable, real-time, uncertainty-aware pose is the target the next generation of methods converges on.



% \newpage

\section{Spatial Structure Construction}
\label{sec:p3-sec6-3-spatial-structure-construction}

\FramedInitial{W}{ith} entities recognized and their positions fixed, what remains is how they hang together. Spatial structure construction turns independently localized entities into an organized configuration: which objects relate to which, how those relations compose into a scene, and how the scene as a whole is laid out. The section builds outward in scope, from pairwise spatial relations, through the scene graphs that assemble them, to room-level layout; it then lets structure take on time and scale, as the event structure of action detection and the whole-environment structure of mapping and SLAM. Structure in real perception is seldom purely static, and what the section yields is the explicit spatial configuration that downstream reasoning and planning act on.






\subsection{Spatial Relation Extraction}
\label{sec:spatial_relation_extraction}

% Task definition
%  - Motivation & Importance
%  - Formal Task Definition
%  - Subtasks, if have:

% Part 2: Method taxonomy.
% modality
%

% Task Definition
%  - Motivation & Importance
% ,liu2024edge，liu2025boost
Spatial Relation Extraction (SRE) identifies the relative spatial configuration of objects, such as ``table-left-of-chair'' or ``car-behind-bus''.
It turns detections or segmentations into spatial relation knowledge that feeds reasoning, scene graph generation~\cite{xu2025llava,li2024pixels,wu2025universal}, autonomous driving~\cite{malawade2022roadscene2vec,zhou2023pit,liu2024edge,liu2025boost}, and embodied AI~\cite{SongBTTSB25,zhang2025vlabench}.
It is narrower than Scene Graph Generation (SGG), which runs the full pipeline from detection to relation prediction and includes non-spatial predicates (``person-holding-book''); here SRE covers only geometric, topological, and projective spatial predicates.


%  - Formal Task Definition
\subsubsection{Task Definition}
Given a visual scene (a 2D image, a video, a point cloud, or an RGB-D stream) and two target object instances specified by tags, boxes, or masks, SRE outputs their spatial relationship either as a triplet $(s, p, o)$ with $p$ drawn from a spatial predicate vocabulary, or as a natural-language sentence describing the relation and its context.
Three subtasks differ in output granularity. \textbf{Visual Spatial Classification (VSC)}~\cite{yang2019spatialsense, zhao2024acm} predicts a predicate from a closed label set (``on'', ``in'', ``next to'') or verifies whether a given predicate holds for the object pair. \textbf{Visual Spatial Description (VSD)}~\cite{zhao2022visual, zhao2023generating} generates a free-form sentence that can carry attributes and context beyond the bare predicate. \textbf{Quantitative spatial understanding}~\cite{cheng2024spatialrgpt,chen2024spatialvlm} replaces qualitative predicates with metric quantities (distances, absolute sizes, 3D orientations), which navigation and manipulation require. The field has moved through these three in order, from closed-set classification to generation to metric quantification.


% Part 2: Method taxonomy.
% (Focus on giving a comprehensive but structured overview; no need to cover every paper in depth.) 规则+学习
\subsubsection{Methodology}

Figure~\ref{fig:p3_s6_3_1_sre_overview} summarizes the task and the three method families surveyed below.

\begin{figure}[!t]
    \centering
    \includegraphics[width=\linewidth]{figures/part3/section6/spatial_relation_extraction.pdf}
    \caption{Spatial relation extraction. \textit{Top}: given a visual scene (image, video, point cloud, or RGB-D) and two target objects $(O_1, O_2)$, predict (a) a closed-set relation label (VSC), (b) a free-form description (VSD), or (c) a quantitative metric estimate. \textit{Bottom}: (i) rule-based, (ii) learning-based decoupled, and (iii) MLLM-based mechanisms with method timelines (chronological spacing is not proportional). Metric outputs require depth or scale information.}
    \label{fig:p3_s6_3_1_sre_overview}
\end{figure}

\paragraph{Rule-based approaches.}
Rule-based pipelines~\cite{malawade2022roadscene2vec,zhao2022visual,zhou2024hktsg} first run off-the-shelf perception (a 2D detector with a bird's-eye-view transform~\cite{liu2008bird}, or a 3D scene extractor~\cite{nie2020total3dunderstanding}) to obtain each object's 3D position, size, orientation, and appearance in a camera-centered frame whose vertical axis is aligned with gravity, and then apply deterministic geometric rules to the centroids to produce $(O_1, p, O_2)$ triplets~\cite{zhao2022visual}.
Let $(x_{O_i}, y_{O_i}, z_{O_i})$ denote the centroid of object $O_i$ in this shared camera-centered frame, with $x$ pointing right and $y$ pointing upward. Then
\begin{equation}
\begin{aligned}
\Delta x &= |x_{O_1} - x_{O_2}|, \quad
\Delta y = |y_{O_1} - y_{O_2}|, \quad
\Delta z = |z_{O_1} - z_{O_2}|, \\
&\Delta y > \tau \,\land\, y_{O_1} > y_{O_2} \;\Rightarrow\; \text{``above''}, \\
&\Delta x > \Delta z \,\land\, \Delta x > \tau \,\land\, x_{O_1} < x_{O_2} \;\Rightarrow\; \text{``left''}, \\
&\Delta x > \Delta z \,\land\, \Delta x > \tau \,\land\, x_{O_1} > x_{O_2} \;\Rightarrow\; \text{``right''},
\end{aligned}
\end{equation}
where $\tau$ is a small threshold (\eg, $0.2$\,m).
Such rules are cheap but brittle: 2D-to-3D back-projection misplaces occluded or oddly posed objects (a rear vehicle whose centroid lands ambiguously on the BEV plane), fixed rules cannot express functional or contextual relations (``person riding bicycle'', a book more likely on a shelf than under it), and every new domain needs its rules rewritten.


\paragraph{Learning-based approaches.}
Zhao \etal~\cite{zhao2022visual} decoupled spatial relation generation from generic captioning and trained it with a joint visual--spatial objective; 3DVSD~\cite{zhao2023generating} moved the same task into 3D scene modeling and added structure-aware sampling that holds up under occlusion.

\paragraph{MLLM-based approaches.}
Multimodal large language models carried the line from image-level predicates toward finer and more metric grounding.
SpatialVLM~\cite{chen2024spatialvlm} used monocular-depth-derived spatial QA supervision to train a vision-language model so the model answers metric queries such as ``How far is A from B?'', outperforming GPT-4V on spatial VQA.
SpatialBot~\cite{cai2024spatialbot} exposes pixel-level depth queries through a depth API and trains progressively, translating spatial reasoning into higher robotic grasping success.
SpatialRGPT~\cite{cheng2024spatialrgpt} sharpens granularity from the whole image to regions, accepting boxes or masks and auto-generating 8.7M region-level spatial samples through scene-subgraph sampling (S$^3$).
SpatialLLM~\cite{ma2025spatialllm} pairs a hybrid CLIP--DINOv2 encoder with staged data injection to hold qualitative and quantitative capabilities in one model.


% Part 3: Benchmarks & metrics.
\subsubsection{Datasets and Metrics}
Table~\ref{tab:sre_dataset} lists the benchmarks. They differ mainly in how relations were collected: manual labeling of robot-captured lab scenes in SpatialAwareRobot~\cite{wang2025spatial}; adversarial crowdsourcing in SpatialSense~\cite{yang2019spatialsense}, where annotators propose relations that fool a baseline and only those count, which makes the dataset deliberately hard; the merge of SpatialSense and Visual Genome~\cite{krishna2017visual} into VSD~\cite{zhao2022visual}, whose VSD-v1 version contains over 143K descriptions; the 34-preposition, English--French annotations of SpatialVOC2K~\cite{belz2018spatialvoc2k}; and the automatically generated, depth-aware SpaceSGG~\cite{xu2025llava}, released as 10K scene descriptions, 20K question--answer pairs, and 20K multi-turn conversations over 20K scenes.

\begin{table}[!t]
\centering
\caption{Benchmark datasets for spatial relation extraction. ``--'' marks an open vocabulary without a fixed category count.}
\label{tab:sre_dataset}
\fontsize{9}{11}\selectfont
\setlength{\tabcolsep}{3mm}
\begin{tabular}{lccccc}
\toprule
\textbf{Dataset} & \textbf{Scene Type} & \textbf{Modality} & \textbf{Annotation} & \textbf{Samples} & \textbf{Relation Categories} \\
\midrule
SpatialAwareRobot~\cite{wang2025spatial} & Robotics & Image & Manual & 1,000 & 7 \\
SpatialSense~\cite{yang2019spatialsense} & Generic & Image & Manual & 11,569 & 9 \\
VSD~\cite{zhao2022visual} & Generic & Image & Manual & 29,272  & -- \\
SpatialVOC2K~\cite{belz2018spatialvoc2k} & Generic & Image & Manual & 2,026  & 34 \\
SpaceSGG~\cite{xu2025llava} & Generic & Image & Automated & 20,000  & --  \\
\bottomrule
\end{tabular}
\end{table}

Each subtask has its own evaluation regime, and the metrics are not comparable across them.
VSC is scored as classification: SpatialSense~\cite{yang2019spatialsense} measures binary accuracy for given object--predicate pairs, with the driving-scene systems Roadscene2Vec~\cite{malawade2022roadscene2vec} and HKT-SG~\cite{zhou2024hktsg} reporting per-relation F1 or accuracy on their own data; results on SpatialSense's adversarial split lag well behind those on conventional crowdsourced data, which is the contextual-relation difficulty taken up under Challenges.
VSD is scored as generation: VSDNet~\cite{zhao2022visual} and 3DVSD~\cite{zhao2023generating} report BLEU-4, METEOR, and CIDEr on the VSD datasets, and LLaVA-SpaceSGG~\cite{xu2025llava} reports QA accuracy on SpaceSGG.
Quantitative understanding is the newest regime: SpatialVLM~\cite{chen2024spatialvlm} and SpatialRGPT~\cite{cheng2024spatialrgpt} evaluate metric estimates using human-annotated spatial questions and answers (SpatialVLM) or questions derived from ground-truth 3D cuboid annotations (SpatialRGPT-Bench), SpatialLLM~\cite{ma2025spatialllm} reports combined qualitative and quantitative VQA accuracy, and SpatialBot~\cite{cai2024spatialbot} adds robotic grasp success as an embodied proxy.
MLLM-based models dominate the qualitative leaderboards but still fall short on metric-scale reasoning, which motivates the unified representations discussed next.


\subsubsection{Challenges}
Spatial relations are not purely geometric. ``A cup on a table'' is plausible and ``a cup on a wall'' is not, yet the two differ only in object identity, so rule-based and purely geometric methods cannot capture functional or commonsense constraints; a model needs to know what objects are for and where they tend to be, not just their boxes or point coordinates.
Metric grounding is the second gap. MLLMs handle contextual and functional predicates well but judge distances, sizes, and angles poorly; injecting geometry and depth during pre- or post-training, or exposing dedicated spatial tools, is the current route, and the longer-term goal is a representation that holds discrete predicates (``near'') and continuous quantities (``2 meters away'') together, since navigation and manipulation need both.
Generalization is the third. SpatialSense~\cite{yang2019spatialsense} and VSD~\cite{zhao2022visual} cover a small slice of real scenes, and models trained on them degrade under new object appearances, scene complexity, lighting, and occlusion; SpatialVLM~\cite{chen2024spatialvlm} and its successors generalize better, but unseen and dynamic environments remain largely untested.





\subsection{Scene Graph Generation}
\label{sec:scene_graph_generation}
% - `Writing task`: 任务定义（TSG、VSG、PSG、HOI / HRI (egocentric)） → 表达与评测 → 方法类别 → 数据与协议 → 难点与趋势
% - Content:
%- Table:  SGG datasets
% - Figure:  三种SG示意图
% https://github.com/ChocoWu/Awesome-Scene-Graph-Generation


Scene graph generation (SGG)~\cite{xu2017scene,chang2021comprehensive,li2024scene} parses visual content into a graph of triplets, each joining a subject and an object through a predicate, for example $(\text{person}, \text{rides}, \text{horse})$.
The graph records not only which objects are present but how they interact, which is what retrieval~\cite{johnson2015image,miao2024scenegraphloc}, captioning~\cite{zhong2020comprehensive,lu2025benchmarking}, autonomous driving~\cite{zhou2022toward,liu2023learning,malawade2022roadscene2vec}, and robotics~\cite{nguyen2025generating,werby2024hierarchical,garg2024robohop} consume downstream.


%  - Formal Task Definition
\subsubsection{Task Definition}
Given a visual scene, SGG produces a graph
\begin{equation}
\mathcal{G} = \{(s_i, p_i, o_i)\}_{i=1}^{M}, \qquad s_i, o_i \in \mathcal{O},\; p_i \in \mathcal{P},
\end{equation}
where each entity in $\mathcal{O}$ is an object instance with a localization (a bounding box, or a mask in the panoptic setting) and a class label, and $p_i$ is a predicate from a relation vocabulary $\mathcal{P}$ describing how $s_i$ and $o_i$ interact.

% 任务分类：
% 1. Textual Scene Graph Generation
% 2. 2D Image Scene Graph Generation
% 2.1  Panoptic Scene Graph Generation
% 3. Spatio-Temporal (Video) Scene Graph Generation
% 4. 3D/ Scene Graph Generation
% 5. Universal Scene Graph Generation
% 6. HOI
% 7. Ego-SGG

The task is instantiated by input modality. \textbf{Image SGG (ISG)}~\cite{krishna2017visual,zellers2018neural} on static 2D images is the original and most studied setting and anchors most benchmarks. \textbf{Video SGG (VSG)}~\cite{teng2021target, cong2021spatial} adds time, so relations become dynamic and temporally extended, and temporal context disambiguates interactions a single frame leaves open. \textbf{Textual SGG (TSG)}~\cite{li2023factual,choi2022scene} parses a sentence (``A person is riding a horse in a field'') into a graph, treating SGG as structured language parsing. \textbf{3D SGG (3DSG)}~\cite{wald2020learning,wu2021scenegraphfusion} operates on point clouds or meshes and models relations in metric 3D space, free of the perspective ambiguity of images. \textbf{Universal SGG (USG)}~\cite{wu2025universal} aims at one graph from several simultaneous modalities (image, video, 3D, text).
Two further variants change the level of detail rather than the modality. \textbf{Human--Object Interaction detection (HOI)}~\cite{li2019transferable, wang2020learning} restricts the subject to humans and has become a field of its own because of pose variability and the subtlety of intent. \textbf{Panoptic SGG (PSG)}~\cite{yang2022panoptic,yang2023panoptic} grounds nodes on panoptic segmentation masks instead of boxes, so relations can involve ``stuff'' (sky, road) and are localized without the imprecision of boxes.

\subsubsection{Methodology}
% 1.按模态分（Image、video、textual、3D、Universal Scene Graph Generation）
Figure~\ref{fig:p3_s6_3_2_sgg_overview} illustrates the five modality settings and the two specializations with a representative input and output each.

\begin{figure}[!t]
    \centering
    \includegraphics[width=\linewidth]{figures/part3/section6/scene_graph.pdf}
    \caption{Scene graph generation by input modality (top) and sub-task setting (bottom). \textit{Top}: image, video, text, 3D, and universal (multimodal) SGG, each with a sample input and output graph. \textit{Bottom}: (i) human--object interaction (HOI), restricted to human subjects at box level, and (ii) panoptic SGG (PSG), grounding ``things'' and ``stuff'' on pixel masks.}
    \label{fig:p3_s6_3_2_sgg_overview}
\end{figure}

% 2.two-stage —> one-stage
\paragraph{Two-stage methods.}
Two-stage methods~\cite{xu2017scene,zellers2018neural,li2021bipartite,lin2020gps,li2023compositional,zhou2024hierarchical} were the standard for years: a pre-trained detector such as Faster R-CNN~\cite{ren2015faster} produces boxes and per-object features, and a separate relation model predicts a predicate for each pair. Specializing each stage made the graphs accurate, but detector errors propagate uncorrected into relation prediction, the two models are trained and tuned separately, and object interactions are only modeled after detection has already committed.

\paragraph{One-stage methods.}
One-stage methods~\cite{im2024egtr,wang2024oed,newell2017pixels,liu2021fully} fold detection and relation prediction into one end-to-end network. Pixel-based~\cite{newell2017pixels} and fully convolutional~\cite{liu2021fully} formulations came first; after DETR~\cite{carion2020end}, most one-stage SGG adopted set prediction with learned queries, and the methods differ in what the queries represent. RelTR~\cite{cong2023reltr} uses paired subject--object queries so the Transformer predicts relations between pairs directly; SGTR~\cite{li2022sgtr} decouples subject, object, and predicate queries and composes them, which lets each component be predicted on its own; and PSGTR~\cite{yang2022panoptic} uses triplet queries that emit scene-graph triplets without a detector at all.

% 3.引入先验知识，LLM
\paragraph{Prior knowledge.}
Because predicates are long-tailed and many relations are not visually decidable, SGG has long imported priors. Language priors~\cite{liang2017deep,lu2016visual,cui2018context} supply semantic embeddings that help rare predicates; statistical priors~\cite{yu2017visual,zhou2023learning,zhou2024hierarchical} model object-pair co-occurrence; knowledge graphs such as ConceptNet~\cite{speer2017conceptnet} add commonsense~\cite{zhang2024hiker,chen2023more}; and LLMs bring world knowledge~\cite{chen2024scene,kim2024llm4sgg}.

\paragraph{Beyond images.}
The non-image settings reuse these designs and add modality-specific machinery. Video SGG tracks object tubes and aggregates relations over time: TRACE~\cite{teng2021target} decomposes relation prediction by temporal segment and STTran~\cite{cong2021spatial} applies a spatio-temporal Transformer over frame-level graphs. Textual SGG is a parsing problem, for which the FACTUAL parser~\cite{li2023factual} fine-tunes a sequence-to-sequence model to emit faithful triples. 3D SGG predicts relations from point-cloud or mesh segments with graph networks~\cite{wald2020learning}, and SceneGraphFusion~\cite{wu2021scenegraphfusion} does so incrementally from an RGB-D stream so the graph grows with the reconstruction. Panoptic SGG is served by the triplet-query PSGTR~\cite{yang2022panoptic} and its video extension~\cite{yang2023panoptic}, and USG~\cite{wu2025universal} shares one graph decoder across image, video, 3D, and text inputs.

% TODO add HOI，Traffic SG dataset.

\subsubsection{Datasets and Metrics}


\begin{table}[!t]
  \centering
\caption{Benchmark datasets for scene graph generation.}
\label{tab:SGG_data}
\fontsize{8}{10}\selectfont
\setlength{\tabcolsep}{2.5mm}
  \begin{tabular}{lcccccc}
    \toprule
    \textbf{Dataset} & \textbf{Modality} &\textbf{ Obj. Class} & \textbf{BBox} & \textbf{Rela. Class} & \textbf{Triplets} & \textbf{Instances} \\
    \midrule
    Visual Phrase~\cite{Sadeghi2011VisualPhrases} & Image & 8 & 3,271 & 9 & 1,796 & 2,769 \\
    Scene Graph~\cite{johnson2015image} & Image & 266 & 69,009 & 68 & 109,535 & 5,000 \\
    VRD~\cite{lu2016visual} & Image & 100 & - & 70 & 37,993 & 5,000 \\
    Open Images v7~\cite{kuznetsova2020open} & Image & 600 & 3,290,070 & 31 & 374,768 & 9,178,275 \\
    Visual Genome~\cite{krishna2017visual} & Image & 5,996 & 3,843,636 & 1,014 & 2,347,187 & 108,077 \\
    GQA~\cite{hudson2019gqa} & Image & 1,740 & - & 330 & - & 113,018 \\
    VrR-VG~\cite{liang2019vrr} & Image & 1,600 & 282,460 & 117 & 203,375 & 58,983 \\
    UnRel~\cite{peyre2017weakly} & Image & - & - & 18 & 76 & 1,071 \\
    SpatialSense~\cite{yang2019spatialsense} & Image & 3,679 & - & 9 & 13,229 & 11,569 \\
    SpatialVOC2K~\cite{belz2018spatialvoc2k} & Image & 20 & 5,775 & 34 & 9,804 & 2,026 \\
    PSG~\cite{yang2022panoptic} & Image (panoptic) & 133 & - & 56 & - & 49K \\
    AUG~\cite{li2024aug} & Image (overhead) & 76 & - & 61 & - & - \\
    STAR~\cite{li2025star} & Image (satellite) & 48 & 219,120 & 58 & 400,795 & 31,096 \\
    ReCon1M~\cite{yan2025recon1m} & Image (satellite) & 60 & 859,751 & 64 & 1,149,342 & 21,392 \\
    SkySenseGPT~\cite{luo2024skysensegpt} & Image (satellite, instruction) & - & - & - & - & - \\
    Traffic Scene Graph~\cite{zhou2022toward} & Image (traffic) & 2,266 & - & 4,272 & - & 451 \\
    ImageNet-VidVRD~\cite{shang2017video} & Video & 35 & - & 132 & 3,219 & 1,000 \\
    VidOR~\cite{shang2019annotating} & Video & 80 & - & 50 & - & 10,000 \\
    Action Genome~\cite{ji2020action} & Video & 35 & 0.4M & 25 & 1.7M & 10,000 \\
    AeroEye~\cite{nguyen2024cyclo} & Video (drone) & 56 & - & 384 & - & 2.2M \\
    PVSG~\cite{yang2023panoptic} & Video (panoptic) & 126 & - & 57 & 4,587 & 400 \\
    ASPIRe~\cite{nguyen2024hig} & Video (interlacements) & - & - & 4.5K & - & 1.5K \\
    Ego-EASG~\cite{rodin2024action} & Video (egocentric) & 407 & - & 235 & - & - \\
    IESG~\cite{liu2023learning} & Video (traffic) & 9 & - & 16 & - & - \\
    3DSSG~\cite{wald2020learning} & 3D & 528 & - & 39 & - & 48K \\
    PSG4D~\cite{yang20234d} & 4D & 46 & - & 15 & - & - \\
    4D-OR~\cite{ozsoy20224d} & 4D (operating room) & 12 & - & 14 & - & - \\
    MM-OR~\cite{ozsoy2025mm} & 4D (operating room) & - & - & - & - & - \\
    EgoExOR~\cite{ozsoy2025egoexor} & 4D (operating room) & 36 & - & 22 & 568,235 & - \\
    TSG Bench~\cite{yang2025llm} & Text & - & - & - & 11,820 & 4,289 \\
    FACTUAL~\cite{li2023factual} & Text, Image & 4,042 & - & 1,607 & 40,149 & 40,369 \\
    DiscoSG-DS~\cite{lin2025discosg} & Text, Image & 4,018 & - & 2,033 & 68,478 & 8,830 \\
    \bottomrule
  \end{tabular}
\end{table}

Table~\ref{tab:SGG_data}\footnote{The table is extended from the GitHub repository: \url{https://github.com/ChocoWu/Awesome-Scene-Graph-Generation}} lists the benchmarks by modality.
The image benchmarks run from the small precursors Visual Phrase~\cite{Sadeghi2011VisualPhrases}, Scene Graph~\cite{johnson2015image}, and VRD~\cite{lu2016visual}, which first exposed the long-tailed predicate distribution, to Visual Genome~\cite{krishna2017visual}, Open Images~\cite{kuznetsova2020open}, and GQA~\cite{hudson2019gqa}, with VrR-VG~\cite{liang2019vrr} and UnRel~\cite{peyre2017weakly} carved out to test visually relevant and zero-shot relations respectively.
Visual Genome is the primary benchmark but its raw release (33,877 object and 40,480 relationship synsets) is noisy and heavily long-tailed, so the community works with cleaned subsets: the broad cleaning in Table~\ref{tab:SGG_data} keeps roughly 6K object and 1K predicate classes, and the de-facto SGG benchmark VG150 restricts these to the 150 most frequent object and 50 most frequent predicate classes.
Spatial-only relations are covered by SpatialSense and SpatialVOC2K (see the preceding subsection), panoptic grounding by PSG~\cite{yang2022panoptic}, and non-standard viewpoints by the overhead AUG~\cite{li2024aug}, the satellite STAR~\cite{li2025star}, ReCon1M~\cite{yan2025recon1m}, and instruction-based SkySenseGPT~\cite{luo2024skysensegpt}, and the traffic-scene graphs of~\cite{zhou2022toward}.
Video benchmarks began with ImageNet-VidVRD~\cite{shang2017video} and the trajectory-annotated VidOR~\cite{shang2019annotating}, then specialized to human activity (Action Genome~\cite{ji2020action}), drone footage (AeroEye~\cite{nguyen2024cyclo}), panoptic video (PVSG~\cite{yang2023panoptic}), fine topological interlacements (ASPIRe~\cite{nguyen2024hig}), egocentric capture (Ego-EASG~\cite{rodin2024action}), and imminent traffic events (IESG~\cite{liu2023learning}).
3DSSG~\cite{wald2020learning} is the standard 3D benchmark; PSG4D~\cite{yang20234d} adds time, and the operating-room sets 4D-OR~\cite{ozsoy20224d}, MM-OR~\cite{ozsoy2025mm}, and EgoExOR~\cite{ozsoy2025egoexor} bring multimodal (RGB-D, audio, robot logs) and egocentric surgical scenes.
On the text side, TSG Bench~\cite{yang2025llm} evaluates LLMs on graph generation from narratives, while FACTUAL~\cite{li2023factual} and DiscoSG-DS~\cite{lin2025discosg} pair images with text for factually consistent, cross-modal graphs.

Evaluation separates what is given from what is predicted. In \textbf{Predicate Classification (PredCls)} boxes and object labels are given and only predicates are predicted; in \textbf{Scene Graph Classification (SGCls)} boxes are given and both object labels and predicates are predicted; in \textbf{Scene Graph Generation (SGGen, also SGDet)} everything is predicted, and a triplet counts only if both boxes overlap their ground truth at IoU $\geq 0.5$, with \textbf{Phrase Detection} relaxing this to the union box and SGGen+ also crediting correctly detected singleton nodes. Panoptic SGG replaces box IoU with mask IoU.
The standard metric is recall of ground-truth triplets among the top-$K$ predictions,
\begin{equation}
    R@K = \frac{|Top_{K} \cap GT|}{|GT|},
\end{equation}
under a graph constraint (one predicate per pair) or without it. Because $R@K$ is dominated by frequent predicates, \textbf{mean Recall@K} ($mR@K$) averages recall per predicate class and has become the long-tail metric of record; the gap between $R@K$ and $mR@K$ drove much of the debiasing literature~\cite{lin2020gps,li2021bipartite,zhou2024hierarchical}. Weighted triplet recall, micro- versus macro-averaging over images, the stricter $mAP_{rel}$ / $mAP_{phr}$ of the Open Images challenge, and retrieval-based R-Precision are reported in specific settings.

Each modality has its own benchmark and protocol, so leaderboards are not comparable across types: video SGG uses trajectory-based evaluation on ImageNet-VidVRD and VidOR, while Action Genome commonly uses frame-level PredCls, SGCls, and SGDet; textual SGG uses set F1 over triples on FACTUAL and LLM-judged action coverage on TSG Bench; 3D SGG uses object-, predicate-, and triplet-level recall on 3DSSG; panoptic SGG uses mask-IoU $R@K$ on PSG and PVSG. Table~\ref{tab:sgg_performance} therefore reports only the image setting, where VG150 SGDet is shared.

\begin{table}[!t]
\centering
\caption{Representative image SGG methods on Visual Genome (VG150), scene graph detection (SGDet) with graph constraint, as reported in the comparison of EGTR~\cite{im2024egtr}. Two-stage methods use a Faster R-CNN detector with a ResNeXt-101-FPN backbone; one-stage methods are DETR-based with the backbone shown. ``--'' = not reported.}
\label{tab:sgg_performance}
\fontsize{9}{11}\selectfont
\setlength{\tabcolsep}{3.0mm}
\begin{tabular}{llccccc}
\toprule
\textbf{Method} & \textbf{Family / backbone} & \textbf{Year} & \textbf{R@50} & \textbf{R@100} & \textbf{mR@50} & \textbf{mR@100} \\
\midrule
Neural Motifs~\cite{zellers2018neural} & Two-stage, X-101-FPN & 2018 & 32.1 & 36.9 & 5.5 & 6.8 \\
BGNN~\cite{li2021bipartite}            & Two-stage, X-101-FPN & 2021 & 31.0 & 35.8 & 10.7 & 12.6 \\
SGTR~\cite{li2022sgtr}                 & One-stage, DETR-R101 & 2022 & 24.6 & 28.4 & 12.0 & 15.2 \\
RelTR~\cite{cong2023reltr}             & One-stage, DETR-R50  & 2023 & 27.5 & -- & 10.8 & -- \\
EGTR~\cite{im2024egtr}                 & One-stage, Deformable DETR-R50 & 2024 & 30.2 & 34.3 & 7.9 & 10.1 \\
\bottomrule
\end{tabular}
\end{table}

Table~\ref{tab:sgg_performance} shows the trade-off that still shapes image SGG. The two-stage Motifs baseline keeps the highest overall recall, but its mean recall is the lowest, which is the frequent-predicate bias in numbers; BGNN's debiasing doubles mean recall at a small cost in $R@K$. The one-stage DETR-style models start several recall points lower, since they must detect and relate in one pass with a weaker backbone, yet SGTR and RelTR reach the best mean recall of the group, and EGTR closes most of the overall-recall gap to two-stage methods while running end to end.


\subsubsection{Challenges}
% 1.单模态SGG -> 统一SGG
% 2.与MLLM结合 构建结构化mllm
% 3.单粒度/层次化场景图 -> 无限粒度场景图
% 4.下游具身应用

\paragraph{From single-modality to universal SGG.}
Methods are built per modality, yet real scenes are multimodal, with images carrying visual detail and text carrying abstract concepts and events. A universal scene graph~\cite{wu2025universal} would express all of them in one cross-modal representation; the hard part is identifying the same objects, events, and relations across modalities and suppressing cross-modal noise, and granularity has to follow the application rather than a fixed ontology, since robotic interaction or driving may need relations at a level of detail that image-level graphs do not provide.

\paragraph{Integrating SGG with MLLMs.}
MLLMs are strong at language and weaker at structured visual understanding. Scene graphs offer a way in: pretraining objectives that mask and predict graph nodes (objects, attributes, relations) across image, video, and text inputs would teach a language model to read structure rather than only the next token, and the scene-graph-guided MLLMs of the preceding subsection~\cite{xu2025llava} are an early instance.

\paragraph{From images to embodied use.}
Most SGG analyzes static images or videos, but the graph is most useful when it is kept live: robots update it as they act~\cite{nguyen2025generating,werby2024hierarchical,garg2024robohop}, and driving systems read relations among lanes, signs, and pedestrians from it~\cite{zhou2022toward,liu2023learning,malawade2022roadscene2vec}. As these uses grow, the scene graph shifts from a product of perception to the state on which decision and interaction operate.





\subsection{Layout Estimation and Room Parsing}
\label{sec:layout_estimation}

Layout estimation recovers the structural geometry of an indoor space, the corners, walls, floor, and ceiling that bound it, as a simplified 3D ``wireframe'' called the room layout.
Room parsing extends this by also detecting, localizing, and labeling the objects inside (furniture and fixtures), giving a complete description of the scene.
For outdoor scenes, the analogous problem is recovering a road layout in bird's-eye view from a street image~\cite{YangLLYMHP21}.
Figure~\ref{fig:p3_s6_3_3_layout_overview} shows the inputs and output representations for the structural part of the task.

\begin{figure}[!t]
    \centering
    \includegraphics[width=0.99\linewidth]{figures/part3/section6/layout_parsing.pdf}
    \caption{Inputs and output representations of structural layout estimation. Top: (i) indoor perspective images, (ii) 360$^\circ$ panoramas, and (iii) outdoor street views. Bottom: (a) polygonal room layouts and (b) schematic indoor/outdoor surface-label maps, illustrating representations that may be predicted separately or jointly. Object-level parsing and outdoor bird's-eye-view outputs are not depicted.}
    \label{fig:p3_s6_3_3_layout_overview}
\end{figure}

\subsubsection{Task Definition}
Given an image $\mathbf{I}\in\mathbb{R}^{H\times W\times 3}$ (a perspective view, a $360^\circ$ panorama, or a multi-view or video capture), layout estimation recovers the dominant surfaces that bound the navigable volume of a room as a \emph{room layout}
$\mathcal{L}=\{P_k\}_{k=1}^{K}$, where each $P_k$ is a planar polygon in 3D; equivalently the output may be plane parameters $(\mathbf{n}_k,d_k)$ or a per-pixel structural-surface label map. Under a vertical-wall model, the layout compresses to a floor-plan polygon with floor and ceiling heights $y_f, y_c$. Outdoor road-layout methods instead predict road regions and vehicle occupancy in bird's-eye view.
Room parsing adds per-object instance masks $\{m_i\}$ or boxes with class labels, drawn from a fixed ontology or, in the open-vocabulary setting, from free-form language.

\subsubsection{Methodology}

\paragraph{Pre-deep-learning approaches.}
Classical room layout estimation rests on the \textbf{Manhattan assumption}~\cite{CoughlanY99}: the scene is built from surfaces aligned with three mutually orthogonal directions, as walls, floors, and ceilings in most built environments are.
Hedau \etal~\cite{HedauHF09} introduced the task, modeling a cluttered indoor image jointly as a 3D box layout (the room as if empty) and per-pixel surface labels (visible walls, floor, ceiling, and objects).
Three mutually orthogonal vanishing points are first estimated with a robustified variant of Rother's method~\cite{Rother02}, in which a line segment $l$ votes for a candidate vanishing point $p$ with
\begin{equation}
v(l,p) = |l| \cdot \exp\!\left(-\frac{\alpha(l,p)^2}{2\sigma^2}\right),
\end{equation}
where $\alpha(l,p)$ is the angle between $l$ and the line from the midpoint of $l$ to $p$, $|l|$ is the segment length, and $\sigma$ is a fixed bandwidth; candidates are selected greedily, removing the lines that supported each selected point.
The vanishing points fix the orientation of the box; sampling pairs of rays through two of them generates box-layout hypotheses, whose remaining edges follow from the third vanishing point, and a learned structured regressor scores the hypotheses and selects the layout.

\paragraph{Deep learning approaches.}
Mallya and Lazebnik~\cite{MallyaL15} first used a fully convolutional network to predict informative edge maps, then kept the classical hypothesis-generation-and-ranking step, scoring hypotheses with features from the learned edges together with line membership and geometric context.
Later deep estimators are best organized by input modality and geometric representation, which vary independently: panorama input does not imply a cuboid output, and non-cuboid geometry need not break the Manhattan assumption. Figure~\ref{fig:p3_s6_3_3_methodology_overview} groups the representative methods of 2015--2021 along these lines.

\begin{figure}[!t]
    \centering
    \includegraphics[width=\linewidth]{figures/part3/section6/layout_parsing_method.pdf}
    \caption{Representative layout-estimation methods, 2015--2021. Each lane pairs a core-mechanism illustration with a publication timeline: (i) perspective corner and edge prediction, (ii) panorama representations, (iii) beyond-cuboid floor plans, and (iv-a, iv-b) indoor plane geometry and outdoor cross-view prediction. Timeline spacing is not proportional; arrows indicate chronology, not method dependencies.}
    \label{fig:p3_s6_3_3_methodology_overview}
\end{figure}

\textbf{(i) Perspective-image corner and edge predictors.}
RoomNet~\cite{lee2017roomnet} regresses an ordered set of room corners with a recurrent encoder--decoder conditioned on one of 11 predefined cuboid topologies, and LayoutNet~\cite{ZouCSH18} predicts corner and boundary probability maps and fits a cuboid to them. The topology constraint simplifies reconstruction, but a single perspective view often leaves part of the room outside the frame.

\textbf{(ii) Panorama-based methods.}
A $360^\circ$ equirectangular panorama removes the field-of-view truncation of a perspective image, though not occlusion. LayoutNet~\cite{ZouCSH18} predicts 2D corner and boundary maps from a gravity-aligned panorama; HorizonNet~\cite{sun2019horizonnet} compresses the layout into three 1D per-column signals (floor-wall boundary, ceiling-wall boundary, and wall-wall corner indicator), adds Pano Stretch augmentation, and reconstructs general Manhattan layouts as well as cuboids; HoHoNet~\cite{sun2021hohonet} learns latent horizontal features decodable into layout, depth, and semantics; and LED2-Net~\cite{wang2021led2net} supervises the layout through differentiable depth rendering, comparing depth rendered from the predicted layout with the target geometry.

\textbf{(iii) Beyond-cuboid and Atlanta-world methods.}
Floor-plan prediction separates the number of walls from the constraints on their orientation. DuLa-Net~\cite{yang2019dulanet} combines equirectangular and ceiling--floor projections to estimate a floor plan under a Manhattan prior, which admits non-rectangular rooms with orthogonal walls. AtlantaNet~\cite{pintore2020atlantanet} projects a gravity-aligned panorama onto horizontal planes above and below the camera and predicts a room footprint and height under the Atlanta-world model, which keeps walls vertical but allows arbitrary horizontal orientation, including curved footprints. Both still assume vertical walls, so slanted walls and ceilings remain out of reach.

\textbf{(iv) Plane-based indoor geometry and outdoor cross-view prediction.}
GeoLayout~\cite{0005ZZ20} predicts pixel-level plane parameters, infers the depth map of each dominant plane, and reads layout boundaries off the intersections of those depth maps, a geometric alternative to regressing polygon vertices; the accompanying Matterport3D-Layout dataset covers cuboid and non-cuboid rooms. For outdoor scenes, Yang \etal~\cite{YangLLYMHP21} use a cross-view transformation with cycle-consistency constraints to predict road layout and vehicle occupancy in bird's-eye view from a monocular front-facing image.

\subsubsection{Datasets and Metrics}
Table~\ref{tab:layout_estimation_room_parsing_dataset} summarizes the benchmarks. Perspective-image evaluation grew from Hedau~\cite{HedauHF09} (308 web and LabelMe images with box layouts and surface labels) through its augmented successor Hedau+ and the more cluttered SUNbox drawn from SUN2012~\cite{XiaoHEOT10}, both built by Mallya and Lazebnik~\cite{MallyaL15} to train FCN edge predictors, to LSUN Room Layout~\cite{yu2015lsun}, the standard perspective benchmark with corner positions and cuboid topologies on over 5K images. Corridor and video data serve navigation: Yang~\cite{YangMS16} labels ground versus wall in corridor images, FreeYourCamera~\cite{FurlanMSLS13} records six-degree-of-freedom smartphone walks that frequently violate box and Manhattan assumptions, and MichiganIndoorCorridor~\cite{TsaiXLK11} provides non-Manhattan corridors with glass and reflections. Panoramic and 3D evaluation rests on PanoContext~\cite{zhang2014panocontext} (cuboid layouts and furniture for about 500 SUN360 panoramas, the standard testbed for LayoutNet, HorizonNet, and LED2-Net), Stanford 2D-3D-S~\cite{armeni2017joint} (70,496 perspective RGB images and 1,413 RGB panoramas, with aligned depth, normals, semantics, and 3D meshes over 270 rooms), and Matterport3D-Layout~\cite{0005ZZ20}, derived from Matterport3D~\cite{chang2017matterport3d} with RANSAC-fitted plane ground truth for cuboid and non-cuboid rooms.

\begin{table*}[!t]
    \centering
    \caption{Layout estimation and room parsing benchmarks. Scale is in images unless noted.}
    \label{tab:layout_estimation_room_parsing_dataset}
    \fontsize{9}{11}\selectfont
    \setlength{\tabcolsep}{2.0mm}
    \begin{tabular}{L{0.20\textwidth} c C{0.08\textwidth} C{0.11\textwidth} C{0.12\textwidth} L{0.27\textwidth}}
    \toprule
    \textbf{Dataset} & \textbf{Year} & \textbf{Scene} & \textbf{Modality} & \textbf{Scale} & \textbf{Annotation} \\
    \midrule
    Hedau \cite{HedauHF09} & 2009 & Room & Image & 308 & Box layout + surface labels \\
    Hedau+ \cite{MallyaL15} & 2015 & Room & Image & 337 & Box layout (Hedau + 75 bedroom images) \\
    SUNbox \cite{MallyaL15} & 2015 & Room & Image & 543 & Box layout, cluttered SUN2012 scenes \\
    LSUN Room Layout \cite{yu2015lsun} & 2015 & Room & Image & 5,394 & Corners + 11 cuboid topologies \\
    Yang \cite{YangMS16} & 2016 & Corridor & RGB-D image & 967 & Ground / non-ground masks \\
    FreeYourCamera \cite{FurlanMSLS13} & 2013 & Mixed & Video (6-DoF) & 10 seq.\ (84 frames) & Layout under non-box, non-Manhattan motion \\
    MichiganIndoorCorridor \cite{TsaiXLK11} & 2011 & Corridor & Video & 11 seq.\ ($\sim$3,300 frames) & Ground + three walls on first frame \\
    PanoContext \cite{zhang2014panocontext} & 2014 & Room & Panorama & $\sim$500 & Cuboid layout + furniture \\
    Matterport3D-Layout \cite{0005ZZ20} & 2020 & Room & RGB-D image & 7,360 & Plane parameters, cuboid + non-cuboid \\
    Stanford 2D-3D-S \cite{armeni2017joint} & 2017 & Building & RGB-D & 70,496 views; 1,413 panoramas & Depth, normals, semantics, mesh \\
    \bottomrule
    \end{tabular}
\end{table*}

Metrics follow the output representation. Perspective benchmarks (Hedau, LSUN) report \emph{pixel error}, the fraction of pixels whose predicted surface label differs from the ground truth, and \emph{corner error}, the mean distance between predicted and true corners normalized by image size; the edge-map methods additionally report the contour-detection F-scores at the optimal dataset scale (ODS, one threshold for the whole set) and optimal image scale (OIS, the best threshold per image). Panoramic and 3D benchmarks (PanoContext and panoramic layout subsets of Stanford 2D-3D-S and Matterport3D) report 3D and 2D \emph{room IoU} between the predicted and ground-truth layout volumes or floor plans, with frequency-weighted IoU and pixel accuracy as segmentation-style complements. A method trained for one input type rarely transfers to the other without re-training, so perspective and panoramic results are not compared directly, and the outdoor road-layout protocol~\cite{YangLLYMHP21} is separate again.

\begin{table}[!t]
    \centering
    \caption{Panorama-based layout estimation, 3D IoU (\%) on PanoContext, Stanford 2D-3D-S, and the 2,295-panorama Matterport3D layout subset (overall), as reported in the comparison of LED2-Net~\cite{wang2021led2net}. Higher is better.}
    \label{tab:layout_performance}
    \fontsize{9}{11}\selectfont
    \setlength{\tabcolsep}{1.5mm}
    \begin{tabular}{L{0.16\textwidth} L{0.28\textwidth} c C{0.10\textwidth} C{0.13\textwidth} C{0.14\textwidth}}
    \toprule
    \textbf{Method} & \textbf{Representation} & \textbf{Year} & \textbf{PanoContext} & \textbf{Stanford 2D-3D-S} & \textbf{Matterport3D} \\
    \midrule
    LayoutNet~\cite{ZouCSH18}              & 2D corner + boundary maps      & 2018 & 74.48 & 76.33 & 75.82 \\
    DuLa-Net~\cite{yang2019dulanet}         & Dual projection, Manhattan plan & 2019 & 77.42 & 79.36 & 75.05 \\
    HorizonNet~\cite{sun2019horizonnet}     & 1D column representation       & 2019 & 82.17 & 79.79 & 78.73 \\
    AtlantaNet~\cite{pintore2020atlantanet} & Top-down footprint, Atlanta world & 2020 & 78.76 & 82.43 & 80.02 \\
    LED2-Net~\cite{wang2021led2net}         & Differentiable depth rendering & 2021 & \textbf{82.75} & \textbf{83.77} & \textbf{81.52} \\
    \bottomrule
    \end{tabular}
\end{table}

Table~\ref{tab:layout_performance} compares the panorama methods under one protocol. The 1D column representation of HorizonNet gained about eight IoU points over LayoutNet's 2D maps on PanoContext, and the geometric supervision of LED2-Net adds a further point or two on every benchmark; AtlantaNet exceeds HorizonNet on the Matterport3D subset, whereas DuLa-Net remains below it, so these results do not show a uniform advantage for beyond-cuboid methods.

\subsubsection{Challenges}
The geometric priors that make layout estimation tractable also bound it. Cuboid models cannot express alcoves or L-shaped plans; Manhattan models keep orthogonal walls even when the footprint is free-form; the Atlanta-world model of AtlantaNet~\cite{pintore2020atlantanet} frees wall orientation but not wall tilt; and plane-based methods such as GeoLayout~\cite{0005ZZ20} approximate curved geometry only with many planes. Slanted enclosures, split levels, and the outdoor bird's-eye-view target of~\cite{YangLLYMHP21}, which is a different output space altogether, all sit outside the standard formulations.
Clutter and occlusion are the second obstacle: furniture hides the corners, boundaries, and vanishing lines that both classical and learned methods read, so estimators must hallucinate hidden structure from layout priors, and confidently wrong boundaries without uncertainty are the typical failure.
Generalization follows from the data: benchmarks are indoor, mostly rooms and corridors, and models trained on one set degrade on another's appearance, architecture, and lighting, while sparse, unposed multi-view capture breaks the pose assumptions of classical pipelines; the feed-forward geometry models of \S\ref{sec:visual_odometry} are the natural starting point for layout from such input.





\subsection{Action Detection}
\label{sec:action_detection}
% \wsq{人体动作行为识别与时序定位}

Action detection (AD) finds human actions in untrimmed video: for each action instance it returns the time interval $[t_s, t_e]$ in which it occurs and its category $c$. Where action recognition classifies a clip that has already been trimmed to one action, detection must also decide when actions start and end, in streams that contain many of them.
It belongs to spatial structure construction because it organizes the entities and relations of a scene along the time axis into events, the temporal counterpart of the scene graphs above and the basis of video SGG and human--object interaction.
Figure~\ref{fig:p3_s6_3_4_ad_overview} illustrates the task.


\begin{figure}[!ht]
    \centering
    \includegraphics[width=\linewidth]{figures/part3/section6/action_detection.pdf}
    \caption{Action detection: each action instance in an untrimmed video is localized to $[t_s,t_e]$ and labeled with class $c$ (top). \textit{Bottom}: three task illustrations: (i) action classification of a trimmed clip, (ii) temporal action localization (illustrated with one interval), and (iii) temporal action detection of multiple instances.}
    \label{fig:p3_s6_3_4_ad_overview}
\end{figure}

\subsubsection{Task Definition}
\label{sec:6_3_4_definition}
Given an untrimmed video (or a continuous multi-sensor stream), AD outputs a set of tuples $(t_s, t_e, c)$ with start time, end time, and category for every action instance.
The following terms emphasize different aspects of the task. \textbf{Action classification (AC)} takes a trimmed clip and predicts one category, with no localization. \textbf{Temporal action localization (TAL)} predicts start and end times in untrimmed video, and may include multiple instances; the term is often used interchangeably with TAD. \textbf{Temporal action detection (TAD)} emphasizes joint localization and classification of multiple instances per video, possibly overlapping or concurrent, across temporal scales, each to be localized and classified.


\begin{figure}[!ht]
    \centering
    \includegraphics[width=0.99\linewidth]{figures/part3/section6/action_detection_method_1.pdf}
    \caption{Representative methods and mechanisms for (i) action classification, (ii) temporal localization with boundary estimation and alternative supervision signals, and (iii) temporal detection with anchor-based, anchor-free, boundary-based, and query-based formulations. Method timelines show publication order, not computational dependencies; spacing is not proportional. Same-year methods are grouped. The localization and detection panels emphasize different aspects of closely related tasks.}
    \label{fig:p3_s6_3_4_methodology_overview}
\end{figure}

\subsubsection{Methodology}

Figure~\ref{fig:p3_s6_3_4_methodology_overview} organizes the methods by the three settings; each builds on the previous one's representations.

\paragraph{Action Classification.}
Two-stream CNNs~\cite{simonyan2014two} model appearance from RGB frames in one stream and motion from dense optical flow in another, fusing the two at the output. 3D CNNs~\cite{diba2017temporal} convolve over space and time jointly, C3D~\cite{tran2015learning} and I3D~\cite{carreira2017quo} being the standard examples, at the cost of a large parameter count that R(2+1)D~\cite{tran2018closer} and P3D~\cite{qiu2017learning} reduce by factoring 3D convolution into 2D spatial and 1D temporal steps. TSN~\cite{wang2016temporal} samples a few snippets sparsely across a long video and supervises at video level, which scales long-range temporal modeling cheaply, and Transformers~\cite{girdhar2019video,yang2022recurring} model global spatio-temporal dependencies among sampled tokens beyond the local reach of convolution.

\paragraph{Temporal Action Localization.}
Early TAL ran a classifier over dense, fixed-size sliding windows, which was both slow and blind to the wide range of action durations. Structured segment networks~\cite{zhao2017temporal} instead model a proposal's starting, ongoing, and ending stages and score its completeness, which refines boundary regression. Boundary-sensitive methods~\cite{lin2018bsn} go further and predict start and end probabilities at every snippet, then pair high-confidence boundary points into proposals; this flexible boundary generation is the component on which later detection systems rest.

Because precise boundary annotation is expensive, weakly supervised TAL (WSTAL)~\cite{shou2018autoloc,zhang2021cola,paul2018w} learns from video-level labels alone. The task becomes multiple-instance learning: a video is positive if at least one segment contains the action, class activation sequences are derived from video features, and foreground segments are thresholded into proposals. Trained only with a classification loss, such models fire on the most discriminative moments rather than the whole action, so proposals are incomplete and high confidence does not imply accurate boundaries. Point-level supervision (P-TAL)~\cite{xia2024realigning,lee2021learning,fu2022compact} adds one labeled frame per instance, which tends to fall near the temporally most salient, roughly central moment, and uses it to realign confidence and infer complete boundaries from that single anchor.


\paragraph{Temporal Action Detection.}
TAD methods differ in how they generate and score temporal proposals.
Anchor-based detection~\cite{gao2017turn} transplants Faster R-CNN and SSD to the time axis: a dense set of multi-scale temporal anchors is classified and regressed, first in two stages (S-CNN~\cite{shou2016temporal}) and then in one (SSAD~\cite{lin2017single}); fixed anchors produce many redundant proposals and limit boundary precision, which anchor-free variants such as AFSD~\cite{lin2021learning} address by regressing boundaries from a temporal center representation.
Boundary-based detection~\cite{xiong2017pursuit,lin2018bsn} is bottom-up: start and end probabilities are predicted over the whole timeline and boundary matching scores every candidate pairing, which gives flexible, precise boundaries; TSA-Net~\cite{gong2019scale} adds multi-dilation temporal convolution for multi-scale context and BGCN~\cite{chen2021boundary} graph convolution over co-occurring proposals.
Query-based detection~\cite{wu2021towards,shi2022react,kim2023self,zhu2024dual} follows DETR: a small set of learned queries attends to global video features in a Transformer encoder--decoder and decodes the action set in parallel, removing anchors, NMS, and hand-designed matching at the cost of heavier computation.
Across the three, anchors give dense parallel scoring but coarse boundaries, boundary methods give precise boundaries but complex matching, and query methods give the simplest end-to-end pipeline; all three keep moving toward longer-range context, continuous temporal representations, and less frame-level annotation.

\subsubsection{Datasets and Metrics}
Table~\ref{tab:action_detection_dataset} lists the benchmarks by role. UCF101~\cite{soomro2012ucf101} and HMDB51~\cite{kuehne2011hmdb} are trimmed-clip classification sets; AVA~\cite{gu2018ava} adds spatio-temporal boxes on movie footage. THUMOS'14~\cite{jiang2014thumos}, whose 20 classes carry dense temporal boundaries, and ActivityNet~\cite{caba2015activitynet}, with 200 classes over 20K untrimmed videos, are the standard TAL and TAD benchmarks, extended in scale and granularity by HACS Segment~\cite{zhao2019hacs} and FineAction~\cite{liu2022fineaction}. MultiTHUMOS~\cite{yeung2018every}, Charades~\cite{sigurdsson2016hollywood}, and TVSeries~\cite{de2016online} supply dense, multi-label, overlapping actions in sports, homes, and television; Breakfast~\cite{kuehne2014language}, 50Salads~\cite{stein2013combining}, MPII Cooking 2~\cite{rohrbach2016recognizing}, and COIN~\cite{tang2019coin} cover procedural cooking and instructional activity for segmentation and step detection; and MEVA~\cite{corona2021meva} (thousands of hours of surveillance) and MoVi~\cite{ghorbani2021movi} (motion-captured actors) sit at the large-scale and controlled extremes.

\begin{table}[!t]
\centering
\caption{Benchmark datasets for action detection.}
\label{tab:action_detection_dataset}
\fontsize{9}{11}\selectfont
\setlength{\tabcolsep}{5.5pt}
\begin{tabular}{lccccc}
\toprule
\multirow{2}{*}{\textbf{Dataset}} & \textbf{Action} & \multirow{2}{*}{\textbf{Videos}} & \textbf{Duration} & \textbf{Multi-label} & \multirow{2}{*}{\textbf{Domain}} \\
& \textbf{Categories} & & \textbf{(hours)} & \textbf{(Per-Frame)} & \\
\midrule
UCF101~\cite{soomro2012ucf101}           & 101  & 13,320 & 27 & --  & Sports, Daily Activities\\
HMDB51~\cite{kuehne2011hmdb}           & 51   & 6,849  & -- & --  & Movies \\
AVA~\cite{gu2018ava}              & 80   & 430    & 108 & Yes & Movies \\
THUMOS'14~\cite{jiang2014thumos}        & 101 (20 w/ boundaries)  & 18,394 & --  & No & Sports \\
MultiTHUMOS~\cite{yeung2018every}  & 65                         & 413             & --                       & Yes         & Sports                 \\
Breakfast~\cite{kuehne2014language}    & 48                         & 1,712           & 77                    & No          & Cooking                \\
50Salads~\cite{stein2013combining}     & 17                         & 50              & 5.3                       & No          & Cooking                \\
MPII cooking 2~\cite{rohrbach2016recognizing} & 59                        & 273             & 27          & No          & Cooking                \\
TVSeries~\cite{de2016online}     & 30                         & 27              & --                   & Yes         & TV Series              \\
HACS Segment~\cite{zhao2019hacs} & 200                        & 50,000             & --                       & No          & Daily Activities       \\
ActivityNet-1.3~\cite{caba2015activitynet}  & 200  & 19,994 & 849   & No & Daily Activities \\
Charades~\cite{sigurdsson2016hollywood}         & 157  & 9,848  & --    & Yes         & Daily Activities \\
MoVi~\cite{ghorbani2021movi}             & 20   & --     & 9    & -- & Daily Activities \\
MEVA~\cite{corona2021meva}             & 37   & --     & 9,300  & -- & Surveillance \\
COIN~\cite{tang2019coin}         & 180 tasks                  & 11,827          & 476                       & No          & Daily Activities       \\
FineAction~\cite{liu2022fineaction}      & 106  & 16,732 & --    & Yes         & Daily Activities \\
\bottomrule
\end{tabular}
\end{table}

Each setting has its own metric. Action classification reports Top-1 accuracy, $\mathrm{Acc}_{\mathrm{Top}\text{-}1} = \frac{1}{N} \sum_{i=1}^{N} \mathbb{I}(\hat{y}_i = y_i)$, with Top-$k$ on large label sets; temporal boundaries are given, so localization is not assessed. Localization quality is measured by the temporal intersection over union of predicted and ground-truth intervals,
\begin{equation}
    \mathrm{tIoU}(I_p, I_g) =
    \frac{|I_p \cap I_g|}{|I_p \cup I_g|},
\end{equation}
a prediction counting as correct above a threshold $\tau$; TAL reports average precision at one or several thresholds. TAD scores detections $(I_p, \hat{y}, s)$ with confidence $s$ by mean average precision over classes and thresholds,
\begin{equation}
    \mathrm{mAP} = \frac{1}{|\mathcal{C}|}
    \sum_{c \in \mathcal{C}}
    \frac{1}{|\{\tau_m\}|}
    \sum_{\tau_m} \mathrm{AP}_c(\tau_m),
\end{equation}
which penalizes misclassification, loose boundaries, and duplicates together; THUMOS'14 uses $\tau \in \{0.3,\dots,0.7\}$ and ActivityNet $\{0.5,\dots,0.95\}$, and proposal methods also report average recall at a fixed proposal budget. Weakly supervised methods are often evaluated on ActivityNet-1.2 rather than 1.3, so their numbers are not interchangeable with fully supervised ones.

\begin{table*}[!t]
    \centering
    \caption{Representative methods under the two shared protocols. \textit{Top}: action classification, Top-1 accuracy (\%) on UCF101 (split 1, two-stream RGB+flow where applicable, with ImageNet or Kinetics pretraining as in each paper). \textit{Bottom}: temporal action detection on THUMOS'14, mAP (\%) at tIoU $=0.5$ and averaged over tIoU $\in\{0.3,\dots,0.7\}$, as reported by the cited papers and the comparisons in AFSD~\cite{lin2021learning}, CoLA~\cite{zhang2021cola}, and DualDETR~\cite{zhu2024dual}; ``weak'' = video-level labels only. ``--'' = not reported.}
    \label{tab:ad_performance}
    \fontsize{9}{11}\selectfont
    \setlength{\tabcolsep}{2.0mm}
    \begin{tabular}{L{0.22\textwidth} L{0.30\textwidth} c C{0.14\textwidth} C{0.16\textwidth}}
    \toprule
    \textbf{Method} & \textbf{Family} & \textbf{Year} & \textbf{UCF101 Top-1}$\uparrow$ & \\
    \midrule
    Two-Stream CNN~\cite{simonyan2014two} & RGB + optical-flow streams & 2014 & 88.0 & \\
    TSN~\cite{wang2016temporal}           & Sparse segment sampling    & 2016 & 94.2 & \\
    R(2+1)D~\cite{tran2018closer}         & Factorized 3D convolution  & 2018 & 97.3 & \\
    I3D~\cite{carreira2017quo}            & Inflated 3D convolution    & 2017 & \textbf{98.0} & \\
    \midrule
    \hiderowcolors \textbf{Method} & \textbf{Family} & \textbf{Year} & \textbf{mAP@0.5}$\uparrow$ & \textbf{Avg.\ mAP@[0.3:0.7]}$\uparrow$ \\
    \midrule
    S-CNN~\cite{shou2016temporal}     & Anchor-based, two-stage          & 2016 & 19.0 & -- \\
    SSN~\cite{zhao2017temporal}       & Segment-structure proposals      & 2017 & 29.8 & -- \\
    BSN~\cite{lin2018bsn}             & Boundary-sensitive proposals     & 2018 & 36.9 & 36.8 \\
    AutoLoc~\cite{shou2018autoloc}    & Weak, outer--inner contrast      & 2018 & 21.2 & -- \\
    CoLA~\cite{zhang2021cola}         & Weak, snippet contrastive        & 2021 & 32.2 & -- \\
    AFSD~\cite{lin2021learning}       & Anchor-free, salient boundary    & 2021 & 55.5 & 52.0 \\
    DualDETR~\cite{zhu2024dual}       & Query-based (dual DETR)          & 2024 & \textbf{70.4} & \textbf{66.8} \\
    \bottomrule
    \end{tabular}
\end{table*}

Table~\ref{tab:ad_performance} shows how far the two settings have diverged. Trimmed-clip classification on UCF101 was effectively solved once Kinetics-pretrained 3D networks arrived, with I3D and R(2+1)D within a point of each other near 98\%. Detection on THUMOS'14 tells a different story: mAP@0.5 moved from 19.0 (S-CNN) to 36.9 (BSN) over the anchor-to-boundary transition, jumped to 55.5 with anchor-free salient-boundary features (AFSD), and reached 70.4 with query-based decoding (DualDETR), while the best weakly supervised method (CoLA, 32.2) matches fully supervised accuracy of only a few years earlier, which quantifies what video-level labels still cost.


\subsubsection{Challenges}
Fully supervised TAD depends on densely annotated boundaries, which scale poorly to long, loosely structured real-world video; weakly and semi-supervised methods relax this with video-level labels or sparse temporal cues~\cite{zhao2025pose,ghoddoosian2023weakly}, at the accuracy cost the table above makes visible.
Actions are also still detected in isolation. Coupling TAD with structured temporal representations such as scene graphs or event schemas~\cite{wu2025universal,lu2025benchmarking} would let a system infer boundaries, temporal relations, and higher-level activity structure jointly, which matters most in instructional video, driving, and human--robot interaction.
Finally, benchmarks and detectors remain visual and closed-set, whereas real action understanding draws on audio, language, gaze, and context and must handle actions outside a fixed vocabulary~\cite{zhou2024learning,bao2025exploiting,liu2026synergistic}; multimodal temporal detection, where audio or language sharpens boundaries and disambiguates semantics, and open-vocabulary formulations that loosen the classification-centric setup are the directions large multimodal video corpora now make feasible.





\subsection{Mapping and SLAM}
\label{sec:mapping_slam}
Mapping builds a spatial representation of an unknown environment from streaming sensor observations, together with an estimate of the agent's trajectory through it.
Doing so from on-board sensors requires the agent to localize itself inside the partial map it is still building: the chicken-and-egg problem known as \emph{Simultaneous Localization and Mapping (SLAM)}. SLAM resolves it by estimating trajectory and map jointly from the same sensor stream, and observations of previously visited places supply the external references that bound the drift of purely proprioceptive odometry (IMU, wheel encoders). The result is a globally consistent spatial model that supports localization, path planning, obstacle avoidance, and exploration.
Figure~\ref{fig:p3_s6_3_5_slam_modules} depicts the pipeline from sensors through front-end odometry and back-end optimization to the final map.



\begin{figure}[!ht]
    \centering
    \includegraphics[width=\linewidth]{figures/part3/section6/slam.pdf}
    \caption{Example SLAM outputs: real-time visual SLAM (left), a keyframe trajectory over a sparse feature map with tracked image features in the inset, and LiDAR SLAM (right), a dense point-cloud map.}
    \label{fig:p3_s6_3_5_slam_overview}
\end{figure}

\subsubsection{Task Definition}
The input is a stream of exteroceptive observations (camera images, RGB-D frames, LiDAR, or radar), optionally fused with proprioceptive readings (IMU, wheel odometry, velocity and steering commands) and priors such as a coarse map, an initial pose, or GPS.
The output is a map plus the estimated trajectory. Maps lie on an abstraction--detail spectrum: a \emph{topological map} is a graph of distinctive places and their traversable connections, with no metric layout; a \emph{sparse metric map} stores registered keypoints or landmarks and the camera-pose graph that observes them, enough for accurate localization but not for geometry; and a \emph{dense metric map} reconstructs a complete surface or volume (point cloud, mesh, TSDF, or occupancy grid) for obstacle avoidance and 3D reconstruction. Two dense representations have become prominent in the past three years: \emph{neural implicit maps} that store geometry and appearance in network weights (an MLP in iMAP~\cite{Sucar2021iMAP}, a voxel-grid encoder in NICE-SLAM~\cite{Zhu2022NICESLAM}) and \emph{3D Gaussian Splatting (3DGS) maps} that represent the scene as explicit, differentiable 3D Gaussians~\cite{Keetha2024SplaTAM,Yan2024GSSLAM,Matsuki2024MonoGS}.


\subsubsection{Methodology}

SLAM systems are classified by the exteroceptive sensor and by the density of the map they build.
\textbf{LiDAR SLAM} measures range directly with a stable error model; KISS-ICP~\cite{Vizzo2023KISSICP} and KISS-SLAM~\cite{Guadagnino2025KISSSLAM} show that a minimal pipeline delivers robust real-time 3D LiDAR odometry and mapping at city scale. \textbf{Visual SLAM} uses monocular, stereo, or RGB-D cameras, exploiting texture for place recognition, with PTAM~\cite{KleinMurray2007PTAM} and the ORB-SLAM family~\cite{mur2015orb,Campos2021ORBSLAM3} as the foundational systems. \textbf{Multi-sensor SLAM} tightly fuses cameras, IMUs, and/or LiDAR (VINS-Mono, OKVIS, FAST-LIO2), and \textbf{event-camera SLAM} uses asynchronous per-pixel brightness changes for high-speed and high-dynamic-range operation~\cite{Huang2026EventSLAMSurvey}.
By map density, sparse keypoint-and-pose-graph systems (ORB-SLAM) sit next to direct methods that minimize photometric error over semi-dense (LSD-SLAM~\cite{engel2014lsd}) or sparse (DSO~\cite{engel2017direct}) image samples, dense geometric systems that fuse point clouds, meshes, or TSDFs (KinectFusion, ElasticFusion), neural implicit systems (iMAP~\cite{Sucar2021iMAP}, NICE-SLAM~\cite{Zhu2022NICESLAM}, Co-SLAM, NICER-SLAM), and 3DGS systems that render and optimize Gaussians jointly with poses (SplaTAM~\cite{Keetha2024SplaTAM}, MonoGS~\cite{Matsuki2024MonoGS}, SplatSLAM).
Read historically, the field moved from feature-based geometry (PTAM, ORB-SLAM) to direct methods (LSD-SLAM, DSO), to tightly coupled multi-sensor systems (OKVIS, VINS-Mono, ORB-SLAM3), to learning-augmented pipelines (DROID-SLAM~\cite{teed2021droid}, SuperPoint and SuperGlue matching~\cite{detone2018superpoint,sarlin2020superglue}), to neural implicit and then 3DGS maps, and most recently to 3D geometry foundation models (DUSt3R~\cite{wang2024dust3r}, MASt3R~\cite{mast3r}, MASt3R-SLAM~\cite{Murai2025MASt3RSLAM}).

\begin{figure}[!ht]
    \centering
    \includegraphics[width=\linewidth]{figures/part3/section6/slam_method.pdf}
    \caption{SLAM module data flow. Sensor observations feed odometry and place recognition; odometry constraints and verified revisit constraints enter back-end optimization. Mapping registers raw observations using optimized poses and landmarks, producing a map and trajectory. The blue edge in the pose graph is a loop-closure constraint between previously visited poses, not an execution-loop arrow. Map representations are alternatives.}
    \label{fig:p3_s6_3_5_slam_modules}
\end{figure}

Whatever the sensor or map, a SLAM system runs the four modules of Figure~\ref{fig:p3_s6_3_5_slam_modules} in a loop.
The \textbf{front end} is visual (or LiDAR) odometry, the incremental motion estimate between consecutive readings treated as a stand-alone task in \S\ref{sec:visual_odometry}; because each pose is estimated relative to the last, its errors accumulate as drift.
\textbf{Back-end optimization} corrects that drift globally. Modern back ends pose the problem as factor-graph optimization: robot poses and landmarks are nodes, and measurements between them (visual reprojection, IMU pre-integration, GPS, loop-closure relative poses) are factors; solving for all nodes at once yields a consistent trajectory and map, factor graphs absorb heterogeneous sensors naturally, and incremental solvers such as iSAM2 update the estimate in real time, with GTSAM, g\textsuperscript{2}o, and OKVIS as representative libraries.
\textbf{Loop closing} supplies the back end with its strongest constraints by recognizing that the agent has returned to a previously visited place, comparing current observations against the history and, on a match, adding a constraint between two temporally distant poses that pulls the accumulated drift of the intervening trajectory back into alignment.
\textbf{Mapping} finally places the sensor data in the global frame using the optimized trajectory, in whatever representation the application needs; for neural implicit and 3DGS maps it also minimizes a rendering loss (photometric or perceptual such as LPIPS) between observed images and views rendered from the current map, optimizing geometry, appearance, and poses jointly~\cite{Keetha2024SplaTAM}.

\begin{table*}[!th]
\centering
\caption{Mapping and SLAM benchmarks. Scale counts sequences unless noted.}
\label{tab:mapping_slam_dataset}
\fontsize{9}{11}\selectfont
\setlength{\tabcolsep}{1.2mm}
\begin{tabular}{L{0.17\textwidth} C{0.10\textwidth} C{0.16\textwidth} C{0.11\textwidth} C{0.13\textwidth} C{0.07\textwidth} C{0.13\textwidth}}
\toprule
\textbf{Dataset} & \textbf{Task} & \textbf{Modality} & \textbf{Agent} & \textbf{Scene} & \textbf{Scale} & \textbf{Ground Truth} \\
\midrule
COLD \cite{PronobisC09} & VSLAM & mono + omni + odom & wheeled & lab & 76 & laser \\
TUM RGB-D \cite{Sturm2012TUMRGBD} & VSLAM & RGB-D & handheld, wheeled & office, work shed & 39 & MoCap \\
KITTI \cite{KITTI-GeigerLU12} & VSLAM, VO & stereo + laser & car & suburban, highway & 22 & GPS/OXTS \\
ICL-NUIM \cite{HandaWMD14} & VSLAM, VO & RGB-D & synthetic & home, office & 8  & simulator \\
EuRoC \cite{Burri2016EuRoC} & MSLAM & stereo + IMU & drone & work shed & 11  & Vicon MoCap/laser \\
SceneNet \cite{McCormacHLD17} & VSLAM & RGB-D & synthetic & home & 15K  & simulator \\
TUM VI \cite{Schubert2018TUMVI} & MSLAM, VO & stereo + IMU & handheld & campus buildings & 28  & MoCap \\
Bonn RGB-D \cite{PalazzoloBLGS19} & VSLAM & RGB-D & handheld & lab & 24  & MoCap \\
InteriorNet \cite{LiSMCTYHTL18} & MSLAM & RGB-D + IMU & synthetic & home & 15K  & simulator \\
OpenLORIS-Scene \cite{ShiLZTTLZSQSGW020} & MSLAM & RGB-D stereo + IMU + odom & wheeled & indoor & 22  & MoCap/laser \\
Replica \cite{straub2019replica}        & dense       & RGB-D     & synthetic      & indoor       & 18    & simulator \\
ScanNet++ \cite{yeshwanth2023scannet++}  & dense, sem  & RGB-D     & handheld       & indoor       & 460   & laser \\
TartanAir \cite{TartanAir}      & MSLAM, VO   & stereo + IMU & synthetic  & diverse      & 1M+ frames   & simulator \\
ETH3D SLAM \cite{ETH3DSLAM}             & VSLAM, VO   & stereo + global shutter & handheld & indoor/outdoor &  -- & LiDAR \\
DSEC \cite{gehrig2021dsec}              & event SLAM  & stereo events & car       & driving      & 53    & LiDAR + GPS \\
CU-Multi \cite{albin2025cu}            & collab.\ SLAM & RGB-D + LiDAR + IMU & wheeled & outdoor & 8 runs & RTK-GNSS \\
\bottomrule
\end{tabular}
\end{table*}

\subsubsection{Datasets and Metrics}
Table~\ref{tab:mapping_slam_dataset} groups the benchmarks by what they stress. The classical trajectory sets are TUM RGB-D~\cite{Sturm2012TUMRGBD} (handheld indoor RGB-D with motion-capture ground truth), KITTI~\cite{KITTI-GeigerLU12} (stereo and LiDAR driving), and the synthetic ICL-NUIM~\cite{HandaWMD14} with exact ground truth, while COLD~\cite{PronobisC09} targets topological localization under changing conditions. Visual-inertial systems are tested on EuRoC~\cite{Burri2016EuRoC} (micro-aerial vehicle), TUM VI~\cite{Schubert2018TUMVI}, and OpenLORIS-Scene~\cite{ShiLZTTLZSQSGW020}, the last built for lifelong SLAM in changing indoor scenes, and Bonn RGB-D~\cite{PalazzoloBLGS19} isolates dynamic objects. Large synthetic corpora (SceneNet~\cite{McCormacHLD17}, InteriorNet~\cite{LiSMCTYHTL18}, Replica~\cite{straub2019replica}, TartanAir~\cite{TartanAir}) supply photorealistic training and controlled evaluation, with Replica the standard testbed for neural implicit and 3DGS SLAM; ScanNet++~\cite{yeshwanth2023scannet++} and ETH3D SLAM~\cite{ETH3DSLAM} add laser-scanned geometry for dense reconstruction quality; and DSEC~\cite{gehrig2021dsec} and CU-Multi~\cite{albin2025cu} extend evaluation to event cameras and multi-robot data association. AirSwarm~\cite{li2025airswarm} is a multi-UAV research platform rather than a benchmark dataset.

Evaluation covers trajectory, map, rendering, and system cost. Trajectory accuracy uses the absolute trajectory error and relative pose error defined in \S\ref{sec:visual_odometry}, after aligning the estimated trajectory to the ground truth by a rigid or, for monocular systems, similarity transform; RMSE is the usual summary, with mean, median, and standard deviation also reported.
Dense maps are scored by \emph{accuracy} (mean distance from reconstructed points to the nearest ground-truth point) and \emph{completeness} (the reverse), summarized by an F-score at a distance threshold, with Chamfer distance common for neural and 3DGS maps. Neural implicit and 3DGS maps add novel-view-synthesis metrics (PSNR, SSIM, LPIPS) on held-out views of the same trajectory. Memory footprint, tracking and mapping frame rates, and end-to-end latency decide embedded deployment~\cite{Guadagnino2025KISSSLAM}.
Mature evaluation tools such as EVO\footnote{\url{https://github.com/MichaelGrupp/evo}} and the RPG Trajectory Evaluation toolbox\footnote{\url{https://github.com/uzh-rpg/rpg_trajectory_evaluation}} implement these protocols.








\subsubsection{Challenges}
Deployment often violates common SLAM assumptions of static scenes, stable illumination, abundant texture, and ample computation. Dynamic objects induce correspondences that drift with their own motion; illumination change, reflective and transparent surfaces, and repetitive texture create perceptual aliasing that confuses data association and place recognition alike; and resource-constrained platforms must run ever heavier front and back ends. Current work answers each directly. Moving features are detected and excluded through semantic segmentation, motion-consistency checks, or learned masks; learned matchers (SuperPoint~\cite{detone2018superpoint}, SuperGlue~\cite{sarlin2020superglue}, LoFTR~\cite{Sun2021LoFTR}) replace hand-crafted descriptors and tolerate larger viewpoint and illumination changes; neural implicit and 3DGS maps are optimized jointly with poses and yield photorealistic reconstructions usable for inspection, telepresence, and policy learning; and the two-view geometry priors of DUSt3R~\cite{wang2024dust3r} and MASt3R~\cite{mast3r} remove the fragile hand-engineered front end, already enabling real-time dense monocular SLAM at the robustness of MASt3R-SLAM~\cite{Murai2025MASt3RSLAM}.

Two broader shifts follow. Lifelong SLAM aims at months of continuous operation by combining relocalization, map update, and place re-identification against changing environments, and is now an evaluation track of its own; collaborative multi-agent SLAM fuses observations across robots for faster coverage and graceful degradation, increasingly together with event cameras and heterogeneous sensors for high-speed, low-light, and high-dynamic-range conditions. What remains open is concrete: 3DGS maps still do not scale to city-level memory and rendering budgets, learned components carry a sim-to-real gap, and incorrect loop closures under strong perceptual aliasing are still the failure that breaks every back end.





% \newpage



\section{Spatio-Temporal Perception}
\label{sec:p3-sec6-4-spatio-temporal-perception}

\FramedInitial{M}{ost} of what precedes has described a single instant. Spatio-temporal perception adds the missing axis, asking not only where entities are but how they move, when events happen, and which identities persist. 
Its three concerns rise from the pixel through the object to the event: scene flow recovers dense, per-point motion between frames; object tracking maintains the identity of an entity as it moves through space and time; and temporal localization fixes when an action or event occurs along a sequence.
Together they give the largely static perception of the preceding sections a temporal thread, the continuity on which motion understanding and long-horizon perception depend.






\subsection{Scene Flow Estimation}
\label{sec:scene_flow}
Scene flow is the dense three-dimensional motion of every visible surface point between two instants: optical flow lifted out of the image plane and joined to the change in depth. The image-based pipelines surveyed here take stereo video, monocular video, or a multi-view sequence and produce per-pixel 3D motion; some refinement modules additionally consume estimated or measured depth. Dynamic scene reconstruction~\cite{park2021nerfies}, motion-aware rendering~\cite{li2021neural}, physical interaction modeling~\cite{Bear2021Physion}, and egocentric world understanding~\cite{grauman2022ego4d} all consume it.

\subsubsection{Task Definition}
Given consecutive frames $I_t$ and $I_{t+1}$ (optionally with a stereo partner or several views, and with intrinsics and relative poses when known), scene flow estimation returns a field $\mathbf{v}(\mathbf{x}) \in \mathbb{R}^3$ that displaces the surface point imaged at pixel $\mathbf{x}$ in frame $t$ to its position at $t+1$, in camera or world coordinates. In the stereo setting this is equivalently a triple $(d_t, d_{t+1}, \mathbf{f})$ of disparity in both frames plus optical flow, which is how the KITTI benchmark represents and scores it. In image-only pipelines, structure and motion are inferred from appearance, parallax, and temporal correspondence. Geometry-aware refinement modules may instead take depth estimated upstream or measured by RGB-D sensors; point-cloud-only methods are outside the scope here.

\subsubsection{Methodology}

\begin{figure}[!ht]
    \centering
    \includegraphics[width=0.99\linewidth]{figures/part3/section6/scene_flow_estimation.pdf}
    \caption{Image- and video-based scene flow. Top row: observation settings and supervision signals. Bottom row: (1) variational / piecewise-rigid energy minimization, (2) stereo and temporal CNNs, (3) recurrent SE(3) refinement with image and depth pairs, (4) monocular estimation using photometric and geometric constraints, and (5) related motion-field approaches to dynamic view synthesis. The columns illustrate distinct mechanisms rather than successive stages of one pipeline.}
    \label{fig:p3_s6_4_3_method_overview}
\end{figure}

Methods differ along two axes: what they observe (stereo, the classical and still dominant setting; monocular video, the self-supervised frontier; multi-view rigs) and what supervises them (synthetic ground truth, photometric and geometric consistency, or view synthesis). Figure~\ref{fig:p3_s6_4_3_method_overview} groups them into five families along these axes.

\paragraph{Variational and energy-based methods.}
\citet{vedula1999three} introduced the term and the problem: recover 3D structure and 3D motion jointly from a calibrated camera array. \citet{huguet2007variational} formulated it using stereo image pairs at two consecutive times, solving disparity and flow jointly in one variational energy; \citet{wedel2011stereoscopic} decoupled the two into a stereo stage and a flow stage, which made real-time estimation feasible; and \citet{basha2013multi} returned to multi-view input under a view-centered energy. The strongest classical formulation is the piecewise-rigid scene model of \citet{vogel20153d} (PRSF/PRSM), which segments the scene into rigidly moving planar patches and was the leading entry on KITTI Scene Flow at the benchmark's release.

\paragraph{End-to-end stereo and temporal CNNs.}
\citet{mayer2016large} recast disparity and flow as supervised regression (DispNet, FlowNet) and released FlyingThings3D, Monkaa, and Driving, the synthetic corpora that made this training regime possible. SENSE~\cite{jiang2019sense} shares one encoder across disparity, optical flow, occlusion, and semantic heads; DRISF~\cite{ma2019deep} keeps the piecewise-rigid prior of Vogel \etal\ but implements it as an instance-level decoder over learned features. DRISF surpassed both listed variational baselines on KITTI Scene Flow, while SENSE improved on OSF but not PRSM (Table~\ref{tab:scene_flow_performance}), shifting the bottleneck from optimization to supervised data.

\paragraph{Recurrent geometric refinement.}
RAFT~\cite{teed2020raft} introduced all-pairs correlation volumes with GRU-based iterative updates for 2D optical flow; RAFT-3D~\cite{teed2021raft} uses image and depth pairs to iteratively refine a per-pixel SE(3) field with rigid-motion embeddings. Depth can be obtained from stereo or RGB-D input; it is an input to the refinement rather than necessarily inferred jointly from appearance. Its scene-flow results are included in Table~\ref{tab:scene_flow_performance}.

\paragraph{Monocular and self-supervised methods.}
Mono-SF~\cite{brickwedde2019mono} was the first monocular entry on KITTI Scene Flow, resolving scale with a learned single-view depth prior. Self-Mono-SF~\cite{hur2020self} jointly predicts disparity and scene flow with a single decoder, using unlabeled stereo video during training and monocular image pairs at test time; stereo geometry provides the scale cue, while photometric and 3D consistency losses provide self-supervision. A separate test-time optimization approach for LiDAR point clouds (Neural Scene Flow Prior~\cite{li2021neuralsfp}) lies outside the image-based scope of this subsection.

\paragraph{Implicit-field and view-synthesis-coupled methods.}
A parallel line never targets scene flow as output but obtains a dense 4D motion field as a byproduct of dynamic view synthesis. Neural Scene Flow Fields~\cite{li2021neural} trains a dynamic NeRF together with an explicit scene-flow MLP; D-NeRF~\cite{pumarola2021d} represents motion as a deformation field into a canonical frame; Nerfies~\cite{park2021nerfies} extends this to casually captured monocular video; and 4D Gaussian Splatting~\cite{wu20244d} attaches time-parameterized motion to explicit Gaussian primitives. These methods are evaluated by rendering quality rather than by motion error, a split returned to under Challenges below.

\subsubsection{Datasets and Metrics}

Table~\ref{tab:scene_flow_datasets} lists the image-based benchmarks. Synthetic sets provide dense, exact ground truth at scale (FlyingThings3D as the standard pretraining corpus; Sintel and Spring for film-grade motion, with Spring at HD resolution and super-resolved ground truth; TartanAir for environment diversity), while KITTI Scene Flow 2015~\cite{menze2015object}, built on the KITTI suite~\cite{KITTI-GeigerLU12}, is the only real-world benchmark and carries sparse LiDAR-derived labels on 200 training and 200 test scenes. LiDAR-based scene-flow splits of nuScenes~\cite{caesar2020nuscenes}, Waymo Open~\cite{WaymoOpen-SunKDCPTGZCCVHN20}, and Argoverse~2~\cite{wilson2023argoverse}, and the event-camera flow of DSEC~\cite{gehrig2021dsec}, serve the point-cloud and event communities and are not image-based benchmarks in the sense used here.

\begin{table}[t]
\centering
\fontsize{9}{11}\selectfont
\setlength{\tabcolsep}{1.0mm}
\caption{Image-based scene-flow benchmarks. Counts are indicative of the commonly used releases.}
\label{tab:scene_flow_datasets}
\begin{tabular}{L{0.26\textwidth} c C{0.10\textwidth} C{0.11\textwidth} C{0.14\textwidth} L{0.24\textwidth}}
\toprule
\textbf{Dataset} & \textbf{Year} & \textbf{Type} & \textbf{Input} & \textbf{Scale} & \textbf{Ground truth} \\
\midrule
KITTI Scene Flow 2015~\cite{menze2015object} & 2015 & Real driving & Stereo video & 200 train / 200 test scenes & Sparse disparity ($t$, $t{+}1$) + flow from LiDAR and CAD fits \\
FlyingThings3D / Monkaa / Driving~\cite{mayer2016large} & 2016 & Synthetic & Stereo video & $\sim$25K / 8.6K / 4.4K frames & Dense disparity, flow, disparity change \\
MPI Sintel~\cite{sintel} & 2012 & Synthetic film & Stereo video & 1,041 train frames, 23 seq. & Dense flow, depth, disparity \\
Spring~\cite{Spring} & 2023 & Synthetic film & Stereo video & $\sim$6K frames, 47 seq., 1920$\times$1080 & Dense stereo, flow, scene flow (super-resolved) \\
TartanAir~\cite{TartanAir} & 2020 & Synthetic & Stereo video & $>$1M frames, 30 environments & Dense depth, flow, pose; indoor and outdoor, weather \\
\bottomrule
\end{tabular}
\end{table}

KITTI scores scene flow in image space. A pixel's disparity in frame $t$ (D1), its disparity in frame $t{+}1$ (D2), and its optical flow (Fl) are each counted as outliers when the error exceeds both 3\,px and 5\% of the ground-truth value, and a scene-flow outlier (SF) is any pixel that fails at least one of the three; the benchmark reports outlier percentages over all pixels and separately over background and foreground. Point-level evaluation on synthetic data instead uses the 3D end-point error
\begin{equation}
    \mathrm{EPE}_{3D} = \frac{1}{N} \sum_{i=1}^{N} \| \hat{\mathbf{v}}_i - \mathbf{v}_i \|_2,
\end{equation}
with $\hat{\mathbf{v}}_i$ and $\mathbf{v}_i$ the predicted and true motion of point $i$, together with accuracy ratios at fixed thresholds (the fraction of points with error below 0.05\,m or 5\%, and below 0.1\,m or 10\%) and an outlier ratio (error above 0.3\,m or 10\%), conventions inherited from point-cloud scene flow.

Table~\ref{tab:scene_flow_performance} collects KITTI Scene Flow 2015 test results for the explicit-output families. RAFT-3D reduces SF-all error from 8.97\% (PRSM) to 5.77\%; the table does not separate foreground and background results. The monocular rows are a different regime: Mono-SF, with a learned depth prior, sits at 23\% SF-all, and Self-Mono-SF without any ground truth at 49.5\% (33.9\% after fine-tuning), the price of giving up stereo. The implicit-field family has no row because it reports rendering quality (PSNR, LPIPS) on its own scene sets rather than motion error.

\begin{table}[!t]
    \centering
    \fontsize{9}{11}\selectfont
    \setlength{\tabcolsep}{2.2mm}
    \caption{KITTI Scene Flow 2015 test set, outlier percentages (lower is better) over all pixels. Stereo rows from the RAFT-3D comparison~\cite{teed2021raft}; monocular rows from the Mono-SF~\cite{brickwedde2019mono} and Self-Mono-SF~\cite{hur2020self} papers (``ft'' = fine-tuned on KITTI). D1 / D2: disparity in frames $t$ / $t{+}1$; Fl: optical flow; SF: scene flow.}
    \label{tab:scene_flow_performance}
    \begin{tabular}{L{0.30\textwidth} L{0.26\textwidth} C{0.07\textwidth} C{0.07\textwidth} C{0.07\textwidth} C{0.07\textwidth}}
    \toprule
    \textbf{Method} & \textbf{Family} & \textbf{D1-all} & \textbf{D2-all} & \textbf{Fl-all} & \textbf{SF-all} \\
    \midrule
    \rowcolor{lightgrey} \multicolumn{6}{l}{$\bullet$ \textit{Stereo}}\\
    OSF~\cite{menze2015object}          & Energy-based, rigid objects   & 5.79 & 7.77 & 7.83 & 10.23 \\
    PRSM~\cite{vogel20153d}             & Energy-based, piecewise rigid & 4.27 & 6.79 & 6.68 & 8.97 \\
    SENSE~\cite{jiang2019sense}         & Shared-encoder CNN            & 2.22 & 5.89 & 7.64 & 9.55 \\
    DRISF~\cite{ma2019deep}             & CNN + rigid instances         & 2.55 & 4.04 & 4.73 & 6.31 \\
    RAFT-3D~\cite{teed2021raft}         & Recurrent SE(3) refinement    & 1.81 & 3.67 & 4.29 & 5.77 \\
    \midrule
    \rowcolor{lightgrey} \multicolumn{6}{l}{$\bullet$ \textit{Monocular}}\\
    Mono-SF~\cite{brickwedde2019mono}   & Learned depth prior           & 16.32 & 19.59 & 12.77 & 23.08 \\
    Self-Mono-SF~\cite{hur2020self}     & Self-supervised               & 34.02 & 36.34 & 23.54 & 49.54 \\
    Self-Mono-SF-ft~\cite{hur2020self}  & Semi-supervised, ft           & 22.16 & 25.24 & 15.91 & 33.88 \\
    \bottomrule
    \end{tabular}
\end{table}

\subsubsection{Challenges}

The supervision landscape is lopsided. Every dense-ground-truth benchmark in Table~\ref{tab:scene_flow_datasets} is synthetic and stereo, so monocular methods are scored only against KITTI's sparse labels, and models pretrained on FlyingThings3D or Sintel still lose accuracy on real driving footage after fine-tuning, most visibly on foreground motion where synthetic random-object trajectories differ from vehicle kinematics. All current benchmarks also emphasize rigid or random-object motion; dense ground truth for non-rigid humans, hands, and animals does not exist at scale, which leaves the family of implicit-field methods that handle such motion best without a quantitative yardstick.

That yardstick problem is the second challenge. Explicit-output methods~\cite{vogel20153d,jiang2019sense,teed2021raft} are scored by KITTI's outlier rates; implicit-field methods~\cite{li2021neural,pumarola2021d,wu20244d} by PSNR and LPIPS of rendered views. No protocol evaluates both on the same input, so the two paradigms are compared only qualitatively, and whether a motion field is good depends on whether scene flow is the end (the KITTI leaderboard) or a means (reconstruction, novel-view synthesis, physical reasoning). The cost profiles diverge as well: RAFT-3D needs several GRU iterations per frame, variational methods minutes per frame, and implicit fields minutes to hours per scene.

Foundation priors are reshaping the monocular line. Depth Anything~\cite{DepthAnything} and DINOv2~\cite{dinov2} supply the per-pixel depth and semantic features that Self-Mono-SF had to learn from scratch, and SAM~2~\cite{ravi2024sam2} segmentation together with the object-pose estimators of Section~\ref{sec:object_pose} makes an object-centric formulation (a per-object SE(3) trajectory plus residual non-rigid motion) practical, a bridge to the tracking problem of Section~\ref{sec:object_tracking}. As implicit-field scene flow merges into 4D reconstruction, with 4D Gaussian Splatting~\cite{wu20244d} the early example, evaluation is likely to follow the downstream task rather than the KITTI leaderboard.
% TODO: anchor §8.x once §8 reconstruction subsections stabilize








\subsection{Object Tracking}
\label{sec:object_tracking}

Tracking localizes entities in every frame and associates the detections across frames so that each entity keeps one identity through the sequence. Where detection and segmentation describe a single instant, tracking yields trajectories, the representation that autonomous driving~\cite{luo2021exploring,chiu2021probabilistic}, surveillance~\cite{ojha2015imageobjectsurveillance}, AR/VR~\cite{chen2021hightracking}, and robotics~\cite{Deitke2020RoboTHOR} build on.

\begin{figure}[!ht]
    \centering
     \includegraphics[width=0.95\linewidth]{figures/part3/section6/object_tracking.pdf}
    \caption{Object-tracking subtasks: (a) single object tracking (SOT) from an initial box, (b) multiple object tracking (MOT) with identity-consistent boxes, (c) multiple object tracking and segmentation (MOTS) with per-frame masks, (d) 3D multi-object tracking (3D-MOT) of oriented boxes from LiDAR or camera input, (e) pose and landmark tracking, and (f) panoptic tracking with semantic classes for stuff and persistent instance IDs for things. Repeated IDs denote the same entity across frames; arrows indicate increasing time. Scene images and overlaid states are illustrative.}
    \label{fig:tracking_task}
\end{figure}

\subsubsection{Task Definition}
Given a frame sequence $\{I_t\}_{t=1}^{T}$ (a video, or RGB-D, LiDAR, or multi-view streams), a tracker returns a set of trajectories $\{\mathcal{T}_k\}$, each $\mathcal{T}_k = \{(t, b_t^k)\}$ a time-indexed series of states $b_t^k$ with one identity $k$ held fixed over the sequence. The state is a bounding box in the basic setting and a segmentation mask, a 3D box, or a set of keypoints in the variants.

The variants, illustrated in Figure~\ref{fig:tracking_task}, differ in how the targets are specified and what state is tracked. Single object tracking (SOT)~\cite{fan2019lasot,muller2018trackingnet} follows one target given its first-frame box; multiple object tracking (MOT)~\cite{milan2016mot16,dendorfer2020mot20,yu2020bdd100k} tracks every instance of given categories, typically pedestrians and vehicles, under identity constraints; MOTS~\cite{voigtlaender2019mots} and video object segmentation replace boxes with masks; 3D-MOT~\cite{AB3DMOT,yin2021center,caesar2020nuscenes,WaymoOpen-SunKDCPTGZCCVHN20} outputs world-coordinate 3D boxes from LiDAR, RGB-D, or multi-view input; pose tracking~\cite{andriluka2018posetrack,iqbal2017posetrack,samanta2016data} and landmark tracking~\cite{zafeiriou20173d} follow skeletal joints or facial keypoints; and panoptic tracking combines pixel-level semantic classes for stuff with temporally consistent instance identities for things. Category-specific trackers for pedestrians~\cite{stadler2021improving}, vehicles~\cite{van2001real}, or animals~\cite{zhang2023animaltrack} specialize any of these.


\subsubsection{Methodology}
Figure~\ref{fig:task_tracking_method} organizes the methods into four families around image-plane box and mask tracking (SOT, MOT, MOTS, VOS). Pose, landmark, and panoptic tracking reuse the same propagate-or-associate machinery with a different state representation and are covered at the benchmark level in Table~\ref{tab:video_object_tracking_dataset}.

\paragraph{Propagation-based methods.}
Because a target moves little between consecutive frames, its previous state is a strong prior for the current one, and the oldest tracking paradigm exploits this by predicting the next state from the current one and correcting it against the new frame. Correlation filters realize this in the frequency domain: MOSSE~\cite{bolme2010visual} learns a minimum-output-sum-of-squared-error filter from a single frame and updates it online at hundreds of frames per second; KCF~\cite{henriques2014high} and MKCF~\cite{tang2015multi} add kernelization and multiple-kernel learning, respectively; ECO~\cite{DanelljanBKF17} adds factorized convolution operators, a compact generative sample model, and conservative updates. These trackers are fast enough for embedded use but drift under large appearance change and clutter. MDNet~\cite{NamH16} answered with online discriminative learning, a shared CNN with per-sequence domain branches that separate target from background, at a high per-frame cost. Siamese trackers moved the representation learning offline: SiamFC~\cite{BertinettoVHVT16} correlates a template embedding with the search-region embedding in one forward pass, and SiamRPN~\cite{LiYWZH18}, ATOM~\cite{DanelljanBKF19}, and DiMP~\cite{BhatDGT19} add region-proposal regression, IoU-prediction refinement, and a discriminatively learned prediction head. Transformer trackers then replaced correlation with attention: TransT~\cite{ChenYZ0YL21} fuses template and search features with self- and cross-attention, STARK~\cite{0002PF0L21} predicts the box with an encoder--decoder that also carries temporal state, and OSTrack~\cite{YeCMSC22} joins feature extraction and relation modeling in a single stream over concatenated template and search tokens, the current baseline design.

Mask propagation applies the same idea at the pixel level for video object segmentation~\cite{pont20172017,xu2018youtube}, with memory of past frames replacing a single template. XMem~\cite{ChengS22XMem} keeps multi-scale spatiotemporal memories that remain stable over videos of more than 10{,}000 frames; AOT~\cite{YangWY21aot} and DeAOT~\cite{YangY22adeaot} encode object identities as tokens so that many objects propagate in one pass; and SAM~2~\cite{ravi2024sam2} adds a streaming memory to the Segment Anything Model~\cite{kirillov2023sam}, giving promptable tracking that transfers across VOS benchmarks.


\begin{figure}[!t]
    \centering
    \includegraphics[width=0.99\linewidth]{figures/part3/section6/object_tracking_method.pdf}
    \caption{Four method families for object tracking: (i) propagation-based, (ii) detect-then-associate, (iii) joint detection--association, and (iv) universal models with a shared backbone and task heads for SOT, MOT, VOS, and MOTS.}
    \label{fig:task_tracking_method}
\end{figure}

\paragraph{Detect-then-associate methods.}
The MOT standard pipeline detects objects in every frame and then solves an assignment problem between detections and existing tracks. SORT~\cite{bewley2016sort} predicts track positions with a Kalman filter~\cite{welch1995kalman} and matches by IoU with the Hungarian algorithm~\cite{Kuhn10}; DeepSORT~\cite{WojkeBP17} adds an appearance embedding so that tracks survive occlusion. Track-RCNN~\cite{voigtlaender2019mots}, JDE~\cite{WangZLLW20jde}, and FairMOT~\cite{ZhangWWZL21FairMOT} train the re-identification embedding jointly with the detector; association remains a separate step, which keeps them in this family. All of these discard low-confidence detections, and since occluded targets are exactly those with low scores, trajectories fragment where they are hardest to keep. ByteTrack~\cite{ZhangSJYWYLLW22} associates nearly every detection and uses motion similarity to the existing tracks to reject the false positives among the low-score boxes, a change in association alone that raised MOT17 HOTA by several points (Table~\ref{tab:tracking_performance}). In 3D, the same pipeline reappears: AB3DMOT~\cite{AB3DMOT} runs a 3D Kalman filter with Hungarian matching on LiDAR detections, and CenterPoint~\cite{yin2021center} tracks 3D object centers through predicted velocities, with nuScenes~\cite{caesar2020nuscenes} using AMOTA/AMOTP and Waymo~\cite{WaymoOpen-SunKDCPTGZCCVHN20} using MOTA/MOTP.

\paragraph{Joint detection--association methods.}
End-to-end trackers fold association into the detector. CenterTrack~\cite{ZhouKK20} represents objects as center points and regresses their offset to the previous frame, so that association is a local lookup; TrackFormer~\cite{MeinhardtKLF22} introduces tracking-by-attention, in which track queries carried from the previous frame and new object queries are decoded jointly by a Transformer; MOTR~\cite{zeng2022motr} extends DETR~\cite{carion2020end} with track queries that persist across the whole sequence and are supervised with a sequence-level loss, which is what gives it its advantage on DanceTrack~\cite{SunCJYBKL22}, where appearance is uninformative and motion is non-linear. The price is training cost and data: these models need long annotated sequences and remain behind detect-then-associate trackers on MOT17 (Table~\ref{tab:tracking_performance}).

\paragraph{Universal models.}
A single model for SOT, MOT, VOS, and MOTS is the most recent direction. UniTrack~\cite{wang2021different} showed that one self-supervised appearance representation serves propagation and association tasks alike; Unicorn~\cite{yan2022unicorn} trains one network end to end for all four tasks; UTT~\cite{MaSZ0XYY22} unifies SOT and MOT in a Transformer; UniVS~\cite{li2024univs} casts every video-segmentation task as prompt-driven segmentation; and OneTracker~\cite{HongYZLZGJCLCZ24} adapts a foundation tracker to RGB, depth, thermal, and event inputs with light adapters. In practice, SAM~\cite{kirillov2023sam} combined with a memory tracker (XMem, DeAOT) already provides promptable multi-object segmentation and tracking from points, boxes, scribbles, or text, and methods such as UASTrack~\cite{UASTrack} and Un-Track~\cite{WuZRV0PGT24UNTrack} extend unified tracking to RGB paired with depth, thermal, or event inputs. Language enters as a specification signal as well: LaSOT~\cite{fan2019lasot} pairs every sequence with a sentence, TNL2K~\cite{wang2021tnl2k} is built for text-initialized tracking, and LaMOT~\cite{li2025lamot} brings language-guided specification to the multi-object setting.

\begin{table}[!ht]
\centering
\caption{Datasets for object tracking with bounding boxes (Bbox), segmentation masks (Mask), or human poses. ``Frames / Boxes'' gives annotated frames, or annotated boxes where reported (MOT16/17, TrackingNet, YouTube-BB). MOT17 reuses its 14 videos with three public detection sets. For pose datasets, the annotation column gives the setting and the number of body landmarks.}
\label{tab:video_object_tracking_dataset}
\fontsize{9}{11}\selectfont
\setlength{\tabcolsep}{2.2mm}
\begin{tabular}{lcccccc}
\toprule
 \textbf{Name} & \textbf{Videos} & \textbf{Frames / Boxes} & \textbf{Categories} & \textbf{Year} & \textbf{Attributes} & \textbf{Annotations}\\
 \midrule
\rowcolor{lightgrey} \multicolumn{7}{l}{$\bullet$ \textit{Single Object Tracking}} \\
VIVID~\cite{collins2005open} &   9 &  16,274 &  - & 2005 & - & Bbox\\
OTB100~\cite{wu2013online} &   98 &  58,610 &  -  & 2015 & 11 & Bbox\\
ALOV300~\cite{smeulders2013visual} &   314 &  89,364 &   64 & 2013 & - & Bbox\\
TC128~\cite{liang2015encoding} &   128 &  55,652 &  - &  2015  & 11 & Bbox\\
VOT16~\cite{roffo2016visual} &   60 &  21,455 &  - & 2016 & - & Bbox\\
DAVIS-16~\cite{perazzi2016benchmark} & 50 &  3,455 & 4 & 2016 & 15& Mask \\
UAV20L~\cite{benchmark2016uav} &   20 &  58,670 &   -  & 2016 & 12 & Bbox\\
UAV123~\cite{benchmark2016uav} &   123 &  113,476 &   9  & 2016 & 12 & Bbox\\
NUS-PRO~\cite{nus_pro} &   365 &  135,305 &  8 & 2016 & - & Bbox\\
VOT17~\cite{alatan2017visual} &   60 &  21,356 & 24 & 2017 & - & Bbox\\
YouTube-BB~\cite{real2017youtube} &   380,000 &  5.6M &  23  & 2017  & - & Bbox\\
NFS~\cite{kiani2017need} &   100 &  383,000 &   17  & 2017 & 9 & Bbox\\
TrackingNet~\cite{muller2018trackingnet} &   30,132 &  14.2M &   27  & 2018 & 15  & Bbox\\
LaSOT~\cite{fan2019lasot}  &  1,400  & 3.52M &  70 & 2019 & 14 & Bbox\\
TNL2K~\cite{wang2021tnl2k} &  2,000  &  1.24M &  - & 2021 & 17 & Bbox\\
GOT-10k~\cite{huang2021got}   & 10,000 & 1.5M & 563  & 2021 & 6 & Bbox\\
\midrule
\rowcolor{lightgrey} \multicolumn{7}{l}{$\bullet$ \textit{Multiple Object Tracking}} \\
MOT16~\cite{milan2016mot16} &   14 &  292,733 & 1 & 2016  & - & Bbox\\
MOT17~\cite{milan2016mot16} &   14 &  564,228 & 1 &   2017 & -  & Bbox\\
BDD100K~\cite{yu2020bdd100k} &   1,600 &  318K &  8  & 2020  & - & Bbox\\
\midrule
\rowcolor{lightgrey} \multicolumn{7}{l}{$\bullet$ \textit{Video Object Segmentation}} \\
DAVIS-17~\cite{pont20172017} & 150 & 10,459 & - & 2017 &-  & Mask\\
YouTube-VOS~\cite{xu2018youtube} & 4,453 & - &  94 & 2018  & -  & Mask\\
\midrule
\rowcolor{lightgrey} \multicolumn{7}{l}{$\bullet$ \textit{RGB-D Object Tracking}} \\
PTB~\cite{song2013tracking} &  100 &  21,542 & - & 2013 & 5 &  Bbox \\
STC~\cite{xiao2017robust} &  36 & 9,009 & -  & 2017 & 12 &  Bbox \\
CDTB~\cite{lukezic2019cdtb} &  80 &  101,956 & - & 2019 &  13 &  Bbox  \\
DepthTrack~\cite{yan2021depthtrack} & 200 &  294,591 & - & 2021  &  15 &  Bbox \\
ARKitTrack~\cite{zhao2023arkittrack} & 300 & 229.7K & 144 & 2023 & 16 & Mask \\
RGBD1K~\cite{zhu2023rgbd1k} &  1,050&  717,900 & - & 2023 &  15 &  Bbox  \\
\midrule
\rowcolor{lightgrey} \multicolumn{7}{l}{$\bullet$ \textit{Human Pose Tracking}} \\
VideoPose2.0~\cite{sapp2011parsing} & 44 & 1,286 &- &  2011 & - & Single-person (-) \\
Poses in the Wild~\cite{cherian2014mixing} & 30 & 830 & - & 2014 & - & Single-person (7) \\
ICDPose~\cite{samanta2016data} & 60 & 2,700 & - & 2016 &- &  Single-person (-)\\
CMU Panoptic~\cite{joo2017panoptic} & 65 & 297K & - & 2017 & - & Multi-person (-) \\
Multi-Person PoseTrack~\cite{iqbal2017posetrack} & 60 & - & - & 2017 &- &  Multi-person (14)\\
PoseTrack18~\cite{andriluka2018posetrack} & 1,138 & 66,374 & - & 2018 &- &  Multi-person (15)\\
\bottomrule
\end{tabular}
\end{table}

\subsubsection{Datasets and Metrics}

Table~\ref{tab:video_object_tracking_dataset} groups the benchmarks by what is tracked. The SOT line runs from short, category-agnostic clips (OTB100, TC128, the VOT challenges) to large, long-term, category-diverse training sets (TrackingNet, LaSOT, GOT-10k, whose test classes are disjoint from training) and to condition-specific tests for aerial footage (UAV123), high frame rate (NFS), and language initialization (TNL2K). MOT benchmarks are pedestrian-only in MOT16/17 and MOT20 and multi-class driving footage in BDD100K; the 3D-MOT counterparts are the nuScenes and Waymo tracking tasks. DAVIS and YouTube-VOS supply mask-level ground truth, the latter with categories unseen at training time; the RGB-D sets add depth and have grown from tens of sequences (PTB, STC) to over a thousand (RGBD1K); and pose tracking moved from single-person upper-body clips (VideoPose2.0, Poses in the Wild, ICDPose) to multi-person full-body sequences with person identities (PoseTrack18, CMU Panoptic).

Each line has its own metrics. SOT uses the success plot: with $\mathrm{SR}(\tau)$ the fraction of frames whose predicted-box IoU exceeds $\tau$, the headline number is its area under the curve,
\begin{equation}
    \mathrm{AUC} = \int_0^1 \mathrm{SR}(\tau)\, d\tau ,
\end{equation}
reported with normalized precision (the fraction of frames whose center error, divided by the target size, is below a threshold) on LaSOT and TrackingNet, and with average overlap (mean IoU over frames) and $\mathrm{SR}_{0.5}$, $\mathrm{SR}_{0.75}$ on GOT-10k; the VOT challenges use expected average overlap (EAO), which folds accuracy and failure rate into one expectation~\cite{roffo2016visual}. MOT scores detection and identity jointly. MOTA~\cite{bernardin2008evaluating} charges false negatives, false positives, and identity switches against the number of ground-truth objects,
\begin{equation}
    \mathrm{MOTA} = 1 - \frac{\mathrm{FN}+\mathrm{FP}+\mathrm{IDS}}{\mathrm{GT}} ,
\end{equation}
and so is dominated by detection; IDF1 is the F1 score of identity-level true positives, $\mathrm{IDF1} = 2\,\mathrm{IDTP} / (2\,\mathrm{IDTP} + \mathrm{IDFP} + \mathrm{IDFN})$, and emphasizes identity preservation; HOTA~\cite{luiten2021hota}, $\sqrt{\mathrm{DetA}\cdot\mathrm{AssA}}$ averaged over localization thresholds, balances the two and has become the primary MOT metric. VOS reports the Jaccard index $\mathcal{J} = |M \cap G| / |M \cup G|$ of predicted and true masks and the contour F-measure $\mathcal{F}$ from boundary precision and recall~\cite{perazzi2016benchmark}, averaged as $\mathcal{J}\&\mathcal{F}$.

Table~\ref{tab:tracking_performance} reports the methods of the Methodology above on the benchmarks each line treats as standard. Within SOT the step from a discriminative head over correlation features (DiMP) to attention-based fusion (TransT, STARK) is worth about ten points of LaSOT AUC, and the one-stream design of OSTrack adds six more over TransT. On MOT17 the detect-then-associate trackers lead, and ByteTrack's gain over FairMOT comes from association alone; the end-to-end trackers (CenterTrack, TrackFormer, MOTR) trail on MOT17 but MOTR leads ByteTrack by 6.5 HOTA on DanceTrack~\cite{zeng2022motr}, where appearance cues are weak. VOS accuracy on DAVIS-17 has moved from the mid-80s (XMem, DeAOT) to above 90 $\mathcal{J}\&\mathcal{F}$ with SAM~2, which is also the first method trained on a corpus far larger than the benchmarks themselves.

\begin{table}[!t]
    \centering
    \fontsize{9}{11}\selectfont
    \setlength{\tabcolsep}{2.2mm}
    \caption{Representative trackers on the standard benchmarks of each line. SOT: AUC on LaSOT and TrackingNet, AO on GOT-10k (from the OSTrack comparison~\cite{YeCMSC22}). MOT: MOT17 test, private detections (from the MOTR~\cite{zeng2022motr} and ByteTrack~\cite{ZhangSJYWYLLW22} comparisons). VOS: $\mathcal{J}\&\mathcal{F}$ on DAVIS-17 val and $\mathcal{G}$ on YouTube-VOS 2019 val (from the SAM~2 comparison~\cite{ravi2024sam2}).}
    \label{tab:tracking_performance}
    \begin{tabular}{L{0.30\textwidth} L{0.26\textwidth} C{0.10\textwidth} C{0.11\textwidth} C{0.10\textwidth}}
    \toprule
    \rowcolor{lightgrey} \multicolumn{5}{l}{$\bullet$ \textit{Single object tracking}}\\
    \textbf{Method} & \textbf{Design} & \textbf{LaSOT AUC} & \textbf{TrackingNet AUC} & \textbf{GOT-10k AO} \\
    \midrule
    DiMP~\cite{BhatDGT19}         & Discriminative prediction head   & 56.9 & 74.0 & 61.1 \\
    TransT~\cite{ChenYZ0YL21}      & Transformer feature fusion       & 64.9 & 81.4 & 67.1 \\
    STARK~\cite{0002PF0L21}        & Encoder--decoder, temporal state & 67.1 & 82.0 & 68.8 \\
    OSTrack-384~\cite{YeCMSC22}    & One-stream Transformer           & 71.1 & 83.9 & 73.7 \\
    \midrule
    \rowcolor{lightgrey} \multicolumn{5}{l}{$\bullet$ \textit{Multiple object tracking (MOT17 test)}}\\
    \hiderowcolors \textbf{Method} & \textbf{Family} & \textbf{HOTA} & \textbf{MOTA} & \textbf{IDF1} \\
    \midrule
    CenterTrack~\cite{ZhouKK20}            & Joint, center offsets        & 52.2 & 67.8 & 64.7 \\
    TrackFormer~\cite{MeinhardtKLF22}      & Joint, tracking-by-attention & --   & 74.1 & 68.0 \\
    MOTR~\cite{zeng2022motr}               & Joint, persistent queries    & 57.8 & 73.4 & 68.6 \\
    FairMOT~\cite{ZhangWWZL21FairMOT}      & Detect + joint Re-ID         & 59.3 & 73.7 & 72.3 \\
    ByteTrack~\cite{ZhangSJYWYLLW22}       & Detect + low-score association & 63.1 & 80.3 & 77.3 \\
    \midrule
    \rowcolor{lightgrey} \multicolumn{5}{l}{$\bullet$ \textit{Video object segmentation}}\\
    \hiderowcolors \textbf{Method} & \textbf{Design} & \textbf{DAVIS-17 $\mathcal{J}\&\mathcal{F}$} & \textbf{YT-VOS 2019 $\mathcal{G}$} & \\
    \midrule
    XMem~\cite{ChengS22XMem}       & Multi-store memory               & 86.0 & 85.6 & \\
    DeAOT~\cite{YangY22adeaot}     & Identity tokens, decoupled propagation & 86.2 & 86.1 & \\
    SAM~2~\cite{ravi2024sam2}      & Promptable, streaming memory     & 90.7 & 89.3 & \\
    \bottomrule
    \end{tabular}
\end{table}

\subsubsection{Challenges}
Identity under occlusion remains the central failure. Motion models and appearance embeddings bridge short occlusions, but long occlusions and dense crowds still produce identity switches and fragmented tracks~\cite{ZhangSJYWYLLW22,SunCJYBKL22}, and memory-based propagation drifts or forgets over long videos even when each step is accurate~\cite{ChengS22XMem}. HOTA and IDF1 measure the outcome but not robustness to occlusion length or the calibration of track confidence, so the metrics reward the easy part of the problem.

Trackers are also tuned to a few curated benchmarks and transfer poorly to UAV footage, thermal and event imagery, adverse weather, and multi-camera surveillance; foundation-model trackers such as OneTracker~\cite{HongYZLZGJCLCZ24} and the RGB-X methods UASTrack~\cite{UASTrack} and Un-Track~\cite{WuZRV0PGT24UNTrack} address this directly but inherit the biases of their pretrained backbones. Long streams add memory and compute limits on top of drift.

The accuracy of Transformer trackers comes at a cost that keeps them off embedded hardware, while the lightweight correlation-filter line gives up robustness; and the unified models that would amortize that cost across tasks meet negative transfer between objectives with different demands, mask precision against identity association in particular~\cite{yan2022unicorn}. Language and segmentation priors widen what a tracker can be asked to do, at the price of prompt sensitivity that current benchmarks do not measure.







\subsection{Temporal Localization}
\label{sec:temporal_localization}

Temporal localization answers \emph{when} something happens in a long, untrimmed video. The condition can be a category from a closed set, which gives temporal action localization (TAL) as treated in Section~\ref{sec:action_detection}, or a free-form sentence, which gives temporal sentence grounding (TSG): return the segment that best matches ``the person opens the microwave''. This subsection sets the two side by side, with the weight on the text-conditioned problem, which turns localization from exhaustive detection into retrieval and is what video search and interactive navigation~\cite{hendricks2017mcn,gao2017tall,krishna2017densecap} rely on.

\begin{figure}[!ht]
    \centering
    \includegraphics[width=\linewidth]{figures/part3/section6/temporal_localization.pdf}
    \caption{Temporal localization methods. Top row (TAL, category-conditioned): (i) window-based, (ii) two-stage proposal-based, (iii) one-stage anchor-free, and (iv) Transformer set prediction. Bottom row (TSG, text-conditioned): (a) cross-modal matching, (b) 2D temporal-map reasoning with a proposal-graph alternative, and (c) Transformer set prediction. Hatched map cells are invalid start--end pairs; graph nodes represent candidate moments. Coarse-to-fine corpus retrieval is shown separately as a scope extension, not a mandatory stage of set prediction. Scene frames and scores are illustrative.}
    \label{fig:p3_s6_4_2_method_overview}
\end{figure}

\subsubsection{Task Definition}
\textbf{Temporal Action Localization (TAL)} takes an untrimmed video $V=\{x_t\}_{t=1}^{T}$, where $x_t$ may be a visual, audio, or multimodal feature, and detects \emph{all} instances of $K$ predefined action classes, returning $\mathcal{A}=\{a_i\}_{i=1}^{N}$ with $a_i=(t_i^{\mathrm{start}},\, t_i^{\mathrm{end}},\, c_i,\, s_i)$: start and end times, a class label $c_i\in\{1,\ldots,K\}$, and a confidence $s_i$ (the $(t_s, t_e, c)$ tuples of Section~\ref{sec:action_detection}).

\textbf{Temporal Sentence Grounding (TSG)}, also called natural language video localization or video moment retrieval, pairs $V$ with a sentence $q$ and returns a ranked list of moments $\mathcal{M}(q,V)=\{m_j\}_{j=1}^{M}$, $m_j=(\tau_j^{\mathrm{start}},\,\tau_j^{\mathrm{end}},\, r_j)$, with boundaries and a relevance score $r_j$; usually one segment is the answer. TAL is thus an exhaustive detection task over a closed label set, TSG a targeted retrieval task guided by open-ended language.

Both are studied under several settings, which Figure~\ref{fig:p3_s6_4_2_task_settings} lays side by side. Offline methods see the whole video; online methods infer causally over the streamed prefix and are judged on low-latency detection of action starts~\cite{shou2018odas,gao2019startnet}. Supervision ranges from frame-accurate boundaries to video-level labels only (weakly supervised TAL via multiple-instance learning, Section~\ref{sec:action_detection}; weakly supervised TSG by aligning the two modalities during training~\cite{ma2020vlanet}), with self-supervised video pretraining~\cite{tong2022videomae,wang2023videomaev2} and open-vocabulary or zero-shot variants~\cite{liberatori2024t3al} reducing the labeled data needed. TSG additionally varies in scope, from a single video~\cite{gao2017tall,hendricks2017mcn} to corpus-level retrieval that first finds the video and then the moment (VCMR)~\cite{lei2020tvr}, and in output, with QVHighlights~\cite{lei2021qvhighlights} asking for a saliency ranking of highlight clips alongside the grounded moment.

\begin{figure}[!ht]
    \centering
    \includegraphics[width=\linewidth]{figures/part3/section6/temporal_localization_task_settings.pdf}
    \caption{Six settings of temporal localization. Top row (TAL): (a) offline detection of $K$ classes; (b) online detection of action starts~\cite{shou2018odas,gao2019startnet}; (c) weakly supervised training from video-level labels. Bottom row (TSG): (d) single-video sentence grounding~\cite{gao2017tall,hendricks2017mcn}; (e) corpus-level moment retrieval~\cite{lei2020tvr}; (f) joint grounding and highlight detection~\cite{lei2021qvhighlights}. The online panel separates event onset, detection time, and the current observation boundary; future frames are unavailable. Dashed intervals in the weak-supervision panel denote inferred hypotheses, not boundary annotations. Scene examples, confidence values, and saliency ranks are illustrative.}
    \label{fig:p3_s6_4_2_task_settings}
\end{figure}

\subsubsection{Methodology}
Figure~\ref{fig:p3_s6_4_2_method_overview} traces both tasks along the same path from hand-designed pipelines to end-to-end Transformers.

The TAL row is the detection lineage of Section~\ref{sec:action_detection}: sliding windows over hand-crafted features such as improved dense trajectories~\cite{wang2013idt}, then two-stage proposal-plus-classification systems (R-C3D~\cite{xu2017rcc3d}, SSN~\cite{zhao2017temporal}, TAL-Net~\cite{chao2018talnet}) with dedicated boundary-sensitive proposal generators (BSN~\cite{lin2018bsn}, BMN~\cite{lin2019bmn}, G-TAD~\cite{xu2020gtad}), then anchor-free dense detectors and query-based set prediction. Three developments go beyond that treatment. Dense single-stage detectors were scaled by ActionFormer~\cite{zhang2022actionformer}, a multi-scale Transformer encoder over the feature pyramid, and TriDet~\cite{shi2023tridet}, which replaces self-attention with a convolutional scalable-granularity-perception module and a trident boundary head; both still regress boundaries at every position and rely on non-maximum suppression. DETR-style set prediction removes anchors and suppression: RTD-Net~\cite{tan2021rtdnet} applies a relaxed Transformer decoder to proposal generation and TadTR~\cite{liu2022tadtr} decodes the full action set end to end~\cite{liu2022e2etad,kim2024tetad}. And masked video pretraining~\cite{tong2022videomae,wang2023videomaev2} now supplies the backbones that these heads are fine-tuned on.

TSG began with cross-modal matching. CTRL~\cite{gao2017tall} embeds the sentence and candidate clips in a shared space and ranks clips, with boundary regression on top; MCN~\cite{hendricks2017mcn} does the same for DiDeMo with global video context. Later models reason over the candidates jointly: MAN~\cite{zhang2019man} refines moments iteratively through a graph over proposals, 2D-TAN~\cite{zhang2020twodtan} scores every (start, end) pair on a two-dimensional temporal-adjacency map so that neighboring candidates inform each other, and VSLNet~\cite{zhang2020vslnet} drops proposals altogether and predicts the start and end positions as span extraction, in analogy to reading comprehension. Moment-DETR~\cite{lei2021qvhighlights} brought set prediction to grounding, decoding moments and saliency scores directly from an encoder--decoder, while TVR~\cite{lei2020tvr} extended grounding to corpus scale with a coarse-to-fine search over subtitles and frames.

\subsubsection{Datasets and Metrics}

Table~\ref{tab:temporal_localization_datasets} lists the benchmarks; the TAL rows are those of Section~\ref{sec:action_detection}, retained for comparison. The TSG benchmarks move from a closed domain (TACoS~\cite{tacos2013}, cooking) to open-domain sentences over indoor activity (Charades-STA~\cite{gao2017tall}), Flickr videos with multiple references per query (DiDeMo~\cite{hendricks2017mcn}), and dense paragraph-level captions (ActivityNet Captions~\cite{krishna2017densecap}); YouCook2~\cite{zhou2017youcook2} adds procedural structure, TVR~\cite{lei2020tvr} multimodal (subtitle) context and corpus-level retrieval, and QVHighlights~\cite{lei2021qvhighlights} saliency labels for joint highlight detection.

\begin{table}[t]
\centering
\caption{Representative datasets for temporal localization; ``Task'' distinguishes category-conditioned TAL from text-conditioned TSG. Counts are indicative of commonly used versions and splits.}
\label{tab:temporal_localization_datasets}
\small
\setlength{\tabcolsep}{6pt}
\begin{tabularx}{\linewidth}{@{}l c c c X@{}}
\toprule
\textbf{Task} & \textbf{Dataset} & \textbf{Year} & \textbf{Scale / Classes} & \textbf{Temporal annotations} \\
\midrule
TAL & THUMOS14~\cite{jiang2014thumos} & 2014 & 200 train + 213 test videos / 20 cls & Frame-accurate, untrimmed \\
TAL & ActivityNet v1.3~\cite{caba2015activitynet} & 2016 & $\sim$20k videos / 200 cls & Untrimmed, diverse \\
TAL & HACS Segments~\cite{zhao2019hacs} & 2019 & 50k videos / 200 cls; 139k segments & Clip+segment, dense \\
TAL & EPIC-KITCHENS-100~\cite{damen2020epic100} & 2020 & 700 egocentric videos / verb-noun labels & Egocentric, narration-aligned \\
TAL & MultiTHUMOS~\cite{yeung2018every} & 2015 & $\sim$400 videos / 65 cls; $\sim$38.7k ann. & Dense, multi-label \\
TAL & FineAction~\cite{liu2022fineaction} & 2022 & 17k videos / 106 cls; 103k instances & Fine-grained, precise \\
\midrule
TSG & TACoS~\cite{tacos2013} & 2013 & 127 cooking videos; $\sim$18.8k query-segment pairs & Cooking, aligned \\
TSG & Charades-STA~\cite{gao2017tall} & 2017 & $\sim$6.6k videos; 16,128 pairs & Charades-based pairs \\
TSG & DiDeMo~\cite{hendricks2017mcn} & 2017 & $\sim$10k Flickr videos; $\sim$40k queries & Multi-reference, ranking \\
TSG & ActivityNet Captions~\cite{krishna2017densecap} & 2017 & 20k videos; $\sim$100k sentences & Dense captions \\
TSG & YouCook2~\cite{zhou2017youcook2} & 2018 & 2k untrimmed cooking videos / 89 recipes & Cooking procedures \\
TSG & TVR~\cite{lei2020tvr} & 2020 & 21.8k clips (60--90 s) from 6 TV shows; 109k queries & Subtitles, multimodal \\
TSG & QVHighlights~\cite{lei2021qvhighlights} & 2021 & $\sim$10k YouTube videos; query+saliency & Query-saliency, highlights \\
\bottomrule
\end{tabularx}
\end{table}

Temporal IoU between a predicted and a ground-truth segment is the primitive for both tasks, and TAL is scored by mAP at fixed tIoU thresholds (per-threshold on THUMOS14, averaged over $[0.5{:}0.95]$ on ActivityNet) and by AR@AN for proposal quality~\cite{anet2019proposals}, as defined in Section~\ref{sec:action_detection}. TSG uses retrieval metrics instead. For queries $\mathcal{Q}$, ranked predictions $\hat{m}_j$, and reference segments $G(q)$,
\begin{equation}
\mathrm{R}@n,\mathrm{IoU}{=}\theta \,=\, \frac{1}{|\mathcal{Q}|}\sum_{q\in\mathcal{Q}} \mathbf{1}\!\left(\max_{1\le j\le n}\;\max_{g\in G(q)} \mathrm{tIoU}(\hat{m}_j,g)\ge \theta\right)
\end{equation}
is the fraction of queries for which some top-$n$ prediction overlaps a reference by at least $\theta$ (commonly $n\in\{1,5\}$, $\theta\in\{0.3,0.5,0.7\}$), and $\mathrm{mIoU} = \frac{1}{|\mathcal{Q}|}\sum_{q} \max_{g\in G(q)} \mathrm{tIoU}(\hat{m}_1,g)$ scores the top prediction alone; QVHighlights adds mAP over tIoU thresholds~\cite{lei2021qvhighlights}. Recall at a threshold binarizes overlap and ignores the quality of the ranking, which motivated ranking-sensitive alternatives such as AxIoU~\cite{togashi2022axiou}, and annotator disagreement about where a described moment begins and ends puts a ceiling on how precisely any of these can be measured.

Table~\ref{tab:temporal_localization_performance} gives R@1 results for the TSG methods above; TAL numbers are in the table of Section~\ref{sec:action_detection}. Relational reasoning over candidates (2D-TAN) and span prediction (VSLNet) improved the listed Charades-STA R@1 from CTRL's 23.6/8.9 to 39.7/23.3 and 54.2/35.2, respectively, and Moment-DETR adds a further step on Charades-STA without proposals or suppression; its QVHighlights numbers are the reference point for joint grounding and highlight detection.

\begin{table}[!t]
    \centering
    \fontsize{9}{11}\selectfont
    \setlength{\tabcolsep}{1.2mm}
    \caption{Temporal sentence grounding, R@1 at tIoU 0.5 / 0.7 (\%). Charades-STA and ActivityNet Captions values as reported in the 2D-TAN~\cite{zhang2020twodtan}, VSLNet~\cite{zhang2020vslnet}, and Moment-DETR~\cite{lei2021qvhighlights} papers (visual features differ across methods as in the originals; MAN uses I3D features); QVHighlights test values from Moment-DETR without pretraining, with mAP averaged over tIoU thresholds.}
    \label{tab:temporal_localization_performance}
    \begin{tabular}{L{0.26\textwidth} L{0.26\textwidth} C{0.13\textwidth} C{0.15\textwidth} C{0.11\textwidth}}
    \toprule
    \textbf{Method} & \textbf{Design} & \textbf{Charades-STA} & \textbf{ActivityNet Captions} & \textbf{QVHighlights} \\
     & & R@1 0.5 / 0.7 & R@1 0.5 / 0.7 & R@1 0.5 / 0.7 (mAP) \\
    \midrule
    CTRL~\cite{gao2017tall}             & Clip--sentence matching      & 23.6 / 8.9  & 29.0 / 10.3 & -- \\
    MAN~\cite{zhang2019man}             & Iterative graph adjustment   & 46.5 / 22.7 & --          & -- \\
    2D-TAN~\cite{zhang2020twodtan}      & 2D temporal-adjacency map    & 39.7 / 23.3 & 44.5 / 26.5 & -- \\
    VSLNet~\cite{zhang2020vslnet}       & Span prediction              & 54.2 / 35.2 & 43.2 / 26.2 & -- \\
    Moment-DETR~\cite{lei2021qvhighlights} & Set prediction              & 53.6 / 31.4 & --          & 52.9 / 33.0 (30.7) \\
    \bottomrule
    \end{tabular}
\end{table}

\subsubsection{Challenges}

Boundaries are soft. Actions span sub-second gestures to multi-minute activities, which no single temporal resolution serves, and annotators disagree about where an instance starts and ends; a sentence adds its own ambiguity about which part of an activity it describes. The tIoU-thresholded metrics treat these fuzzy boundaries as exact, so part of the measured error is irreducible~\cite{togashi2022axiou}, and probabilistic boundary modeling with calibrated confidences is the natural response.

The two tasks are converging. A closed label set limits TAL, while TSG is bounded by the vocabulary of its query sets; text prompts in place of class labels~\cite{liberatori2024t3al} and set prediction for moments~\cite{lei2021qvhighlights} point to a single architecture that grounds both categories and sentences. Doing this over minutes of video, across camera viewpoints (egocentric footage in particular), with audio, narration, or noisy subtitles as further evidence~\cite{lei2020tvr,lei2021qvhighlights}, and causally for streaming use~\cite{shou2018odas} is where the current systems fall short: memory-efficient long-context encoders and corpus-scale coarse-to-fine retrieval are the two ingredients most actively developed.








% \newpage

\section{Multimodal Spatial Perception}
\label{sec:p3-sec6-5-multimodal-spatial-perception}
% \wsq{将视觉与语言、音频、触觉等模态结合，增强空间感知范围与粒度。}

\FramedInitial{T}{he} preceding sections perceive space largely through vision. Yet vision alone is often incomplete: it is blind past occluders, silent about what lies outside the frame, and mute on which entity a person actually means. Multimodal spatial perception closes these gaps by grounding spatial understanding in signals beyond the image. The section takes two complementary couplings: audio-visual spatial localization, which uses sound to place events that vision may miss, and vision-language grounding, which ties linguistic reference to the specific entity or region it denotes. Neither adds a new spatial primitive; each supplies a channel of evidence that makes perception more complete where a single modality falls short.


\subsection{Audio-Visual Spatial Localization}
\label{sec:av_localization}
% \wsq{声源定位与空间目标匹配（如声音中找人）}

Humans localize a sound by combining auditory cues (interaural level and time differences, spectral shaping) with visual context (lip motion, object dynamics)~\cite{blauert1997spatial}. Audio-visual spatial localization (AV-SL) asks a machine to do the same: find the pixel region, object, or 3D point that a sound comes from, under reverberation, noise, occlusion, and several simultaneous sources~\cite{arandjelovic2017look,owens2018audio}. The two modalities are complementary. Binaural or array audio supplies directional cues, while distance requires additional geometric, near-field, or learned evidence; vision says which entities are plausible sounders and separates acoustically confusable ones (two similar instruments, only one playing). Spatial audio for AR/VR~\cite{morgado2018self}, teleconferencing~\cite{tao2021talknet}, and audio-guided navigation~\cite{chen2020soundspaces} are the main consumers.

\subsubsection{Task Definition}
The input is an audio stream $A$ (mono, binaural, or multichannel, optionally ambisonic~\cite{zotter2019ambisonics}) time-aligned with a video $V$ (RGB, depth, or multi-view) of the same interval; the audio may contain zero, one, or several sources, and the video may show silent objects that look like sounders. The output localizes the sounding entity in a form that depends on the variant: a 2D heatmap or mask, a face or body box with a speaking state, an event class with its temporal extent and direction of arrival (DOA, azimuth and elevation) and optionally distance~\cite{politis2023dataset}, a 2D location with depth, or a 3D point or trajectory in camera or world coordinates. For visible-region localization, off-screen sources or audio with no visible source should yield a null prediction; 3D acoustic localization may still estimate an off-screen source.

Five variants are in common use. Visual sound source localization (VSSL) highlights the sounding region of a frame or clip as a class-agnostic heatmap, in single- and multi-source settings and increasingly with negative or off-screen audio as a robustness test. Active speaker detection (ASD) decides who is speaking and localizes the talker as a face track, often under streaming constraints~\cite{tao2021talknet}. Audio-visual sound event localization and detection (AV-SELD) couples what and when with where: event class, temporal extent, DOA, and in recent benchmarks distance~\cite{politis2023dataset}. The 2.5D variant adds depth or range to a 2D localization from monocular depth~\cite{DepthAnything}, stereo, or echoes~\cite{gao2020visualechoes}; and 3D DOA estimation returns a spatial bearing, whereas source-position or trajectory estimation additionally uses range or triangulation constraints from calibrated cameras, arrays, or egocentric rigs~\cite{grauman2022ego4d}.

\subsubsection{Methodology}
Figure~\ref{fig:p3_s6_5_1_method_overview} maps the five variants onto four method families.

\begin{figure}[!t]
    \centering
    \includegraphics[width=\linewidth]{figures/part3/section6/audio_video_spatial_localization.pdf}
    \caption{Audio-visual spatial localization. \textit{Top}: task outputs for (i) VSSL, (ii) ASD, (iii) AV-SELD, (iv) 2.5D localization, and (v) 3D localization and tracking. \textit{Bottom}: method families (A) cross-modal correspondence, (B) audio--lip synchrony, (C) joint spatial event prediction, and (D) geometric lifting or fusion; neutral connectors mark principal task--family associations, not computational flow. Directional arrows inside panels indicate prediction or information flow. The 2.5D panel links a selected image location to depth; 3D position additionally requires range or triangulation constraints. Scene imagery and overlays are illustrative.}
    \label{fig:p3_s6_5_1_method_overview}
\end{figure}

\paragraph{Cross-Modal Correspondence (VSSL).}
VSSL is learned almost entirely without location labels, from the co-occurrence of sound and image. An audio encoder maps the clip to an embedding $\mathbf{a}$, a visual encoder maps the frame to a feature map $\mathbf{v}(\mathbf{x})$, and the localization map is the per-position similarity
\begin{equation}
    S(\mathbf{x}) = \frac{\mathbf{a}^{\top}\mathbf{v}(\mathbf{x})}{\|\mathbf{a}\|\,\|\mathbf{v}(\mathbf{x})\|},
\end{equation}
trained so that matching audio-image pairs score higher than mismatched ones~\cite{senocak2018learning,owens2018audio}; Sound of Pixels~\cite{zhao2018sound} instead mixes two clips and asks the network to separate them conditioned on pixel features, which forces the same alignment. Later work targets the failure modes of co-occurrence. LVS~\cite{chen2021localizing} mines hard negatives inside the same image (regions that look like sounders but do not match the audio) and introduces the VGG-SS localization benchmark; deep multimodal clustering~\cite{hu2019deep} and coarse-to-fine mixture objectives~\cite{qian2020multiple} handle several sources at once; and SLAVC~\cite{mo2022closer} shows that earlier methods depend on early stopping and adds momentum encoders and extreme visual dropout so that the alignment does not collapse onto background.

\paragraph{Audio--Lip Alignment (ASD).}
ASD runs a face detector and tracker and then decides, per track and per frame, whether the visible mouth produced the audio. The decision rests on audio-visual synchrony: a visual encoder over the lip region and an audio encoder over the spectrogram are compared across time~\cite{afouras2020self}, and TalkNet~\cite{tao2021talknet} adds long-term temporal modeling of both streams with cross-attention, which is what separates a brief pause from a non-speaker. Active Speakers in Context~\cite{alcazar2020active} scores each candidate jointly with the other faces in the frame, and ASDNet~\cite{kopuklu2021how} makes the three stages explicit (per-frame audio-visual encoding, inter-speaker relation modeling, temporal modeling), with each stage trained for its role.

\paragraph{Joint Event and DOA Prediction (AV-SELD).}
SELDnet~\cite{adavanne2018sound} set the template: a convolutional recurrent network over multichannel spectrogram features with two heads, one classifying active events per frame and one regressing their DOA. ACCDOA~\cite{shimada2021accdoa} folds the two into one output, a vector per class whose direction is the DOA and whose length is the activity, which removes the need to balance two losses; the DCASE challenge tracks~\cite{politis2020overview} standardized the protocol, and STARSS23~\cite{politis2023dataset} added video and source distance.

\paragraph{Geometric Lifting (2.5D and 3D).}
Lifting a 2D localization to metric space uses geometry rather than learned correspondence alone. VisualEchoes~\cite{gao2020visualechoes} shows that echoes of an emitted chirp carry depth: the time and spectrum of the return, fused with the image, improve monocular depth estimation and give the 2.5D variant its range. Audio-visual spatial alignment~\cite{morgado2020learning} learns representations by predicting whether 360$^\circ$ video and ambisonic audio are spatially consistent, a pretext that encodes direction, and self-supervised spatial audio generation~\cite{morgado2018self} learns the inverse mapping from 360$^\circ$ video to ambisonics. Full 3D systems combine classical array processing (time differences of arrival from GCC-PHAT~\cite{knapp1976generalized}, steered-response-power beamforming~\cite{dibiase2001robust}) with calibrated cameras and triangulation, and in simulation (SoundSpaces~\cite{chen2020soundspaces,chen2022soundspaces}) with the agent's own motion.

\subsubsection{Datasets and Metrics}

Table~\ref{tab:p3_s6_5_1_avsl_tasks} pairs each variant with its benchmarks and metrics; there is no shared benchmark across variants because their outputs differ.

\begin{table}[!ht]
\centering
\caption{AV-SL task variants with representative datasets, outputs, and metrics. Newer VSSL test suites add negative or off-screen audio to measure robustness.}
\label{tab:p3_s6_5_1_avsl_tasks}
\small
\begin{tabular}{@{}L{0.12\textwidth} L{0.34\textwidth} L{0.20\textwidth} L{0.24\textwidth}@{}}
\toprule
\textbf{Task} & \textbf{Representative Datasets} & \textbf{Primary Outputs} & \textbf{Typical Metrics} \\
\midrule
VSSL (2D) & Flickr-SoundNet~\cite{senocak2018learning}, MUSIC~\cite{zhao2018sound}, VGG-SS~\cite{chen2021localizing} (from VGG-Sound~\cite{chen2020vggsound}), AVSBench~\cite{zhou2022avs} & heatmap / mask & cIoU, AUC; AP with negatives \\
ASD & AVA-ActiveSpeaker~\cite{roth2020ava} & talker boxes + speaking state & mAP \\
AV-SELD & STARSS23~\cite{politis2023dataset}, DCASE SELD tracks~\cite{politis2020overview} & event + time + DOA (+ distance) & F-score, error rate, DOA error, distance error \\
2.5D & Replica / Matterport3D echo splits~\cite{gao2020visualechoes} & 2D + depth & AbsRel, RMSE \\
3D DOA / trajectory & STARSS23 positions~\cite{politis2023dataset}, SoundSpaces~\cite{chen2020soundspaces} & DOA / 3D point / trajectory & angular error, positional error \\
\bottomrule
\end{tabular}
\end{table}

VSSL is scored by thresholding the localization map and computing consensus intersection over union (cIoU) against the annotations; reported metrics include the fraction of test images with cIoU above 0.5 (localization accuracy) and the area under the IoU-threshold curve (AUC)~\cite{senocak2018learning}; the extended VGG-SS and Flickr protocols add clips whose audio has no visible source and report detection-style average precision requiring both visible-source presence and correct localization~\cite{chen2021localizing,mo2022closer}. ASD uses detection-style mAP over face tracks with speaking labels~\cite{roth2020ava}. SELD combines detection (F-score, error rate) with localization (DOA error, localization recall)~\cite{politis2020overview}, while STARSS23 additionally provides source-distance annotations~\cite{politis2023dataset}; 2.5D inherits the depth metrics of Section~\ref{sec:depth_estimation}, and 3D settings report angular and positional error.

Table~\ref{tab:p3_s6_5_1_avsl_performance} collects the two variants with a stable leaderboard. On VSSL the negative-audio protocol changes the ranking: LVS leads the earlier methods on the standard Flickr test but drops to 9.8 AP when silent clips are mixed in, whereas SLAVC holds both. ASD on AVA-ActiveSpeaker has moved from 79.2 mAP for the dataset baseline to above 93 with explicit temporal and inter-speaker modeling, and the remaining differences among top systems are about latency and streaming rather than accuracy. AV-SELD has no such table: each DCASE edition changes the data and the metric, and distance-aware evaluation is not directly comparable to the DOA-only tracks that SELDnet and ACCDOA were evaluated on.

\begin{table}[!t]
    \centering
    \fontsize{9}{11}\selectfont
    \setlength{\tabcolsep}{1.2mm}
    \caption{Visual sound source localization and active speaker detection. VSSL: cIoU (\%) on the Flickr-SoundNet and VGG-SS test sets, and AP (\%) on their extended versions with negative audio, all models trained on VGG-Sound 144k, as reported in the SLAVC comparison~\cite{mo2022closer}. ASD: mAP (\%) on AVA-ActiveSpeaker val, from the TalkNet~\cite{tao2021talknet} and ASDNet~\cite{kopuklu2021how} papers.}
    \label{tab:p3_s6_5_1_avsl_performance}
    \begin{tabular}{L{0.26\textwidth} L{0.26\textwidth} C{0.09\textwidth} C{0.09\textwidth} C{0.09\textwidth} C{0.09\textwidth}}
    \toprule
    \rowcolor{lightgrey} \multicolumn{6}{l}{$\bullet$ \textit{Visual sound source localization}} \\
    \textbf{Method} & \textbf{Design} & \textbf{Flickr cIoU} & \textbf{VGG-SS cIoU} & \textbf{Ext.\ Flickr AP} & \textbf{Ext.\ VGG-SS AP} \\
    \midrule
    Attention10k~\cite{senocak2018learning} & Audio-guided attention   & 42.3 & 18.5 & --   & -- \\
    DMC~\cite{hu2019deep}                   & Multimodal clustering    & 55.6 & 23.9 & --   & -- \\
    LVS~\cite{chen2021localizing}           & Hard negatives           & 71.9 & 34.4 & 9.8  & 5.2 \\
    SLAVC~\cite{mo2022closer}               & Momentum encoder, dropout & 83.2 & 37.2 & 51.6 & 33.0 \\
    \midrule
    \rowcolor{lightgrey} \multicolumn{6}{l}{$\bullet$ \textit{Active speaker detection (AVA-ActiveSpeaker val, mAP)}}\\
    \hiderowcolors \textbf{Method} & \textbf{Design} & \textbf{mAP} & & & \\
    \midrule
    AVA baseline~\cite{roth2020ava}         & Audio-visual fusion      & 79.2 & & & \\
    ASC~\cite{alcazar2020active}            & Inter-speaker context    & 87.1 & & & \\
    TalkNet~\cite{tao2021talknet}           & Long-term temporal, cross-attention & 92.3 & & & \\
    ASDNet~\cite{kopuklu2021how}            & Three-stage architecture & 93.5 & & & \\
    \bottomrule
    \end{tabular}
\end{table}

\subsubsection{Challenges}
Co-occurrence is not causation. A model that has learned which objects usually make which sounds will localize a guitar whether or not it is playing, and the negative-audio protocols~\cite{chen2021localizing,mo2022closer} exist because the standard tests rewarded exactly this; the open problem is grounding on the acoustic evidence itself, which matters most for off-screen and silent cases where the right answer is no localization.

Going from 2D saliency to metric 3D requires semantics and geometry together. Depth foundation models~\cite{DepthAnything} (Section~\ref{sec:depth_estimation}) and self-supervised visual features~\cite{dinov2} now supply priors that AV-SL pipelines previously learned from scratch, and STARSS23~\cite{politis2023dataset} provides joint event, time, DOA, and distance annotations, supporting distance-aware extensions; multi-object objectives for overlapping events are the current frontier there.

Multiple simultaneous sources, identity through occlusion, and on-screen/off-screen transitions over long videos~\cite{qian2020multiple,tian2018audio} are the AV-SL form of the identity problem of Section~\ref{sec:object_tracking} and the pose-continuity problem of Section~\ref{sec:object_pose}. Streaming use adds causal encoders, low-latency fusion, calibration drift, and graceful fallback when one modality is missing~\cite{tao2021talknet,kopuklu2021how}.
% TODO: anchor §8.x once §8 (Spatial Reconstruction and Generation) subsections stabilize --- AV-SL is a natural consumer of 4D scene reconstruction (spatial audio rendering, dynamic-scene physical reasoning) per the §6.4.1 (Scene Flow) §Challenges ¶3 convergence narrative.



\subsection{Vision-Language Grounding}
\label{sec:vl_grounding}
% \wsq{``把手放在右边的杯子上'' -> 找到右边的杯子; 与 RefCOCO / Refer-Youtube-VOS 等任务相关}

Vision-language grounding ties a linguistic reference to the element of a scene it denotes. Given \textit{``Put your hand on the small mug on the right side of the table behind the blue book''}, a grounded system must resolve the noun phrase, its attributes, and its spatial relations to one particular mug in the image; this referential link is what instruction following, multimodal dialogue, and human-robot interaction build on.

\subsubsection{Task Definition}
The input is a scene $S$ (an image, a video, or a 3D environment) and a query $Q$, from a bare noun phrase (``the red cup'') to a sentence with attributes, relations, actions, or an implicit goal (``the person who is waving'', ``something to open the box with''). The output $\mathcal{R}(Q,S)$ is the referred region in the form the variant requires: a box, a mask, a set of regions (possibly empty), a per-frame mask track, or a 3D box. Two formulations recur through the methods below: ranking a set of candidate regions $r$ by a cross-modal score, or regressing the region directly,
\begin{equation}
    \hat{r} = \arg\max_{r \in \mathcal{C}(S)} s\big(\phi(Q), \psi(S, r)\big)
    \qquad \text{or} \qquad
    \hat{r} = f_\theta(S, Q),
\label{eq:p3_s6_5_2_grounding}
\end{equation}
with $\phi$ and $\psi$ the text and region encoders and $\mathcal{C}(S)$ a proposal set.

\begin{figure}[!ht]
    \centering
    \includegraphics[width=0.99\linewidth]{figures/part3/section6/visual_language_graounding_task.pdf}
    \caption{Vision-language grounding task variants: (i) REC / RES~\cite{yu2016modeling,mao2016generation}; (ii) generalized REC / RES with zero, one, or multiple targets~\cite{liu2023gres}; (iii) reasoning grounding~\cite{lai2024lisa}; (iv) video referring and reasoning grounding~\cite{luo2026drV}; (v) 3D and embodied grounding~\cite{chen20263d}.}
    \label{fig:p3_s6_5_2_method_overview_task}
\end{figure}

Figure~\ref{fig:p3_s6_5_2_method_overview_task} lays out the five variants, which differ in scene modality, output granularity, and query complexity.

\textbf{Referring expression comprehension and segmentation (REC / RES).} One image, one expression, one target, returned as a box (REC) or a mask (RES). RefCOCO and RefCOCO+~\cite{yu2016modeling} (short phrases, the latter without location words) and RefCOCOg~\cite{mao2016generation} (longer descriptions), all using MS~COCO~\cite{lin2014microsoft} images, with expressions collected through ReferItGame~\cite{kazemzadeh2014referitgame} or Mechanical Turk, have been the standard benchmarks for a decade; Flickr30k Entities~\cite{plummer2015flickr30k,young2014image} grounds every entity mention of a caption at once.

\textbf{Generalized REC / RES (GREC / GRES).} An expression may match zero, one, or several objects (``the red carrots'' versus ``the unicorn''), which single-target methods cannot represent. gRefCOCO~\cite{liu2023gres}, Ref-ZOM~\cite{hu2023beyond}, PhraseCut~\cite{wu2020phrasecut}, and OmniLabel~\cite{schulter2023omnilabel} evaluate this setting, and MaskGroups-HQ~\cite{cao2025refer} generalizes the prompt itself to arbitrary groups of masks specified by text and reference masks. The output is a set, so the metric must be as well (see Datasets and Metrics below).

\textbf{Reasoning grounding.} Some queries require inference beyond noun matching: ``the loser'' needs a judgment about the people shown, ``turn on the TV'' needs the remote. CLEVR-Ref+~\cite{liu2019clevr} diagnoses this on rendered scenes with ground-truth reasoning programs; ReasonSeg~\cite{lai2024lisa} poses implicit, multi-step queries on real images with mask output, extended by MMR~\cite{jang2025mmr} (194K implicit queries, multi-target and part-level), ReasonSeg-Inst~\cite{yang2023lisa++} (instance level), and MC-Bench~\cite{xu2024mcbench} (instances across image pairs). Ref-Adv~\cite{dong2026ref} keeps the classic REC format but makes expressions minimally sufficient, negation-heavy, and paired with same-category distractors.

\textbf{Video referring and reasoning grounding.} The referent must be found in every frame through motion, appearance change, and occlusion, and the query may name it by what it does (``the man waving his hand''). Refer-Youtube-VOS~\cite{seo2020urvos}, Refer-DAVIS17~\cite{khoreva2018video}, A2D- and JHMDB-Sentences~\cite{gavrilyuk2018actor}, MeViS~\cite{ding2023mevis}, ReVOS~\cite{yan2024visa}, and ReasonVOS~\cite{bai2024one} add language to existing segmentation and actor-action datasets; Ref-VPS~\cite{bagchi2025refereverything} extends the referent from objects to processes such as smoke spreading or waves breaking, which have no object boundary at all. Tracking machinery is usually folded in to carry identity across frames~\cite{fang2025decoupled}.

\textbf{3D and embodied grounding.} The scene is a point cloud or a set of views, and the query refers to an object in 3D (``the chair next to the window''). ScanRefer~\cite{chen2020scanrefer} requires detecting the referent as a 3D box and Multi3DRefer~\cite{zhang2023multi3drefer} allows zero to many targets; ReferIt3D~\cite{achlioptas2020referit3d} provides the object instances and scores selection among them; Instruct3D~\cite{he2024segpoint} and Multi3DRes~\cite{wu20243d} ask for point-level masks from instruction-style queries. Embodied grounding makes the output an action: an agent told ``go pick up the red cup on the right'' must identify the cup and reach it~\cite{zhen20243d}.

\subsubsection{Methodology}

\begin{figure}[!t]
    \centering
    \includegraphics[width=0.99\linewidth]{figures/part3/section6/visual_language_graounding_method.pdf}
    \caption{Vision-language grounding method families: (i) proposal ranking, (ii) compositional reasoning, (iii) end-to-end grounding, (iv) grounded pre-training, (v) foundation models, and (vi) tool-augmented agents. Panels (i)-(iv) correspond to families (A)-(D) in the text, (v) to (E)-(F), and (vi) to (G); family (H), generation for grounding, is not depicted.}
    \label{fig:p3_s6_5_2_method_overview_method}
\end{figure}

The methods below fall into eight families, (A) to (H); Figure~\ref{fig:p3_s6_5_2_method_overview_method} illustrates (A) through (G). They are trajectories rather than eras: the oldest paradigm has resurfaced at foundation-model scale.

\paragraph{(A) Proposal Ranking.}
The first grounding models instantiate the ranking form of Eq.~\eqref{eq:p3_s6_5_2_grounding}: detect region proposals, embed them and the query in a joint space, and pick the best-scoring region. Deep Fragment Embeddings~\cite{karpathy2014deep} aligned sentence fragments with detected regions, and scene-graph CRFs~\cite{johnson2015image} grounded structured queries to regions for retrieval. Proposal quality bounded these methods and context-dependent expressions defeated them. The same two-stage shape has returned with far stronger proposals: RAS~\cite{cao2025refer} has a segmentation foundation model propose candidate masks and a mask-centric multimodal LLM \emph{select} the group that satisfies the prompt, and this non-autoregressive selection beats mask generation in both accuracy and latency.

\paragraph{(B) Compositional Reasoning.}
Modular networks parse the expression and route its parts to dedicated modules. CMN~\cite{hu2017modeling} decomposes an expression into a subject--relationship--object triplet handled by localization and relationship modules; MAttNet~\cite{yu2018mattnet} soft-parses phrases into appearance, location, and relationship attention modules; dynamic graph attention~\cite{yang2019dynamic} reasons over the linguistic structure on a scene graph. These models disambiguate compositional queries (``the man behind the woman in red'') but inherit every parsing error.

\paragraph{(C) End-to-End Grounding.}
The regression form of Eq.~\eqref{eq:p3_s6_5_2_grounding} removes proposals: a single network fuses the sentence embedding with the image feature map and predicts the box. One-Stage Grounding~\cite{yang2019fast} did this with a detector head, TransVG~\cite{deng2021transvg} with a pure Transformer in which a learnable token attends over visual and textual tokens and regresses coordinates. On the segmentation side, CRIS~\cite{wang2022cris} transfers CLIP's image-text alignment to the pixel level with text-to-pixel contrastive learning and made end-to-end RES the standard formulation.

\paragraph{(D) Grounded Pre-Training.}
Phrase-region alignment learned at scale turned grounding into an open-vocabulary capability. MDETR~\cite{kamath2021mdetr} trains a DETR-style set predictor on text-aligned detection corpora so that every box is tied to a span of the caption; GLIP~\cite{li2022grounded} reformulates detection as phrase grounding and scales to 27M region-caption pairs, giving zero-shot transfer; RegionCLIP~\cite{zhong2022regionclip} adds region-level contrastive alignment, and Grounding DINO~\cite{liu2024grounding} deep language fusion at every stage of a DINO detector, which made it the default open-vocabulary detector of later tool chains (family (G)). The residual weakness is the \emph{entity shortcut}: models latch onto object names and ignore the description. DesCo~\cite{li2023desco} counters it in the data, pre-training on LLM-generated rich descriptions with context-sensitive negative queries, which raises description-conditioned detection on LVIS and OmniLabel.

\paragraph{(E) Grounding MLLMs.}
The first multimodal LLMs built for grounding gave the language backbone a region interface. Kosmos-2~\cite{peng2023kosmos} emits location tokens that hyperlink phrases to boxes inside the generated text; Ferret~\cite{you2023ferret} and Ferret-v2~\cite{zhang2024ferret} encode boxes, points, and free-form regions as inputs at any resolution. Mask output comes from attaching a decoder: LISA~\cite{lai2024lisa} decodes a dedicated segmentation token's embedding into a mask, which is what makes reasoning segmentation possible, and LISA++~\cite{yang2023lisa++} extends it to instances within a dialogue; GLaMM~\cite{rasheed2024glamm} unifies phrase grounding, referring segmentation, and free-form conversation under a grounded-conversation objective; Sa2VA~\cite{yuan2025sa2va} couples a LLaVA-style backbone to SAM~2~\cite{ravi2024sam2} and carries dense grounding from images to video referring segmentation. Moondream~3's segmentation extension~\cite{reid2026moondream} dispenses with the decoder and writes the mask autoregressively as a vector path, refined by a lightweight head and trained with reinforcement learning on mask overlap.

\paragraph{(F) Generalist MLLMs.}
A parallel line keeps no grounding component at all: the model writes coordinates as plain text, and grounding is one capability among many. Florence-2~\cite{xiao2024florence} is the compact case, a sequence-to-sequence model trained on the 5.4B-annotation FLD-5B corpus that casts captioning, detection, grounding, and polygon-vertex segmentation into one text interface, with mask accuracy below the decoder-based models of family (E). The frontier generalists, Gemini~\cite{reid2024gemini}, Qwen2.5-VL and Qwen3-VL~\cite{bai2025qwen2.5vl}, InternVL~2.5 and InternVL~3~\cite{chen2024expanding,zhu2025internvl3}, and Seed1.5-VL~\cite{guo2025seed1.5vl}, now match or exceed the specialists on the RefCOCO family from prompting alone (Table~\ref{tab:p3_s6_5_2_grounding_performance}); Llama~3~\cite{grattafiori2024llama} folds phrase-to-box and phrase-to-mask annotations into its multimodal pre-training but reports no REC numbers.

\paragraph{(G) Tool-Augmented Agents.}
A third line delegates localization to specialist models under an outer control loop, which has moved from fixed pipelines to generated programs to learned tool use. Grounded SAM~\cite{ren2024grounded} chains Grounding DINO boxes into SAM~\cite{kirillov2023sam} masks without training; ViperGPT~\cite{suris2023vipergpt} has a code LLM write a Python program that calls perception APIs, so that logic and arithmetic are symbolic and perception is delegated, with interpretable intermediate steps. SegAgent~\cite{zhu2025segagent} treats mask annotation as a sequential decision process, imitating the click sequences of human annotators on an interactive segmenter; DeepEyes~\cite{zheng2025deepeyes} internalizes the tool, using reinforcement learning to make the model emit coordinates that trigger zoom-in crops inside its own chain of thought. SAM~3 Agent~\cite{carion2025sam} wraps SAM~3, whose concept prompts return every instance matching a noun phrase, with an MLLM that proposes phrases, inspects the returned masks, and selects those satisfying the full expression, which gives zero-shot referring and reasoning segmentation without grounding-specific training.

\paragraph{(H) Generation for Grounding.}
The newest family treats a large text-to-image or text-to-video generator as a grounding model, on the premise that generative pre-training already induces a fine-grained map from language to spatial and temporal structure. VD-IT~\cite{zhu2024exploring} uses a frozen text-to-video diffusion model as the feature extractor for referring video segmentation and finds its features more temporally consistent and better aligned with language than discriminative video backbones; ReferEverything~\cite{bagchi2025refereverything} keeps the generator intact and fine-tunes it to predict mask latents decoded by the frozen VAE, which transfers zero-shot to rare objects, ``stuff'', and the amorphous processes of Ref-VPS. Vision Banana~\cite{gabeur2026image} instruction-tunes a state-of-the-art image generator to render the task output as an image, segmentation as painting the referred region a prescribed color, and matches specialist systems zero-shot while keeping its generative ability. Grounding quality in this family rises with the underlying generator.

\subsubsection{Datasets and Metrics}

Table~\ref{tab:p3_s6_5_2_grounding_tasks} lists the benchmarks by variant. Annotation cost rises steeply from a 2D box to a per-frame mask to a 3D box, so image-level REC/RES corpora dominate in size and maturity, video referring sets are an order of magnitude smaller, and 3D corpora smaller still. Within the image benchmarks a generational split has opened. RefCOCO/+/g now serve mainly as training corpora: specialists and generalist MLLMs report Acc@0.5 above 90\% on them, an audit found label errors in 14\%, 24\%, and 5\% of RefCOCO, RefCOCO+, and RefCOCOg~\cite{chen2025revisiting}, and much of the score survives shuffling the word order of the expression~\cite{dong2026ref}. The evaluation-oriented replacements are built to expose what saturation hides: Ref-L4~\cite{chen2025revisiting} (cleaned labels, long expressions, a 365-category vocabulary), Ref-Adv~\cite{dong2026ref} (shortcut-suppressed, where the same models fall to roughly 60\%), HC-RefLoCo~\cite{wei2024hcrefloco}, GigaGrounding~\cite{ma2024gigaground}, and GroundingME~\cite{li2026groundingme}, which scores discrimination, spatial reasoning, limited visibility, and the rejection of ungroundable queries, and on which the best model resolves 45\% of samples while most score zero on rejection.

\begin{table}[!ht]
\centering
\caption{Common vision-language grounding tasks with representative datasets and typical evaluation metrics; several tasks share datasets (\eg, RefCOCO underlies both REC and RES).}
\label{tab:p3_s6_5_2_grounding_tasks}
\fontsize{9}{11}\selectfont
\setlength{\tabcolsep}{4pt}
\begin{tabular}{L{0.18\textwidth} L{0.34\textwidth} L{0.18\textwidth} L{0.22\textwidth}}
\toprule
\textbf{Task} & \textbf{Representative Datasets} & \textbf{Primary Outputs} & \textbf{Typical Metrics} \\
\midrule
REC / RES & RefCOCO/+/g, Flickr30k Entities, Ref-L4, Ref-Adv, GroundingME & box / mask & Acc@0.5; gIoU / cIoU (RES) \\
GREC / GRES & gRefCOCO, OmniLabel, PhraseCut, Ref-ZOM, MaskGroups-HQ & multi-box / multi-mask & F1, N-acc, gIoU / cIoU, AP$_{\text{descr}}$ \\
Reasoning grounding & CLEVR-Ref+, ReasonSeg, ReasonSeg-Inst, MMR, MC-Bench & mask (or box) & gIoU / cIoU, Acc@0.5 \\
Video referring & MeViS, Refer-YT-VOS, Refer-DAVIS17, ReVOS, A2D-Sent., Ref-VPS & per-frame mask & $\mathcal{J}\&\mathcal{F}$; mAP (A2D-Sent.) \\
3D / Embodied & ScanRefer, ReferIt3D, Multi3DRefer, Instruct3D & 3D box / 3D mask & Acc@0.25/0.5, F1@0.5, mIoU \\
\bottomrule
\end{tabular}
\end{table}

Every metric here is built on the IoU of predicted and reference regions (Section~\ref{sec:object_tracking} gives the mask form). REC reports Accuracy@0.5, the fraction of expressions whose box reaches IoU $\ge 0.5$, with stricter thresholds and the threshold-averaged mAcc~\cite{chen2025revisiting} increasingly added; phrase grounding, where several proposals may be returned, reports Recall@$k$. For set-valued outputs (GREC / GRES) the prediction is compared with the ground-truth set through per-expression F1 over boxes or masks~\cite{hu2023beyond} and a no-target accuracy for expressions with no referent~\cite{liu2023gres}, and OmniLabel's AP$_{\text{descr}}$ computes average precision over free-form descriptions including negative ones~\cite{schulter2023omnilabel}. Segmentation reports two aggregates that differ in weighting: gIoU (also called mIoU) averages per-sample IoU and so weights every expression equally, whereas cIoU accumulates intersection and union over the whole set before dividing and so weights large objects more; reasoning and referring segmentation papers report both~\cite{lai2024lisa}, often with Precision@$X$ for the distribution of mask quality and a boundary IoU for edge accuracy~\cite{reid2026moondream}. Video referring segmentation uses the $\mathcal{J}\&\mathcal{F}$ composite of Section~\ref{sec:object_tracking}.

Table~\ref{tab:p3_s6_5_2_grounding_performance} reports the two variants on which the families can be compared. On REC the step from proposal ranking and modular reasoning (MAttNet) to end-to-end regression (TransVG) and then to grounded pre-training (MDETR, Grounding DINO) is about 14 points of RefCOCO Acc@0.5, and the generalist MLLMs (Qwen2.5-VL, InternVL~2.5) sit at or above the best specialist without benchmark-specific training. On RES the grounding MLLMs with a mask decoder have passed the end-to-end specialists (CRIS to Sa2VA is eleven cIoU points on RefCOCO), and Sa2VA's video columns show the same model carrying over to referring video segmentation; on ReasonSeg, LISA-7B reaches 52.9 gIoU after fine-tuning where open-vocabulary segmenters stay below 30~\cite{lai2024lisa}. The tool-augmented and generative families are evaluated zero-shot on subsets of these benchmarks and are not directly comparable here. What the table cannot show is cross-variant breadth: methods strong on image REC routinely degrade on video, reasoning, and 3D grounding, and no single architecture leads all five variants.

\begin{table}[!ht]
    \centering
    \fontsize{9}{11}\selectfont
    \setlength{\tabcolsep}{1.5mm}
    \caption{Representative grounding methods. REC: Acc@0.5 (\%) on RefCOCO val / testA / testB and RefCOCOg val, from the Grounding DINO~\cite{liu2024grounding}, Ferret~\cite{you2023ferret}, and Qwen2.5-VL~\cite{bai2025qwen2.5vl} comparisons. RES: cIoU (\%) on the RefCOCO / RefCOCO+ / RefCOCOg val splits and $\mathcal{J}\&\mathcal{F}$ on MeViS and Ref-DAVIS17, from the Sa2VA comparison~\cite{yuan2025sa2va} (CRIS from its own paper, ResNet-101).}
    \label{tab:p3_s6_5_2_grounding_performance}
    \begin{tabular}{L{0.26\textwidth} L{0.22\textwidth} C{0.10\textwidth} C{0.10\textwidth} C{0.10\textwidth} C{0.10\textwidth}}
    \toprule
    \rowcolor{lightgrey} \multicolumn{6}{l}{$\bullet$ \textit{Referring expression comprehension (Acc@0.5)}} \\
    \multirow{2}{*}{\textbf{Method}} &  \multirow{2}{*}{\textbf{Family}} & \multicolumn{3}{c}{\textbf{RefCOCO}} & \textbf{RefCOCOg} \\
    & & val & testA & testB &  val \\
    \midrule
    MAttNet~\cite{yu2018mattnet}             & (B) modular attention      & 76.7 & 81.1 & 70.0 & 66.6 \\
    TransVG~\cite{deng2021transvg}           & (C) end-to-end Transformer & 81.0 & 82.7 & 78.4 & 68.7 \\
    MDETR~\cite{kamath2021mdetr}             & (D) grounded pre-training  & 86.8 & 89.6 & 81.4 & 81.6 \\
    Grounding DINO-L~\cite{liu2024grounding} & (D) grounded pre-training  & 90.6 & 93.2 & 88.2 & 86.1 \\
    Ferret-13B~\cite{you2023ferret}          & (E) grounding MLLM         & 89.5 & 92.4 & 84.4 & 85.8 \\
    Qwen2.5-VL-72B~\cite{bai2025qwen2.5vl}   & (F) generalist MLLM        & 92.7 & 94.6 & 89.7 & 89.9 \\
    InternVL2.5-78B~\cite{chen2024expanding} & (F) generalist MLLM        & 93.7 & 95.6 & 92.5 & 92.7 \\
    \midrule
    \rowcolor{lightgrey} \multicolumn{6}{l}{$\bullet$ \textit{Referring segmentation (cIoU) and video referring segmentation ($\mathcal{J}\&\mathcal{F}$)}}\\
    \hiderowcolors \textbf{Method} & \textbf{Family} & \textbf{RefCOCO} & \textbf{RefCOCO+} & \textbf{RefCOCOg} & \textbf{MeViS / Ref-DAVIS17} \\
    \midrule
    CRIS~\cite{wang2022cris}                 & (C) end-to-end, CLIP       & 70.5 & 62.3 & 59.9 & -- \\
    LISA-7B~\cite{lai2024lisa}               & (E) mask-decoder MLLM      & 74.9 & 65.1 & 67.9 & -- \\
    GLaMM-7B~\cite{rasheed2024glamm}         & (E) grounded conversation  & 79.5 & 72.6 & 74.2 & -- \\
    VISA-13B~\cite{yan2024visa}              & (E) video reasoning MLLM   & --   & --   & --   & 44.5 / 70.4 \\
    Sa2VA-8B~\cite{yuan2025sa2va}            & (E) MLLM + SAM~2           & 81.9 & 76.5 & 78.9 & 46.9 / 75.2 \\
    \bottomrule
    \end{tabular}
\end{table}

\subsubsection{Challenges}

\textbf{Ambiguity and interaction.} ``The guy on the right'' depends on the viewpoint, and ``the red cup'' presumes one salient red cup. People resolve this through shared context or a follow-up question; grounding models make an arbitrary choice or fall back on a dataset bias such as the largest instance. Recognizing that an expression matches several entities, abstaining under uncertainty, and asking for clarification are outside almost every current benchmark, which is one-shot (one query, one image), while real use is a dialogue with memory of what was grounded before.

\textbf{Reasoning and commonsense.} ``The tool needed to open the box'' cannot be resolved from the image alone, and adversarial evaluation shows that much of current REC accuracy rests on shortcut matching rather than reasoning~\cite{dong2026ref}. LLM-based models carry the knowledge but may answer ``a key'' when none is visible, and vision-only models lack it entirely; reasoning-driven segmentation~\cite{lai2024lisa} and agentic active perception~\cite{zheng2025deepeyes} are the two bridges. Video adds temporal and causal structure (``the man who picks up the phone and then leaves'', ``the car that caused the accident'') that current video grounding, which attends mostly to appearance and simple motion, does not capture.

\textbf{Open world and 3D.} Open-vocabulary detectors and CLIP-based models have blind spots on rare objects, fine-grained distinctions (viola versus violin), proper nouns, and figurative references (``the blue crown'' for a cap), and per-category analyses of large multimodal models show a strong bias toward common training categories~\cite{chen2025revisiting}. Description-conditioned pre-training~\cite{li2023desco} and generative pre-training~\cite{bagchi2025refereverything,gabeur2026image} are the routes to long-tail coverage that do not require labeling it. Functional references (``the driver's seat'') also imply spatial context that only 3D understanding resolves, and the 3D-language work surveyed above~\cite{hong20233d,zhen20243d} will grow with augmented reality and robotics.

\textbf{Efficiency and evaluation.} The strongest systems are multi-model tool chains or LLMs with vision adapters and do not fit AR glasses or robots; compact models~\cite{xiao2024florence,reid2026moondream} and non-autoregressive selection~\cite{cao2025refer} show that accuracy does not require scale, but the accuracy-efficiency frontier is unmapped. Evaluation has the saturation problem described under Datasets and Metrics above, and the new benchmarks, each with its own protocol, risk a cycle of leaderboard chasing; a single protocol spanning image, video, 3D, and dialogue, scoring rejection and clarification as well as localization, does not yet exist.



\section{Summary and Outlook}
\label{sec:p3_s6_6}

This chapter followed spatial perception from the elemental to the structured: recognizing the entities in a scene, fixing them in metric geometry, binding them into relations, layouts, and maps, carrying that structure through time, and supplementing vision with sound and language. Three observations recur across these tasks.

\textbf{Priors fill gaps that observation leaves open.} Monocular depth, surface normals, and camera motion are underdetermined by a single image: scale, reflectance, and the geometry behind an occluder are not fixed by the pixels (§\ref{sec:depth_estimation}, §\ref{sec:visual_odometry}). Learned models increasingly resolve this ambiguity with strong priors, which makes them confident on familiar scenes and fragile on transparent or reflective materials, moving objects, adverse conditions, and domains their training data under-represent. Distinguishing what a percept observes from what a prior supplies is a requirement the later chapters inherit, and Chapter~\ref{chapter:p5_discussions} returns to it as a question of observability.

\textbf{Foundation models move the bottleneck rather than remove it.} Open-vocabulary detectors and segmenters, zero-shot scene classifiers, and 3D geometry foundation models have widened what perception covers without task-specific training, yet the hardest cases persist at the long tail: small and rare objects, unseen instances in clutter, and metric judgments of distance, size, and angle, on which multimodal models remain weak even when their semantic predictions are strong (§\ref{sec:spatial_relation_extraction}). The accuracy--efficiency tension noted for detection recurs in every task of the chapter.

\textbf{Evaluation rewards the easier part of each problem.} Many benchmarks are synthetic, saturated, or confined to a few curated scenes; tracking metrics do not measure robustness to occlusion length (§\ref{sec:object_tracking}), temporal-localization metrics treat fuzzy boundaries as exact, and grounding benchmarks pose a single query against a single image rather than an interaction (§\ref{sec:vl_grounding}).

What the chapter delivers is structured spatial evidence: entities with positions, relations, trajectories, and maps. Chapter~\ref{sec:reasoning} takes that evidence as input and asks what can be inferred and planned from it, and Chapter~\ref{sec:generation} runs the mapping in the opposite direction.